\chapter{Hecke operators, newforms and twists}\label{ch:hecke}

The paper works throughout with the Hecke algebra $\T$ of $S_2(\Gamma_0(N))$, with its eigencharacters, with newforms and their oldforms, with the operators $U_p$
at primes dividing the level, with depleted forms and with quadratic twists. This chapter supplies all of that. It defines the Hecke operators by double cosets and
computes them on $q$-expansions; proves commutativity, multiplicativity and the recursion; proves that $\T$ is free over $\Z$ of rank $\dim S_2(\Gamma_0(N))$ and that
the pairing $(g,t)\mapsto a_1(tg)$ between integral forms and $\T$ is perfect; computes Petersson adjoints; states the theory of Atkin, Lehner and Li (newforms,
multiplicity one, strong multiplicity one); works out $U_p$ on oldforms, the stabilised and depleted forms and the level $M^\Sigma$ of the paper's Section 4.1; proves
the decomposition $\T\otimes\Q=\Q\times\T'$ used in Lemma lem:crit; and treats twisting, including the theorem of Atkin and Li and the transport of congruences along
the twist map $\iota$ of Theorems thm:twist and thm:quant. It ends with the $L$-function of a newform and the Rankin-Selberg residue. It uses \Cref{ch:rings} (Gauss
sums, Dirichlet characters), \Cref{ch:modules} (lattices, saturation, duality of free modules, idempotents of finite algebras, generalised eigenspaces, congruence
ideals), \Cref{ch:ant} (algebraic integers, orders) and, above all, \Cref{ch:modforms} (the slash operator, $q$-expansions, the Petersson product, the Atkin-Lehner
involutions, the integral structure $S_2(\Gamma_0(N);\Z)$ and Sturm's bound). Galois representations of newforms (\Cref{ch:modgalrep}) appear only in sketches and
remarks.

\section*{Core results at a glance}

Throughout, $N\ge1$ is the level, $k\ge2$ an even weight (the paper needs $k=2$), $\Gamma=\Gamma_0(N)$, $S_k(N)=S_k(\Gamma_0(N))$, $M_k(N)=M_k(\Gamma_0(N))$,
$(f|_k\gamma)(z)=(\det\gamma)^{k/2}(cz+d)^{-k}f(\gamma z)$, $V_d\colon h(z)\mapsto h(dz)$, $\langle\ ,\ \rangle$ the unnormalised Petersson product, $t^\dagger$ the adjoint of an
operator, $\T=\T_N$ the full Hecke algebra (the $U_p$ included), $\T^{(N)}$ the anemic one, $\chi_{q^*}$ the quadratic character of conductor $q$ and $f^\Sigma$ the
$\Sigma$-depleted form.

\begin{enumerate}[label=\textbf{\thechapter.\arabic*}]
\item \textbf{Conventions.} The standing notation; the double coset operator, normalised as
$[\Gamma\alpha\Gamma]_kf=(\det\alpha)^{k/2-1}\sum_if|_k\alpha_i$, equal to Diamond-Shurman's; $T_p$ at $p\mid N$ is written $U_p$, and the level is written as a superscript when two levels occur, $T_p^{(M)}h=U_p^{(Mp)}h+p^{k-1}h(pz)$; naive twists $f\otimes\chi$ and depletion
$f^\Sigma$. \Cref{conv:hecke:standing}, \Cref{conv:hecke:double-coset}, \Cref{conv:hecke:level}, \Cref{conv:hecke:twists}.
\item \textbf{Double coset operators} (definition). For $\alpha\in\GL_2^+(\Q)$, $\Gamma\alpha\Gamma=\bigsqcup_{i=1}^r\Gamma\alpha_i$ with $r=[\Gamma:\Gamma\cap\alpha^{-1}\Gamma\alpha]$ finite, and
$[\Gamma\alpha\Gamma]_k$ is well defined and preserves $M_k(N)$ and $S_k(N)$. \emph{(proved here)} \Cref{def:hecke:double-coset}, \Cref{prop:hecke:double-coset}.
\item \textbf{Coset decompositions.} For $p\nmid N$, $\Gamma\sm100p\Gamma=\bigsqcup_{j=0}^{p-1}\Gamma\sm1j0p\sqcup\Gamma\sm p001=\Gamma\sm p001\Gamma$; for $p\mid N$,
$\Gamma\sm100p\Gamma=\bigsqcup_{j=0}^{p-1}\Gamma\sm1j0p$. $T_p=[\Gamma\sm100p\Gamma]_k$, written $U_p$ when $p\mid N$. \emph{(proved here)} \Cref{prop:hecke:coset-decomposition}.
\item \textbf{$q$-expansions.} $a_n(T_pf)=a_{pn}(f)+p^{k-1}a_{n/p}(f)$ for $p\nmid N$ and $a_n(U_pf)=a_{pn}(f)$ for $p\mid N$, for all $n\ge0$. \emph{(proved here)}
\Cref{prop:hecke:Tp-q-expansion}.
\item \textbf{The operators $\mathcal U_p$, $V_p$.} On $q$-series $\mathcal U_pV_p=1$, $V_p\mathcal U_p$ is the projection onto the multiples of $p$, and the operators at different primes (and $V_d$ with
$(d,p)=1$) commute. \emph{(proved here)} \Cref{lem:hecke:UV}.
\item \textbf{The operators $T_n$.} With $T_{p^{r+1}}=T_pT_{p^r}-p^{k-1}T_{p^{r-1}}$ ($p\nmid N$), $T_{p^r}=U_p^r$ ($p\mid N$) and $T_n=\prod_{p^r\|n}T_{p^r}$: all $T_n$ commute;
$T_mT_n=\sum_{d\mid(m,n),(d,N)=1}d^{k-1}T_{mn/d^2}$; $a_m(T_nf)=\sum_{d\mid(m,n),(d,N)=1}d^{k-1}a_{mn/d^2}(f)$, so $a_1(T_nf)=a_n(f)$; $\sum_rT_{p^r}X^r=(1-T_pX+p^{k-1}X^2)^{-1}$
or $(1-U_pX)^{-1}$. \emph{(proved here)} \Cref{thm:hecke:Tn}.
\item \textbf{The Hecke algebra} (definition and proposition). $\T=\Z[T_n:n\ge1]$ is commutative, equal to the $\Z$-span of the $T_n$, preserves $S_2(N;\Z)$, and is a free
$\Z$-module of finite rank; $\T^{(N)}$ is the $\Z$-span of the $T_n$ with $(n,N)=1$. \emph{(proved here, given the integral basis theorem of \Cref{ch:modforms})}
\Cref{def:hecke:hecke-algebra}, \Cref{prop:hecke:hecke-algebra}.
\item \textbf{Eigenforms} (definition and proposition). A simultaneous eigenvector of $\T$ has $a_1\ne0$; normalised ($a_1=1$) it satisfies $T_nf=a_n(f)f$, $a_{mn}=a_ma_n$ for
$(m,n)=1$, $a_{p^{r+1}}=a_pa_{p^r}-p^{k-1}a_{p^{r-1}}$ ($p\nmid N$), $a_{p^r}=a_p^r$ ($p\mid N$); $\phi_f\colon\T\to\C$, $T_n\mapsto a_n(f)$, is a ring homomorphism; distinct normalised
eigenforms are linearly independent. \emph{(proved here)} \Cref{def:hecke:eigenform}, \Cref{prop:hecke:eigenforms}.
\item \textbf{Eigenform criterion.} A cusp form $g=\sum b_nq^n$ of level $L$ with $b_1=1$ is a normalised eigenform of $\T_L$ iff $(b_n)$ is multiplicative and satisfies the
prime-power relations of level $L$. \emph{(proved here)} \Cref{lem:hecke:eigen-criterion}.
\item \textbf{Coefficient fields.} For a normalised eigenform $f$ in $S_2(N)$ the $a_n(f)$ are algebraic integers, $K_f=\Q(a_n(f))$ is a number field, $\phi_f(\T)$ is an order of
$\mathcal O_{K_f}$, and $f^\sigma=\sum\sigma(a_n)q^n$ is again a normalised eigenform for $\sigma\in\Aut(\C)$. \emph{(proved here, given \Cref{ch:ant})}
\Cref{prop:hecke:coefficient-field}.
\item \textbf{The perfect pairing.} $S=S_2(N;\Z)$, $d=\dim S_2(N)$: $\langle g,t\rangle=a_1(tg)$ is a perfect pairing $S\times\T\to\Z$ with $\langle ug,t\rangle=\langle g,ut\rangle$; $\T$ is free of rank
$d$, $S\simeq\Hom(\T,\Z)$, $\T\simeq\Hom(S,\Z)$, and the pairing stays perfect after $\otimes A$ for every ring $A$. $\T$ is spanned by $T_1,\dots,T_B$, $B$ Sturm's bound,
but in general not by $T_1,\dots,T_d$. \emph{(proved here, given the integral basis theorem; the last part given Sturm's bound)} \Cref{thm:hecke:perfect-pairing}.
\item \textbf{Eigenforms and homomorphisms.} For every commutative ring $A$, $S_2(N;A)\simeq\Hom_\Z(\T,A)$, and normalised eigenforms with coefficients in $A$ correspond to ring
homomorphisms $\phi\colon\T\to A$, $\phi\mapsto\sum\phi(T_n)q^n$; maximal ideals of residue characteristic $\ell$ correspond to mod $\ell$ eigenforms up to conjugacy.
\emph{(proved here, given \Cref{thm:hecke:perfect-pairing})} \Cref{thm:hecke:eigen-dictionary}.
\item \textbf{Consequences of the pairing.} A form is determined by its pairing with $\T$; $S[\mf a]\simeq\Hom(\T/\mf a,\Z)$ is saturated; $S[\ker\phi]=\Z f_\phi$ for $\phi\colon\T\to\Z$; if $I\cap J=0$, $\T/I$ and $\T/J$ are torsion-free and $\T/(I+J)$ is finite, then $S/(S[I]\oplus S[J])$ has the order of $\T/(I+J)$. \emph{(proved here)} \Cref{cor:hecke:pairing-consequences}.
\item \textbf{Congruence ideal and congruence number} (definition). $\eta(\phi)=\phi(\Ann_{\T_{\mathcal O}}\ker\phi)$; $r_f$ the largest $r$ with $f\equiv g\pmod r$ for an integral
$g\perp f$; $r_f$ is finite. \emph{(proved here)} \Cref{def:hecke:congruence-ideal}, \Cref{prop:hecke:rf-finite}.
\item \textbf{Adjoints of double cosets.} $\langle[\Gamma\alpha\Gamma]_kf,g\rangle=\langle f,[\Gamma\alpha'\Gamma]_kg\rangle$ with $\alpha'=(\det\alpha)\alpha^{-1}$. \emph{(proved here)}
\Cref{prop:hecke:double-coset-adjoint}.
\item \textbf{Self-adjointness.} $T_n$ is self-adjoint on $S_k(N)$ for $(n,N)=1$; $S_k(N)$ is the orthogonal sum of the eigenspaces of $\T^{(N)}$, whose eigenvalues are real.
\emph{(proved here)} \Cref{thm:hecke:self-adjoint}.
\item \textbf{The adjoint of $U_p$.} $U_p^\dagger=w_NU_pw_N$ for $p\mid N$; $U_p$ is in general neither self-adjoint nor normal; for a newform $h$ of level $M$ and $p\nmid M$,
at level $Mp$, $\langle h(pz),h\rangle=\frac{p^{1-k}a_p(h)}{p+1}\langle h,h\rangle$ and $\langle h(pz),h(pz)\rangle=p^{-k}\langle h,h\rangle$. \emph{(proved here)} \Cref{lem:hecke:trace},
\Cref{prop:hecke:Up-adjoint}.
\item \textbf{Stable complements.} If $V$ is stable under $\T$ and $w_N$ then so is $V^\perp$; in particular $f^\perp$ is $\T$-stable for an eigenform $f$ of $\T$ that is an
eigenvector of $w_N$. \emph{(proved here)} \Cref{cor:hecke:perp-stable}.
\item \textbf{Oldforms and newforms} (definition and proposition). $S_k^{\mathrm{old}}(N)=\sum_{p\mid N}(S_k(N/p)+V_pS_k(N/p))$, $S_k^{\mathrm{new}}(N)$ its orthogonal complement;
both are stable under $\T$ and $w_N$, and $w_N$ maps $V_dS_k(M)$ onto $V_{N/(Md)}S_k(M)$. \emph{(proved here)} \Cref{def:hecke:old-new}, \Cref{prop:hecke:old-new}.
\item \textbf{Atkin-Lehner theory.} $S_k^{\mathrm{new}}(N)$ has an orthogonal basis of newforms; multiplicity one: the $\T^{(L)}$-eigenspace in $S_k(L)$ of the system of a newform
$f$ of level $M\mid L$ is $\operatorname{span}\{f(dz):d\mid L/M\}$; $S_k(N)=\bigoplus_{M\mid N}\bigoplus_{d\mid N/M}V_dS_k^{\mathrm{new}}(M)$ with basis the $h(dz)$; $f^\perp=\bigoplus_{h\ne f}V_h$.
\emph{(stated; reference, with a sketch)} \Cref{thm:hecke:atkin-lehner}.
\item \textbf{Strong multiplicity one.} Two newforms with $a_p(f)=a_p(g)$ for almost all $p$ (in weight $2$: for a set of primes of density one) are equal; a generalised
eigenform of $\T^{(L)}$ whose system differs from that of a newform $f$ at infinitely many $p$ is orthogonal to every $f(dz)$. \emph{(stated; reference, with a sketch;
the consequence proved here)} \Cref{thm:hecke:strong-multiplicity-one}.
\item \textbf{Local behaviour of newforms.} A newform $f$ of level $N$ satisfies $f|_kW_N=\varepsilon_N(f)f$, $\varepsilon_N(f)=\pm1$ \emph{(proved here, given
\Cref{thm:hecke:atkin-lehner})}; $a_p(f)$ is real for $p\mid N$ \emph{(proved here)}; $a_p(f)=0$ if $p^2\mid N$ and $a_p(f)=-p^{k/2-1}\varepsilon_p(f)$ if $p\,\|\,N$ \emph{(stated;
reference, with a sketch)}. \Cref{prop:hecke:newform-local}.
\item \textbf{Galois conjugates.} The old and new subspaces are defined over $\Q$; conjugates of newforms are newforms; $K_f$ is totally real for a newform $f$. \emph{(proved
here, given \Cref{thm:hecke:atkin-lehner})} \Cref{cor:hecke:galois-newforms}.
\item \textbf{Dimensions.} With $\beta=\mu*\mu$:
\[
\dim S_k(N)=\sum_{M\mid N}\sigma_0(N/M)\dim S_k^{\mathrm{new}}(M),\qquad\dim S_k^{\mathrm{new}}(N)=\sum_{M\mid N}\beta(N/M)\dim S_k(M).
\] \emph{(proved
here, given \Cref{thm:hecke:atkin-lehner})} \Cref{cor:hecke:dimensions}.
\item \textbf{$U_p$ on oldforms.} For a newform $h$ of level $M$, on $V_{h,p,e}=\operatorname{span}\{h(p^jz):0\le j\le e\}$: $U_ph(p^jz)=h(p^{j-1}z)$, $U_ph=a_ph-p^{k-1}h(pz)$ ($p\nmid M$) or
$a_ph$ ($p\mid M$); $h(p^ez)$ is cyclic; characteristic polynomial $X^{e-1}(X^2-a_pX+p^{k-1})$ or $X^e(X-a_p)$; eigenvectors: the stabilisations $h-\beta_ph(pz)$, $h-\alpha_ph(pz)$,
$h$, and the depleted forms; semisimple iff the characteristic polynomial is squarefree. \emph{(proved here, given \Cref{thm:hecke:atkin-lehner})} \Cref{prop:hecke:Up-oldspan}.
\item \textbf{Eigenvalues of $U_p$.} In weight $2$ every eigenvalue of $U_p$ on $S_2(N)$ is $0$, $\pm1$ or a root of $X^2-a_p(h)X+p$; the nonzero ones are units at every
prime not above $p$. \emph{(proved here, given \Cref{thm:hecke:atkin-lehner} and \Cref{prop:hecke:newform-local})} \Cref{cor:hecke:Up-eigenvalues}.
\item \textbf{Depletion} (definition and proposition). $\delta_p=2,1,0$ for $p\nmid M$, $p\,\|\,M$, $p^2\mid M$; $f^\Sigma=\prod_{p\in\Sigma}D_pf$; $f^\Sigma$ has least level
$M^\Sigma=M\prod_{p\in\Sigma}p^{\delta_p}$, is a normalised eigenform of $\T_{M^\Sigma}$ with $U_pf^\Sigma=0$ ($p\in\Sigma$), and its generalised eigenspace in $S_k(M^\Sigma)$ is $\C f^\Sigma$.
\emph{(proved here, given \Cref{thm:hecke:atkin-lehner})} \Cref{def:hecke:depletion}, \Cref{prop:hecke:depletion}.
\item \textbf{The structure of $\T$.} $S_k(N)=\bigoplus_hV_h$, $V_h\simeq\bigotimes_{p\mid N}V_{h,p,e_p}$ as a module over the $U_p$; $\T\otimes\Qbar\simeq\prod_h\bigotimes_p\Qbar[X]/(m_{h,p})$; $S_k(N)$ is
free of rank one over $\T\otimes\C$; $\T\otimes\Q$ is reduced iff every $m_{h,p}$ is squarefree. \emph{(proved here, given \Cref{thm:hecke:atkin-lehner})} \Cref{thm:hecke:structure}.
\item \textbf{The newform splitting.} For a newform $f\in S_2(N)$ with rational coefficients, $f^\perp$ is $\T$-stable and $\T\otimes\Q=\Q\times\T'$, the first factor acting on $\Q f$
through $\phi_f$ and $\T'$ acting faithfully on $f^\perp$; the idempotent of the first factor lies in $\T\otimes\Q$ and kills $f^\perp$. \emph{(proved here, given
\Cref{thm:hecke:atkin-lehner})} \Cref{thm:hecke:newform-splitting}.
\item \textbf{$\T$-stable versus Petersson complements.} For an eigencharacter $\pi$ whose generalised eigenspace is a line $\C g$: $S_k(N)=\C g\oplus g^c$ with $g^c$ the $\T$-stable
complement, $\T\otimes\Q=\Q\times\T''$ when $\pi$ is rational, $g^c\cap S_2(N;\Z)$ saturated of rank $d-1$; $g^c=g^\perp$ for newforms, not in general (level $44$).
\emph{(proved here)} \Cref{prop:hecke:complements}.
\item \textbf{Local components.} $\T_\ell=\prod_{\mf m}\T_{\mf m}$; the pairing is perfect on each $e_{\mf m}S_2(N;\Zl)\times\T_{\mf m}$; a homomorphism $\phi\colon\T\to\mathcal O_\lambda$ factors
through one $\T_{\mf m}$, the other factors lie in $\ker\phi$, and $\eta(\phi)$ is computed in $\T_{\mf m}$; $e_{\mf m}$ preserves integrality, $\T$-stable subspaces and
congruences. \emph{(proved here)} \Cref{prop:hecke:local-components}.
\item \textbf{Congruences.} For normalised eigenforms of $\T_N$: congruence at all $n$ $\iff$ at all primes $\iff$ $\phi_f\equiv\phi_g$ on $\T$; this implies, and when $\T/\mf m_{f,\lambda}=\F_\ell$ is
equivalent to, $\mf m_{f,\lambda}=\mf m_{g,\lambda}$; for newforms, congruence at the primes outside $\Sigma$ gives $f^\Sigma\equiv g^\Sigma$ at all $n$. \emph{(proved here)} \Cref{prop:hecke:congruences}.
\item \textbf{The twist operator.} For primitive $\chi$ of conductor $r$: $f\otimes\chi=\mathfrak g(\bar\chi)^{-1}\sum_u\bar\chi(u)f|_k\sm1{u/r}01$ lies in $S_k(\Gamma_0(\lcm(M,r^2)),\chi^2)$; for
quadratic $\chi$ it preserves integrality, intertwines $T_p$ with $\chi(p)T_p$ ($p\nmid r$), kills $U_p$ for $p\mid r$, and $(f\otimes\chi)\otimes\chi=f^{\{p\mid r\}}$; the
paper's $\iota=\iota_S$. \emph{(proved here)} \Cref{prop:hecke:twist-operator}.
\item \textbf{Atkin-Li.} For a newform $g$ of weight $2$ on $\Gamma_0(M)$ and $\chi=\chi_{q^*}$ with $q^2\nmid M$, $g\otimes\chi$ is a newform on $\Gamma_0(\lcm(M,q^2))$; for $q^2\mid M$ it
is the $q$-depletion of a newform of level dividing $M$. \emph{(stated; reference, with a sketch; the case $q^2\mid M$ proved here)} \Cref{thm:hecke:atkin-li}.
\item \textbf{Transport along $\iota$.} $\iota$ maps $g^\perp$ into the forms orthogonal to every $f(dz)$, $f$ the newform with system $\chi_S(p)a_p(g)$; congruences
$g\equiv h$ become $f\equiv\iota(h)$ when $f=\iota(g)$; conversely $\iota(f)=g^S$ and $\iota(f^\perp)$ lies in the $\T$-stable complement of $g^S$. \emph{(proved here, given
\Cref{thm:hecke:atkin-lehner} and \Cref{thm:hecke:strong-multiplicity-one})} \Cref{prop:hecke:twist-transport}.
\item \textbf{Euler products.} $L(f,s)=\prod_{p\nmid N}(1-a_pp^{-s}+p^{k-1-2s})^{-1}\prod_{p\mid N}(1-a_pp^{-s})^{-1}$ for a normalised eigenform of $\T_N$; $L(f^\Sigma,s)=L(f,s)\prod_{p\in\Sigma}L_p(f,s)^{-1}$.
\emph{(proved here)} \Cref{prop:hecke:euler-product}.
\item \textbf{Functional equation.} If $f|_kW_N=\varepsilon_Nf$ then $\Lambda(f,s)=N^{s/2}(2\pi)^{-s}\Gamma(s)L(f,s)$ is entire and $\Lambda(f,s)=i^k\varepsilon_N\Lambda(f,k-s)$; sign $-\varepsilon_N$ in
weight $2$. \emph{(proof sketched here)} \Cref{thm:hecke:functional-equation}.
\item \textbf{Rankin-Selberg.} $\operatorname{Res}_{s=k}\sum|a_n(f)|^2n^{-s}=\frac{3(4\pi)^k}{\pi\Gamma(k)\mu(N)}\langle f,f\rangle$; for a rational newform of weight $2$ and squarefree level
$\langle f,f\rangle=N\,L(\Sym f,2)/(8\pi^3)$; the Rankin series is not the symmetric square $L$-function. \emph{(proof sketched here; the local identities proved)}
\Cref{thm:hecke:rankin-selberg}.
\end{enumerate}

\section*{Where the paper uses this chapter}

\begin{itemize}
\item \textbf{Section 3, preamble} ($\T$ is generated by all $T_n$, ``including the operators $U_p$ for $p\mid N$''): \Cref{prop:hecke:coset-decomposition},
\Cref{thm:hecke:Tn}, \Cref{prop:hecke:hecke-algebra}.
\item \textbf{Lemma lem:crit} (congruence criterion): $\phi_f\colon\T\to\Z$, \Cref{prop:hecke:eigenforms}; the perfect pairing and its version over $\Zl$,
\Cref{thm:hecke:perfect-pairing}; the forms killed by $I$ and by $J$ and the index of $\Z f\oplus L$, \Cref{cor:hecke:pairing-consequences}; $\T\otimes\Q=\Q\times\T'$ and the
$\T$-stability of $f^\perp$, \Cref{thm:hecke:newform-splitting}, \Cref{cor:hecke:perp-stable}; congruences on all of $\T$ and the maximal ideal of $(g,\lambda)$,
\Cref{prop:hecke:congruences}, \Cref{thm:hecke:eigen-dictionary}. The lemma itself is proved in \Cref{ch:modgalrep}.
\item \textbf{Lemma lem:old} (oldforms): the $\T$-stable span of the $h(dz)$ orthogonal to $f$, \Cref{thm:hecke:atkin-lehner}; every formula for $U_p$, the stabilisations, the
depleted forms and the non-semisimplicity, \Cref{prop:hecke:Up-oldspan}; $a_p=0$ for $p^2\mid M$, \Cref{prop:hecke:newform-local}; the tensor structure,
\Cref{thm:hecke:structure}; the passage from primes to all $n$, \Cref{prop:hecke:congruences}. The lemma itself is proved in \Cref{ch:modgalrep}.
\item \textbf{Lemma lem:twist (iii)} (Atkin-Li): \Cref{thm:hecke:atkin-li}.
\item \textbf{Theorem thm:qual (b), (c)}: twists of newforms, $h\otimes\chi$ as a depletion, strong multiplicity one, \Cref{prop:hecke:twist-operator},
\Cref{thm:hecke:atkin-li}, \Cref{thm:hecke:strong-multiplicity-one}.
\item \textbf{Section 4.1}: $f^\Sigma$, $M^\Sigma$, $\delta_p$, \Cref{prop:hecke:depletion}; the Hecke action and $w_{M^\Sigma}T_nw_{M^\Sigma}^{-1}$ as adjoint, whose Petersson
analogue is \Cref{prop:hecke:Up-adjoint}.
\item \textbf{Theorem thm:twist}: the twist map $\iota$, generalised eigenspaces and the transport $g^\perp\to f^\perp$, \Cref{prop:hecke:twist-operator},
\Cref{prop:hecke:twist-transport}, \Cref{ex:hecke:twist-pairs}. The theorem is proved in \Cref{ch:congruence}.
\item \textbf{Theorem thm:quant, Step 1}: strong multiplicity one, the eigenvalue $0$ of $U_q$ with multiplicity one, units among the $U_q$-eigenvalues, the $\T$-stable
complement of $g^S$, idempotents of $\T_\ell$, \Cref{prop:hecke:depletion}, \Cref{cor:hecke:Up-eigenvalues}, \Cref{prop:hecke:complements},
\Cref{prop:hecke:local-components}, \Cref{prop:hecke:twist-transport}. The theorem is proved in \Cref{ch:congruence}.
\item \textbf{Lemma lem:hecke and the paragraph after it} (the congruence ideal computed in $\T_{\mf m}$): \Cref{prop:hecke:local-components}.
\item \textbf{Section 1 and the remark after the Example of Section 3}: the congruence number $r_E$, \Cref{def:hecke:congruence-ideal}, \Cref{ex:hecke:congruence-numbers}
(99a1, 539d1); the Petersson norm and Rankin-Selberg, \Cref{thm:hecke:rankin-selberg}.
\item \textbf{The Example of Section 3}: $f_{\text{11a}}\otimes\chi_{-11}$ is the form of 121d, \Cref{thm:hecke:atkin-li}, \Cref{ex:hecke:level-121}.
\end{itemize}

\section{Introduction and conventions}\label{sec:hecke:intro}

This section fixes notation. Its purpose is to make every statement of the chapter unambiguous about three things that the paper handles silently: the
normalisation of the Hecke operators, the level at which an operator acts, and the difference between a twisted $q$-series and a newform. The physics reader
meets three familiar objects. The Hecke recursion $T_{p^{r+1}}=T_pT_{p^r}-pT_{p^{r-1}}$ is the three-term recursion of the Chebyshev polynomials of the second kind,
and the product formula for $T_{p^a}T_{p^b}$ is the Clebsch-Gordan rule for characters of $\mathrm{SU}(2)$ (\Cref{rem:hecke:chebyshev}). The commuting family of
self-adjoint operators $T_n$, $(n,N)=1$, is diagonalised by an orthogonal basis, as in the spectral theorem for commuting Hermitian operators
(\Cref{thm:hecke:self-adjoint}). The functional equation of $L(f,s)$ comes from the Mellin transform of $f(iy)$ (\Cref{thm:hecke:functional-equation}).

\begin{convention}[standing notation]\label{conv:hecke:standing}
(a) $N\ge1$ is the level, $k\ge2$ an even integer, the weight; everything the paper needs is $k=2$, and the general-$k$ form is given where it costs nothing.
$S_k(N)=S_k(\Gamma_0(N))$ and $M_k(N)=M_k(\Gamma_0(N))$, with trivial character. For a commutative ring $A$, $S_2(N;A)=S_2(N;\Z)\otimes_\Z A$, where $S_2(N;\Z)$ is the group of
cusp forms with integer $q$-expansion at $\infty$ (\Cref{ch:modforms}). The letter $q$ is the variable $e^{2\pi iz}$ of $q$-expansions and also an odd prime (in
Sections~\ref{sec:hecke:twists} and~\ref{sec:hecke:examples}); the two never meet in one display.
(b) $\T=\T_N=\Z[T_n:n\ge1]\subset\End_\C S_k(N)$ is the \emph{full} Hecke algebra of level $N$ (the paper's $\T$), $\T^{(N)}=\Z[T_n:(n,N)=1]$ the \emph{anemic} one,
$\T_A=\T\otimes_\Z A$ and $\T_\ell=\T\otimes\Zl$. Unless said otherwise $\T$ acts on cusp forms.
(c) An \emph{eigenform} is a nonzero simultaneous eigenvector of all of $\T$ (the $U_p$ included); an \emph{anemic eigenform} is one for $\T^{(N)}$ only; \emph{normalised}
means $a_1=1$; a \emph{newform} is a normalised eigenform in the new subspace (\Cref{def:hecke:old-new}). A \emph{system of eigenvalues} is a ring homomorphism
$\phi\colon\T\to\Qbar$ or $\T^{(N)}\to\Qbar$; it is written with an upright $\phi$, the curly $\varphi$ being reserved for maps of curves. Its \emph{generalised eigenspace}
$V(\phi)$ is the set of $v$ killed by a power of $t-\phi(t)$ for every $t$ (\Cref{ch:modules}); the \emph{$\T$-stable complement} of $V(\phi)$ is the sum of the other
generalised eigenspaces.
(d) $\langle f,g\rangle=\int_{\Gamma_0(N)\backslash\HH}f\bar g\,y^k\,dx\,dy/y^2$ is the Petersson product of \Cref{ch:modforms}, unnormalised and linear in the first variable; PARI's
\texttt{mfpetersson} returns $\langle f,g\rangle/\mu(N)$ with $\mu(N)=[\SL_2(\Z):\Gamma_0(N)]=N\prod_{p\mid N}(1+p^{-1})$. Adjoints of operators are written $t^\dagger$. The pairing
$\langle g,t\rangle=a_1(tg)$ of \Cref{sec:hecke:pairing} is $\Z$-bilinear; it is a different object, written with the same brackets only inside that section and always
with one argument in $\T$.
(e) ``$f\equiv g\pmod{\mf a}$'' means $a_n(f)\equiv a_n(g)\pmod{\mf a}$ for \emph{all} $n\ge1$; a congruence only at the primes outside a set is always written out as such.
$\lambda$ is a prime of the ring of integers $\mathcal O$ of a number field, above $\ell$.
(f) $W_N=\sm0{-1}N0$ and $w_N$ is the operator $f\mapsto f|_kW_N$ (\Cref{ch:modforms}); for even $k$ it is an involution, since $W_N^2=-N\cdot1$ and scalars act
trivially in even weight. $\varepsilon_N(f)=\pm1$ denotes its eigenvalue on a newform.
\end{convention}

\begin{convention}[double coset operators]\label{conv:hecke:double-coset}
For $\alpha\in\GL_2^+(\Q)$ with $\Gamma\alpha\Gamma=\bigsqcup_i\Gamma\alpha_i$ we put
\[
[\Gamma\alpha\Gamma]_kf=(\det\alpha)^{k/2-1}\sum_if|_k\alpha_i,
\]
with the slash operator $(f|_k\gamma)(z)=(\det\gamma)^{k/2}(cz+d)^{-k}f(\gamma z)$ of \Cref{ch:modforms}. Diamond and Shurman \cite[Section 5.1]{DiamondShurman} use the slash
$(\det\gamma)^{k-1}(cz+d)^{-k}f(\gamma z)$ and no prefactor; since $(\det\alpha)^{k/2-1}(\det\alpha)^{k/2}=(\det\alpha)^{k-1}$, the two operators coincide. In weight $2$ the prefactor
is $1$.
\end{convention}

\begin{convention}[names of Hecke operators and their level]\label{conv:hecke:level}
At level $N$ the operators are $T_n$, $n\ge1$. For a prime $p\mid N$ the operator $T_p$ is written $U_p$, as in the paper, and $T_{p^r}=U_p^r$. The operator $T_p$
depends on the level: for $h\in M_k(M)$ and $p\nmid M$, the operator $T_p$ of level $M$ and the operator $U_p$ of level $Mp$ differ,
\[
T_p^{(M)}h=U_p^{(Mp)}h+p^{k-1}h(pz)
\]
(\Cref{prop:hecke:Tp-q-expansion}). Whenever two levels are in play the level is written as a superscript. On formal $q$-series we use the operators
\[
\mathcal U_p\colon\sum a_nq^n\mapsto\sum a_{pn}q^n,\qquad V_d\colon\sum a_nq^n\mapsto\sum a_nq^{dn},
\]
so that $V_dh(z)=h(dz)$; by \Cref{prop:hecke:Tp-q-expansion} the operator $U_p$ of any level divisible by $p$ is $\mathcal U_p$ on $q$-expansions, which is why the
level of $U_p$ can often be left implicit.
\end{convention}

\begin{convention}[twists and depletion]\label{conv:hecke:twists}
(a) For a Dirichlet character $\chi$ of conductor $r$, the \emph{(naive) twist} is $f\otimes\chi=\sum_n\chi(n)a_n(f)q^n$, with $\chi(n)=0$ when $(n,r)>1$. It is a $q$-series operation;
whether $f\otimes\chi$ is a newform is a theorem (\Cref{thm:hecke:atkin-li}). (b) For an odd prime $q$, $q^*=(-1)^{(q-1)/2}q$ and $\chi_{q^*}=\bigl(\frac{q^*}{\cdot}\bigr)$ is the quadratic
character of conductor $q$, the character of $\Q(\sqrt{q^*})$; the paper writes it $\chi_q$, and the book never does, $\chi_\ell$ being the cyclotomic character. For a finite
set $S$ of odd primes, $D=\prod_{q\in S}q$ and $\chi_S=\prod_{q\in S}\chi_{q^*}$, a primitive quadratic character of conductor $D$. Values used below: $\chi_{-7}(2)=1$ and
$\chi_{-11}(2)=-1$. For an elliptic curve, $E^{(q)}=E\otimes\chi_{q^*}$ is the twist by $\Q(\sqrt{q^*})$. (c) For a finite set $\Sigma$ of primes and a $q$-series $h=\sum a_nq^n$,
the \emph{$\Sigma$-depletion} is $h^\Sigma=\sum_{(n,\Sigma)=1}a_nq^n$, where $(n,\Sigma)=1$ means that no prime of $\Sigma$ divides $n$; for one prime we write $h^{\{p\}}$. A parenthesised
superscript never denotes depletion: it is reserved for twists of curves and for removed Euler factors ($\zeta^{(N)}$, $\T^{(N)}$).
\end{convention}

\section{Hecke operators: double cosets and $q$-expansions}\label{sec:hecke:operators}

This section defines the operators $T_n$ and $U_p$ and proves the formulas that the paper uses on every page of its Section 3: the action on $q$-expansions,
$U_ph(p^jz)=h(p^{j-1}z)$, commutativity, multiplicativity and the recursion. The definition by double cosets is what makes the operators preserve modularity and
what computes their Petersson adjoints (\Cref{sec:hecke:adjoints}); the $q$-expansion formulas are what one computes with.

\begin{definition}[double coset operators]\label{def:hecke:double-coset}
Let $\Gamma=\Gamma_0(N)$ and $\alpha\in\GL_2^+(\Q)$. The \emph{double coset} $\Gamma\alpha\Gamma=\{\gamma\alpha\gamma':\gamma,\gamma'\in\Gamma\}$ is a union of right cosets $\Gamma\beta$. When it is a finite
disjoint union $\Gamma\alpha\Gamma=\bigsqcup_{i=1}^r\Gamma\alpha_i$, the operator $[\Gamma\alpha\Gamma]_k$ on functions $\HH\to\C$ is $[\Gamma\alpha\Gamma]_kf=(\det\alpha)^{k/2-1}\sum_{i=1}^rf|_k\alpha_i$
(\Cref{conv:hecke:double-coset}).
\end{definition}

\begin{proposition}\label{prop:hecke:double-coset}
Let $\Gamma=\Gamma_0(N)$, $\alpha\in\GL_2^+(\Q)$ and $\Gamma_\alpha=\Gamma\cap\alpha^{-1}\Gamma\alpha$.
\begin{enumerate}[label=(\roman*)]
\item $\Gamma_\alpha$ has finite index in $\Gamma$, and if $\Gamma=\bigsqcup_{i=1}^r\Gamma_\alpha\gamma_i$ then $\Gamma\alpha\Gamma=\bigsqcup_{i=1}^r\Gamma\alpha\gamma_i$. In particular $r=[\Gamma:\Gamma_\alpha]$.
\item $[\Gamma\alpha\Gamma]_kf$ does not depend on the choice of the representatives $\alpha_i$, and $[\Gamma\alpha\Gamma]_k$ maps $M_k(N)$ to $M_k(N)$ and $S_k(N)$ to $S_k(N)$.
\end{enumerate}
\end{proposition}

\begin{proof}
(i) Multiplying $\alpha$ by a positive integer changes neither $\Gamma_\alpha$ nor the cosets, so we may assume $\alpha$ integral; let $m=\det\alpha$. For
$\gamma\in\Gamma(Nm)$ we have $\alpha\gamma\alpha^{-1}=\alpha\gamma\cdot m^{-1}\alpha'$ with $\alpha'=m\alpha^{-1}$ integral, and $\gamma\equiv1\pmod{Nm}$ gives $\alpha\gamma\alpha'\equiv\alpha\alpha'=m\pmod{Nm}$, so
$\alpha\gamma\alpha^{-1}\equiv1\pmod N$ is an integral matrix of determinant one: it lies in $\Gamma(N)\subseteq\Gamma$. Hence $\Gamma(Nm)\subseteq\Gamma_\alpha$, and $\Gamma(Nm)$ has finite
index in $\SL_2(\Z)$. For the decomposition, the map $\Gamma_\alpha\gamma\mapsto\Gamma\alpha\gamma$ from $\Gamma_\alpha\backslash\Gamma$ to the set of right cosets in $\Gamma\alpha\Gamma$ is well
defined (for $\delta\in\Gamma_\alpha$, $\alpha\delta\gamma=(\alpha\delta\alpha^{-1})\alpha\gamma$ with $\alpha\delta\alpha^{-1}\in\Gamma$), surjective (every element of $\Gamma\alpha\Gamma$ is $\gamma'\alpha\gamma$), and
injective ($\Gamma\alpha\gamma=\Gamma\alpha\gamma'$ means $\alpha\gamma'\gamma^{-1}\alpha^{-1}\in\Gamma$, that is, $\gamma'\gamma^{-1}\in\Gamma_\alpha$).

(ii) If $\Gamma\alpha_i=\Gamma\beta_i$ then $\beta_i=\delta\alpha_i$ with $\delta\in\Gamma$ and $f|_k\beta_i=f|_k\alpha_i$ for $f$ invariant under $\Gamma$; so the sum is independent of the
representatives. For $\gamma\in\Gamma$, the cosets $\Gamma\alpha_i\gamma$ are again the cosets of $\Gamma\alpha\Gamma$, permuted (right multiplication by $\gamma$ is a bijection of
$\Gamma\alpha\Gamma$ preserving the partition into right cosets); hence $([\Gamma\alpha\Gamma]_kf)|_k\gamma=(\det\alpha)^{k/2-1}\sum_if|_k\alpha_i\gamma=[\Gamma\alpha\Gamma]_kf$. Each $f|_k\alpha_i$ is
holomorphic on $\HH$. At the cusps: for $\sigma\in\SL_2(\Z)$, write $\alpha_i\sigma=\sigma'\beta$ with $\sigma'\in\SL_2(\Z)$ and $\beta$ upper triangular in $\GL_2^+(\Q)$ (choose $\sigma'$
with first column proportional to the first column of $\alpha_i\sigma$). Then $f|_k\alpha_i\sigma=(f|_k\sigma')|_k\beta$, and $f|_k\sigma'$ has an expansion
$\sum_{n\ge0}b_ne^{2\pi inz/h}$ (with $b_0=0$ if $f$ is cuspidal) by \Cref{ch:modforms}; slashing with $\beta=\sm ab0d$ replaces $z$ by $(az+b)/d$ and multiplies by a
constant, which keeps the function bounded on $\operatorname{Im}z\ge1$ (and tending to $0$ if $b_0=0$). So $[\Gamma\alpha\Gamma]_kf$ is holomorphic (cuspidal) at every cusp.
\end{proof}

The next proposition computes the double coset that defines $T_p$. We use the set
\[
\Delta_n^N=\Bigl\{\Mtwo abcd\in\Mat_2(\Z):ad-bc=n,\ N\mid c,\ \gcd(a,N)=1\Bigr\}.
\]
It is stable under left and right multiplication by $\Gamma_0(N)$: if $\gamma=\sm xyuv\in\Gamma_0(N)$ then $x$ is a unit modulo $N$ (as $xv\equiv1\pmod N$), the lower left
entries of $\gamma\delta$ and $\delta\gamma$ are $\equiv0\pmod N$, and their upper left entries are $\equiv xa$ and $\equiv ax\pmod N$.

\begin{proposition}[coset decompositions]\label{prop:hecke:coset-decomposition}
Let $p$ be a prime and $\Gamma=\Gamma_0(N)$.
\begin{enumerate}[label=(\roman*)]
\item The right cosets $\Gamma\delta$ with $\delta\in\Delta_p^N$ are the $\Gamma\sm1j0p$, $0\le j<p$, which are distinct, together with $\Gamma\sm p001$ if $p\nmid N$.
\item If $p\nmid N$ then $\Gamma\sm100p\Gamma=\Gamma\sm p001\Gamma=\Delta_p^N=\bigsqcup_{j=0}^{p-1}\Gamma\sm1j0p\sqcup\Gamma\sm p001$, a union of $p+1$ cosets.
\item If $p\mid N$ then $\Gamma\sm100p\Gamma=\Delta_p^N=\bigsqcup_{j=0}^{p-1}\Gamma\sm1j0p$, a union of $p$ cosets.
\end{enumerate}
We define $T_p=[\Gamma\sm100p\Gamma]_k$ on $M_k(N)$, and write $U_p$ for it when $p\mid N$ (\Cref{conv:hecke:level}).
\end{proposition}

\begin{proof}
(i) Let $\delta=\sm abcd\in\Delta_p^N$ and $g=\gcd(a,c)$. As $\gcd(a,N)=1$, $g$ is prime to $N$, so $N\mid c/g$; the pair $(u,v)=(-c/g,a/g)$ is coprime with $N\mid u$, so it is the
bottom row of some $\gamma\in\Gamma_0(N)$, and $\gamma\delta$ has lower left entry $ua+vc=0$. So $\Gamma\delta$ contains an upper triangular matrix $\sm{a'}{b'}0{d'}$, with $a'd'=p$; multiplying
by $-1$ we may take $a',d'>0$, and multiplying on the left by $\sm1t01$ we may take $0\le b'<d'$. This gives $\sm1j0p$ with $0\le j<p$, or $\sm p001$; the latter has upper left
entry $p$, which is prime to $N$ only if $p\nmid N$. Distinctness: $\sm1j0p\sm1{j'}0p^{-1}=\sm1{(j-j')/p}01$ is integral only if $j=j'$, and $\sm1j0p\sm p001^{-1}=\sm{1/p}j0p$ is not
integral.

(ii) $\Delta_p^N$ contains $\sm100p$ and is stable under $\Gamma$ on both sides, so it contains $\Gamma\sm100p\Gamma$, and each of the cosets of (i) lies in it. Conversely,
$\sm1j0p=\sm100p\sm1j01$ lies in $\Gamma\sm100p\Gamma$; and since $\gcd(p,N)=1$ there are integers $w,y$ with $pw-Ny=1$, so $\gamma=\sm py N w\in\Gamma_0(N)$, and
$\sm100p\gamma=\sm py{pN}{pw}=\sm1y N{pw}\sm p001$ with $\sm 1yN{pw}\in\Gamma_0(N)$ (its determinant is $pw-Ny=1$). So $\sm p001\in\Gamma\sm100p\Gamma$ and all $p+1$ cosets lie in the double
coset. Two double cosets are equal or disjoint, so $\Gamma\sm p001\Gamma=\Gamma\sm100p\Gamma$.

(iii) As in (ii), $\Gamma\sm100p\Gamma\subseteq\Delta_p^N$, which by (i) is the union of the $p$ cosets $\Gamma\sm1j0p$, each contained in $\Gamma\sm100p\Gamma$.
\end{proof}

\begin{proposition}[$q$-expansions]\label{prop:hecke:Tp-q-expansion}
Let $f=\sum_{n\ge0}a_nq^n\in M_k(N)$ and let $p$ be prime. For every $n\ge0$,
\[
a_n(T_pf)=a_{pn}(f)+p^{k-1}a_{n/p}(f)\quad(p\nmid N),\qquad a_n(U_pf)=a_{pn}(f)\quad(p\mid N),
\]
with $a_{n/p}=0$ when $p\nmid n$. In weight $2$, $a_n(T_pf)=a_{pn}+p\,a_{n/p}$. Equivalently $T_p=\mathcal U_p+p^{k-1}V_p$ ($p\nmid N$) and $U_p=\mathcal U_p$ ($p\mid N$) on $q$-expansions, and
for $h\in M_k(M)$ with $p\nmid M$, $T_p^{(M)}h=U_p^{(Mp)}h+p^{k-1}h(pz)$.
\end{proposition}

\begin{proof}
$f|_k\sm1j0p(z)=p^{k/2}p^{-k}f\bigl(\frac{z+j}p\bigr)$, so with the prefactor $p^{k/2-1}$ the coset $\sm1j0p$ contributes $p^{-1}f\bigl(\frac{z+j}p\bigr)=p^{-1}\sum_na_ne^{2\pi in(z+j)/p}$. Summing over
$j$, $\sum_{j=0}^{p-1}e^{2\pi inj/p}$ is $p$ if $p\mid n$ and $0$ otherwise, so these cosets give $\sum_{p\mid n}a_nq^{n/p}=\sum_ma_{pm}q^m$. The coset $\sm p001$ gives
$p^{k/2-1}\cdot p^{k/2}f(pz)=p^{k-1}\sum_ma_mq^{pm}$. Apply \Cref{prop:hecke:coset-decomposition}. The last statement compares the two formulas: at level $Mp$ the prime $p$
divides the level, so $U_p^{(Mp)}=\mathcal U_p$, while $T_p^{(M)}=\mathcal U_p+p^{k-1}V_p$.
\end{proof}

\begin{lemma}[the operators $\mathcal U_p$ and $V_p$]\label{lem:hecke:UV}
On formal $q$-series: $\mathcal U_pV_p=1$; $V_p\mathcal U_p$ is the projection $\sum a_nq^n\mapsto\sum_{p\mid n}a_nq^n$; and for primes $p\ne p'$ and integers $d$ prime to $p$, $\mathcal U_p$ and $V_p$ commute
with $\mathcal U_{p'}$, $V_{p'}$ and $V_d$.
\end{lemma}

\begin{proof}
Compare $n$-th coefficients: $\mathcal U_pV_p$ sends $(a_n)$ to $(a_{pn/p})=(a_n)$; $V_p\mathcal U_p$ to $(a_{p\cdot n/p})$ for $p\mid n$ and $0$ otherwise; $\mathcal U_pV_d$ and $V_d\mathcal U_p$ both have
$n$-th coefficient $a_{pn/d}$ (read as $0$ if $d\nmid n$, which for $(d,p)=1$ is the same condition as $d\nmid pn$); the other cases are similar.
\end{proof}

\begin{theorem}[the operators $T_n$]\label{thm:hecke:Tn}
Put $T_1=1$, $T_{p^{r+1}}=T_pT_{p^r}-p^{k-1}T_{p^{r-1}}$ for $p\nmid N$ and $r\ge1$, $T_{p^r}=U_p^r$ for $p\mid N$, and $T_n=\prod_{p^r\|n}T_{p^r}$. Then on $M_k(N)$:
\begin{enumerate}[label=(\roman*)]
\item for $p\nmid N$, $T_{p^r}=\sum_{i=0}^rp^{i(k-1)}V_p^i\,\mathcal U_p^{r-i}$ on $q$-expansions, and for $p\mid N$, $T_{p^r}=\mathcal U_p^r$;
\item all the $T_n$ commute;
\item for $f\in M_k(N)$ and $m\ge1$,
\[
a_m(T_nf)=\sum_{d\mid(m,n),\ (d,N)=1}d^{k-1}a_{mn/d^2}(f),\qquad\text{in particular}\qquad a_1(T_nf)=a_n(f);
\]
\item $T_mT_n=\sum_{d\mid(m,n),\ (d,N)=1}d^{k-1}T_{mn/d^2}$, in particular $T_mT_n=T_{mn}$ when $(m,n)=1$;
\item $\sum_{r\ge0}T_{p^r}X^r=(1-T_pX+p^{k-1}X^2)^{-1}$ for $p\nmid N$ and $(1-U_pX)^{-1}$ for $p\mid N$, as formal power series with coefficients in the commutative ring $\Z[T_p]$.
\end{enumerate}
\end{theorem}

\begin{proof}
(i) For $p\mid N$ this is \Cref{prop:hecke:Tp-q-expansion}. For $p\nmid N$ write $\mathcal U=\mathcal U_p$, $V=V_p$, $c=p^{k-1}$ and $R_r=\sum_{i=0}^rc^iV^i\mathcal U^{r-i}$. Then $R_0=1$,
$R_1=\mathcal U+cV=T_p$, and, using $\mathcal U V^i=V^{i-1}$ for $i\ge1$ (\Cref{lem:hecke:UV}),
\[
T_pR_r=(\mathcal U+cV)\sum_{i=0}^rc^iV^i\mathcal U^{r-i}=\mathcal U^{r+1}+\sum_{i=1}^rc^iV^{i-1}\mathcal U^{r-i}+\sum_{i=0}^rc^{i+1}V^{i+1}\mathcal U^{r-i}=R_{r+1}+cR_{r-1}.
\]
So $R_r$ satisfies the recursion defining $T_{p^r}$, with the same initial values, and $R_r=T_{p^r}$.

(ii) By (i), $T_{p^r}$ is a polynomial in $\mathcal U_p$ and $V_p$, and by \Cref{lem:hecke:UV} these commute with $\mathcal U_{p'}$, $V_{p'}$ for $p'\ne p$; so $T_{p^r}$ and $T_{p'^s}$ commute.
Powers of one prime commute because $T_{p^r}$ is a polynomial in $T_p$. Hence all $T_n$ commute (and the order in the product defining $T_n$ is irrelevant).

(iii) For $n=p^r$ and $p\nmid N$: the $m$-th coefficient of $V^i\mathcal U^{r-i}f$ is $a_{p^{r-i}m/p^i}$ if $p^i\mid m$ and $0$ otherwise; with $d=p^i$ this is $a_{mn/d^2}$ for $d\mid(m,n)$. So
(i) gives the formula. For $p\mid N$ only $d=1$ occurs and $a_m(\mathcal U_p^rf)=a_{p^rm}$. If $n=n_1n_2$ with $(n_1,n_2)=1$ and the formula holds for $n_1$ and $n_2$, then
\[
a_m(T_{n_1}T_{n_2}f)=\sum_{d_1\mid(m,n_1)}d_1^{k-1}\sum_{d_2\mid(mn_1/d_1^2,\,n_2)}d_2^{k-1}a_{mn_1n_2/(d_1d_2)^2}(f)
\]
(sums over divisors prime to $N$), and since $(n_1,n_2)=1$ the condition $d_2\mid(mn_1/d_1^2,n_2)$ is $d_2\mid(m,n_2)$; the pairs $(d_1,d_2)$ correspond to the divisors $d=d_1d_2$ of $(m,n)$
prime to $N$. Induction on the number of prime factors gives (iii).

(iv) By (ii) and the definition, both sides are multiplicative in the coprime parts of $m$ and $n$, so it suffices to treat $m=p^a$, $n=p^b$. For $p\mid N$ this is
$U_p^aU_p^b=U_p^{a+b}$. For $p\nmid N$ write $\tau_r=T_{p^r}$ (so $\tau_0=T_1=1$ and $\tau_1=T_p$), $c=p^{k-1}$, $R(a,b)=\sum_{i=0}^{\min(a,b)}c^i\tau_{a+b-2i}$, and prove $\tau_a\tau_b=R(a,b)$ for $b\ge a$ by induction on $a$; the case
$b<a$ follows by commutativity, $R$ being symmetric. The cases $a=0$ and $a=1$ are the definition. For the step let $b\ge a+1$; then $\tau_{a+1}\tau_b=\tau_1\tau_a\tau_b-c\tau_{a-1}\tau_b=\tau_1R(a,b)-cR(a-1,b)$.
In $R(a,b)$ every index $a+b-2i\ge b-a\ge1$, so $\tau_1\tau_{a+b-2i}=\tau_{a+b+1-2i}+c\tau_{a+b-1-2i}$, and
\[
\tau_1R(a,b)=\sum_{i=0}^ac^i\tau_{a+1+b-2i}+\sum_{i=1}^{a+1}c^i\tau_{a+b+1-2i},\qquad cR(a-1,b)=\sum_{i=1}^{a}c^i\tau_{a+b+1-2i}.
\]
The difference is $\sum_{i=0}^ac^i\tau_{a+1+b-2i}+c^{a+1}\tau_{b-a-1}=R(a+1,b)$.

(v) The coefficient of $X^{r+1}$ in $(\sum_rT_{p^r}X^r)(1-T_pX+p^{k-1}X^2)$ is $T_{p^{r+1}}-T_pT_{p^r}+p^{k-1}T_{p^{r-1}}=0$ for $r\ge1$, that of $X$ is $T_p-T_p=0$ and that of $X^0$ is $1$; for $p\mid N$
the series is geometric.
\end{proof}

\begin{remark}[Chebyshev polynomials and Clebsch-Gordan]\label{rem:hecke:chebyshev}
Let $U_r$ be the Chebyshev polynomial of the second kind, $U_r(\cos\theta)=\sin((r+1)\theta)/\sin\theta$, with generating function $\sum_rU_r(x)X^r=(1-2xX+X^2)^{-1}$. Comparing with
\Cref{thm:hecke:Tn}(v), $T_{p^r}=p^{r(k-1)/2}U_r\bigl(T_p/(2p^{(k-1)/2})\bigr)$ for $p\nmid N$. If $T_p$ has the eigenvalue $2p^{(k-1)/2}\cos\theta$ (the Ramanujan-Petersson bound, \Cref{ch:modforms},
says $\theta$ can be taken real), then $T_{p^r}$ has the eigenvalue $p^{r(k-1)/2}\sin((r+1)\theta)/\sin\theta$: the character of the spin-$r/2$ representation of $\mathrm{SU}(2)$ at the rotation by
$2\theta$. The product formula $T_{p^a}T_{p^b}=\sum_ip^{i(k-1)}T_{p^{a+b-2i}}$ of \Cref{thm:hecke:Tn}(iv) is then the Clebsch-Gordan rule $V_{a/2}\otimes V_{b/2}=\bigoplus_iV_{(a+b)/2-i}$ for
these characters.
\end{remark}

\begin{example}[Eisenstein series and $\Delta$]\label{ex:hecke:eisenstein-and-delta}
(i) Level one. The Eisenstein series $E_k=1+c_k\sum_{n\ge1}\sigma_{k-1}(n)q^n$ ($k\ge4$) satisfies $T_pE_k=\sigma_{k-1}(p)E_k$. Indeed the constant term of $T_pE_k$ is $1+p^{k-1}=\sigma_{k-1}(p)$,
and for $n\ge1$ one needs $\sigma_{k-1}(pn)+p^{k-1}\sigma_{k-1}(n/p)=\sigma_{k-1}(p)\sigma_{k-1}(n)$. Write $n=p^am$ with $p\nmid m$; by multiplicativity of $\sigma_{k-1}$ both sides are $\sigma_{k-1}(m)$
times $\sigma_{k-1}(p^{a+1})+p^{k-1}\sigma_{k-1}(p^{a-1})$, respectively $\sigma_{k-1}(p)\sigma_{k-1}(p^a)$, which agree (for $a=0$ read $\sigma_{k-1}(p^{-1})=0$; for $a\ge1$ expand the geometric sums).
Since $\dim S_{12}(1)=1$ (\Cref{ch:modforms}), every $T_n$ preserves the line of $\Delta=\sum\tau(n)q^n$, so $T_n\Delta=\tau(n)\Delta$ by \Cref{thm:hecke:Tn}(iii), and $\tau$ is multiplicative
and satisfies $\tau(p^{r+1})=\tau(p)\tau(p^r)-p^{11}\tau(p^{r-1})$: Mordell's theorem.

(ii) Level $11$. The weight-two Eisenstein series $E_2^{(11)}=-10-24\sum_{n\ge1}(\sigma(n)-11\sigma(n/11))q^n\in M_2(11)$ of \Cref{ch:modforms} has coefficients
$\sigma(n)-11\sigma(n/11)=\sigma(m)$ for $n=11^am$, $11\nmid m$ (as $\sigma(11^a)-11\sigma(11^{a-1})=1$). By the computation in (i) with $k=2$, $T_pE_2^{(11)}=(1+p)E_2^{(11)}$ for $p\ne11$, and
$U_{11}E_2^{(11)}=E_2^{(11)}$ because $\sigma(m)$ does not see the power of $11$ (and the constant term is fixed by $\mathcal U_{11}$). By \Cref{ch:modforms}, $f_{11}\equiv-\frac1{24}E_2^{(11)}\pmod5$
coefficientwise for $n\ge1$; so the eigenvalues of $f_{11}$ satisfy $a_p\equiv1+p\pmod5$ ($p\ne11$) and $a_{11}=1$. In terms of the algebra of the next section, $\T_{11}=\Z$ (acting on the
line $\C f_{11}$), and its maximal ideal $(5)$ contains $T_p-1-p$ ($p\ne11$) and $U_{11}-1$: it is an \emph{Eisenstein} maximal ideal (\Cref{ch:moddeg}). (Checked with PARI for
$n\le2000$.) This is the case of reducible $\bar\rho$ that the paper sets aside in Section sec:eis.
\end{example}

\begin{remark}[Hecke correspondences]\label{rem:hecke:correspondences}
In weight $2$, $f\mapsto f\,dz$ identifies $S_2(N)$ with the holomorphic differentials on $X_0(N)$ (\Cref{ch:modforms}), and $T_p$ becomes the action on differentials of the
Hecke correspondence $X_0(N)\leftarrow X_0(Np)\rightarrow X_0(N)$ given by $z\mapsto z$ and $z\mapsto pz$ (\Cref{ch:curves}). This is how the same ring $\T$ acts on $J_0(N)$ and on
$H^1(X_0(N),\Z)$ (\Cref{ch:moddeg}), and why the Hecke algebra of the paper's Section 4.1 is the one defined here. Nothing in this chapter depends on the remark
\cite[Section 6.3]{DiamondShurman}, \cite[Section 7.2]{Shimura71}.
\end{remark}

\begin{paperlink}\label{pl:hecke:lem-old-formulas}
The proof of Lemma lem:old writes ``$U_ph(p^jz)=h(p^{j-1}z)$ for $j\ge1$, and $U_ph=a_p(h)h-p\,h(pz)$ if $p\nmid M$, $U_ph=a_p(h)h$ if $p\mid M$''. Here $U_p$ is the operator of the
level $N$, $p\mid N$, which on $q$-expansions is $\mathcal U_p$ (\Cref{prop:hecke:Tp-q-expansion}). The first formula is $\mathcal U_pV_p=1$ (\Cref{lem:hecke:UV}) applied to $h(p^{j-1}z)$; the
second is $\mathcal U_ph=T_p^{(M)}h-pV_ph=a_p(h)h-p\,h(pz)$; the third is $\mathcal U_ph=U_p^{(M)}h$. The last sentence of the first paragraph of that proof, ``the congruence for all
$n$ then follows from the multiplicativity of the coefficients, from $a_{p^k}=a_p^k$ for $p\mid N$, and from the recursion'', uses \Cref{thm:hecke:Tn}(iv) in the form of
\Cref{prop:hecke:eigenforms}(ii) below.
\end{paperlink}

\section{The Hecke algebra and eigenforms}\label{sec:hecke:algebra}

This section defines the paper's Hecke algebra $\T$ (``the subring of $\operatorname{End}S_2(\Gamma_0(N))$ generated over $\Z$ by all $T_n$, $n\ge1$, including the operators $U_p$
for $p\mid N$'', Section 3), its eigenforms and the eigencharacter $\phi_f\colon\T\to\Z$ of Lemma lem:crit. The key point for the paper is that $\T$ is the $\Z$-span of the
$T_n$, so that a statement about all $a_n$ is a statement about all of $\T$.

\begin{definition}[Hecke algebras]\label{def:hecke:hecke-algebra}
$\T=\T_N$ is the subring of $\End_\C S_k(N)$ generated by the $T_n$, $n\ge1$, and $\T^{(N)}$ the subring generated by the $T_n$ with $(n,N)=1$. For a ring $A$, $\T_A=\T\otimes_\Z A$.
\end{definition}

\begin{proposition}\label{prop:hecke:hecke-algebra}
\begin{enumerate}[label=(\roman*)]
\item $\T$ is commutative.
\item $\T$ is the $\Z$-span of $\{T_n:n\ge1\}$, and $\T^{(N)}$ is the $\Z$-span of $\{T_n:(n,N)=1\}$.
\item In weight $2$, $\T$ preserves $S_2(N;\Z)$. Given the integral basis theorem ($S_2(N;\Z)$ is a lattice spanning $S_2(N)$, \Cref{ch:modforms}), restriction identifies $\T$
with a subring of $\End_\Z S_2(N;\Z)\simeq\Mat_d(\Z)$, $d=\dim S_2(N)$; so $\T$ is a free $\Z$-module of finite rank, and $\T\otimes A$ acts on $S_2(N;A)$ for every ring $A$.
\end{enumerate}
\end{proposition}

\begin{proof}
(i) is \Cref{thm:hecke:Tn}(ii). (ii) By \Cref{thm:hecke:Tn}(iv) the $\Z$-span of the $T_n$ is closed under products, and it contains $T_1=1$; so it is the ring they
generate. The same formula with $(m,n)=(m,n,N)=1$ only involves $T_{mn/d^2}$ with $(mn/d^2,N)=1$. (iii) By \Cref{thm:hecke:Tn}(iii) with $k=2$, $a_m(T_nf)$ is an integer combination of
the $a_j(f)$. Since $S_2(N;\Z)$ spans $S_2(N)$, an operator that kills $S_2(N;\Z)$ is zero; so restriction is injective, and $\T$ is a subgroup of the free group $\Mat_d(\Z)$.
\end{proof}

\begin{paperlink}\label{pl:hecke:span}
Lemma lem:crit concludes from ``$a_n(g)\equiv a_n(f)\pmod\lambda$ for every $n\ge1$'' that ``$\phi_g\equiv\phi_f\pmod\lambda$ on all of $\T$''. This is \Cref{prop:hecke:hecke-algebra}(ii): every
$t\in\T$ is $\sum c_nT_n$ with $c_n\in\Z$, so $\phi_g(t)-\phi_f(t)=\sum c_n(a_n(g)-a_n(f))\in\lambda$. The same fact makes the identity $t^t=w^*tw^*$ of \Cref{ch:curves} hold for
all $t\in\T$ once it holds for the $T_n$.
\end{paperlink}

\begin{definition}[eigenforms and eigencharacters]\label{def:hecke:eigenform}
An \emph{eigenform} (of level $N$) is a nonzero $f\in S_k(N)$ with $tf\in\C f$ for every $t\in\T$; it is \emph{normalised} if $a_1(f)=1$. Its \emph{eigencharacter} is the map
$\phi_f\colon\T\to\C$ with $tf=\phi_f(t)f$. An \emph{anemic eigenform} is a nonzero simultaneous eigenvector of $\T^{(N)}$.
\end{definition}

\begin{proposition}\label{prop:hecke:eigenforms}
Let $f\in S_k(N)$ be an eigenform with $T_nf=\lambda_nf$.
\begin{enumerate}[label=(\roman*)]
\item $a_n(f)=\lambda_na_1(f)$ for all $n$; hence $a_1(f)\ne0$, and $f/a_1(f)$ is a normalised eigenform with $T_n(f/a_1(f))=a_n(f/a_1(f))\cdot f/a_1(f)$ and $q$-expansion $\sum\lambda_nq^n$.
\item A normalised eigenform satisfies $a_{mn}=a_ma_n$ for $(m,n)=1$, $a_{p^{r+1}}=a_pa_{p^r}-p^{k-1}a_{p^{r-1}}$ for $p\nmid N$, $r\ge1$, and $a_{p^r}=a_p^r$ for $p\mid N$; more generally
$a_ma_n=\sum_{d\mid(m,n),(d,N)=1}d^{k-1}a_{mn/d^2}$.
\item $\phi_f\colon\T\to\C$ is a ring homomorphism with $\phi_f(T_n)=a_n(f)$ if $f$ is normalised. The map $f\mapsto\phi_f$ is injective on normalised eigenforms, and distinct normalised
eigenforms are linearly independent.
\end{enumerate}
\end{proposition}

\begin{proof}
(i) $a_n(f)=a_1(T_nf)=\lambda_na_1(f)$ by \Cref{thm:hecke:Tn}(iii); if $a_1(f)=0$ all coefficients vanish and $f=0$. (ii) Apply $\phi_f$ to \Cref{thm:hecke:Tn}(iv) and use
$\phi_f(T_n)=a_n$. (iii) $\phi_f(tu)f=tuf=\phi_f(u)tf=\phi_f(t)\phi_f(u)f$, and similarly for sums and $1$. A normalised eigenform is $\sum\phi_f(T_n)q^n$, so it is determined by $\phi_f$.
If $\sum_{i=1}^sc_if_i=0$ is a shortest nontrivial relation among normalised eigenforms with distinct characters, pick $t$ with $\phi_{f_1}(t)\ne\phi_{f_2}(t)$; then
$\sum_ic_i(\phi_{f_i}(t)-\phi_{f_1}(t))f_i=0$ is shorter and nontrivial, a contradiction.
\end{proof}

\begin{lemma}[eigenform criterion]\label{lem:hecke:eigen-criterion}
Let $L\ge1$ and let $g=\sum_{n\ge1}b_nq^n\in S_k(L)$ with $b_1=1$. Then $g$ is a normalised eigenform of $\T_L$ if and only if $b_{mn}=b_mb_n$ for $(m,n)=1$, $b_{p^{r+1}}=b_pb_{p^r}-p^{k-1}b_{p^{r-1}}$
for $p\nmid L$, $r\ge1$, and $b_{p^r}=b_p^r$ for $p\mid L$. Its eigenvalues are then $T_ng=b_ng$.
\end{lemma}

\begin{proof}
``Only if'' is \Cref{prop:hecke:eigenforms}(ii). Conversely, the relations say that the function $n\mapsto b_n$ satisfies the same multiplicativity and prime-power recursions as
$n\mapsto T_n$; the proof of \Cref{thm:hecke:Tn}(iv) used nothing else (it is a computation in the commutative ring generated by the $T_{p^r}$, valid for any family satisfying
these relations), so $b_mb_n=\sum_{d\mid(m,n),(d,L)=1}d^{k-1}b_{mn/d^2}$. By \Cref{thm:hecke:Tn}(iii), $a_m(T_ng)=\sum_{d}d^{k-1}b_{mn/d^2}=b_nb_m=a_m(b_ng)$ for all $m\ge1$; both are
cusp forms, so $T_ng=b_ng$.
\end{proof}

\begin{warning}[anemic eigenforms]\label{warn:hecke:anemic}
An anemic eigenform need not have $a_1\ne0$, and \Cref{prop:hecke:eigenforms}(i) fails for it: $f_{11}(2z)\in S_2(22)$ is an eigenvector of every $T_n$ with $(n,22)=1$ (they commute with
$V_2$ by \Cref{lem:hecke:UV}), and $a_1(f_{11}(2z))=0$. Only eigenforms of the full $\T$ can be normalised. This is why the paper insists on ``a normalised eigenform $g\in
S_2(\Gamma_0(N))\otimes\overline\Q$ for all of $\T$'' in Lemma lem:crit.
\end{warning}

\begin{proposition}[coefficient fields]\label{prop:hecke:coefficient-field}
Let $f$ be a normalised eigenform in $S_2(N)$.
\begin{enumerate}[label=(\roman*)]
\item Every $a_n(f)$ is an algebraic integer, $K_f=\Q(a_n(f):n\ge1)$ is a number field, and $\phi_f(\T)=\Z[a_n(f):n\ge1]$ is an order of $\mathcal O_{K_f}$.
\item For $\sigma\in\Aut(\C)$, $f^\sigma=\sum\sigma(a_n(f))q^n$ is again a normalised eigenform in $S_2(N)$, with $\phi_{f^\sigma}=\sigma\circ\phi_f$.
\end{enumerate}
\end{proposition}

\begin{proof}
(i) $\T$ is a commutative ring, finitely generated as a $\Z$-module (\Cref{prop:hecke:hecke-algebra}(iii)), and $\phi_f$ a ring homomorphism to $\C$; \Cref{ch:ant} shows that the
image of such a homomorphism is an order in the ring of integers of the number field it spans. The image is $\Z[a_n(f)]$ by \Cref{prop:hecke:hecke-algebra}(ii).
(ii) Let $g_1,\dots,g_d$ be a $\Z$-basis of $S_2(N;\Z)$; the matrices $(m^{(n)}_{ji})$ of the $T_n$ in this basis are integral. Write $f=\sum_ic_ig_i$ and put $f^\sigma:=\sum_i\sigma(c_i)g_i\in S_2(N)$;
its coefficients are $\sigma(a_n(f))$ because those of the $g_i$ are integers. From $T_nf=a_nf$ we get $\sum_im^{(n)}_{ji}c_i=a_nc_j$ for all $j$; applying $\sigma$, $T_nf^\sigma=\sigma(a_n)f^\sigma$.
\end{proof}

\begin{warning}[non-real coefficients]\label{warn:hecke:nonreal}
A normalised eigenform of the full $\T$ that is not a newform can have non-real coefficients. At level $22$ the operator $U_2$ acts on $\operatorname{span}\{f_{11}(z),f_{11}(2z)\}$ by the
matrix $\sm{-2}1{-2}0$ (columns: images of $f_{11}(z)$, $f_{11}(2z)$), with characteristic polynomial $X^2+2X+2$; its eigenvectors $f_{11}-\beta f_{11}(2z)$, $\alpha,\beta=-1\pm i$, are
normalised eigenforms of $\T_{22}$ with $a_2=\alpha\in\{-1+i,-1-i\}$ (\Cref{prop:hecke:Up-oldspan}). ``$K_f$ is totally real'' holds for newforms (\Cref{cor:hecke:galois-newforms}),
and for the eigenvalues of $\T^{(N)}$ (\Cref{thm:hecke:self-adjoint}), not for these.
\end{warning}

\section{The integral structure: the perfect pairing}\label{sec:hecke:pairing}

This section proves the statement that the proof of Lemma lem:crit cites as \cite[Theorem 2.2]{Ribet83}, \cite[Lemma 2.1]{AU96}: ``The pairing $S_2(\Gamma_0(N);\Z)\times\T\to\Z$,
$(g,t)\mapsto a_1(tg)$, is perfect'', together with ``Every step holds verbatim over $\Z_\ell$, the pairing being perfect after tensoring with $\Z_\ell$'' and the pairing
$S_2(\Gamma_0(N);\Z_\ell)\times\T_\ell\to\Z_\ell$ of Step 1 of the proof of Theorem thm:quant. The only input from outside is the integral basis theorem of \Cref{ch:modforms}: $S_2(N;\Z)$
is a lattice spanning $S_2(N)$. Sturm's bound enters only in part (iv). Throughout the section $k=2$, $S=S_2(N;\Z)$ and $d=\dim S_2(N)$.

\begin{theorem}[the perfect pairing]\label{thm:hecke:perfect-pairing}
Define $\langle g,t\rangle=a_1(tg)$ for $g\in S$ and $t\in\T$.
\begin{enumerate}[label=(\roman*)]
\item $\langle\ ,\ \rangle$ is $\Z$-bilinear with values in $\Z$, and $\langle ug,t\rangle=\langle g,ut\rangle$ for $u\in\T$; moreover $\langle g,T_n\rangle=a_n(g)$.
\item The pairing is perfect: $S\to\Hom_\Z(\T,\Z)$, $g\mapsto\langle g,\cdot\rangle$, and $\T\to\Hom_\Z(S,\Z)$, $t\mapsto\langle\cdot,t\rangle$, are isomorphisms.
\item $\T$ is a free $\Z$-module of rank $d$; $S\simeq\Hom(\T,\Z)$ and $\T\simeq\Hom(S,\Z)$ as $\T$-modules; for every commutative ring $A$ the induced pairing $S_2(N;A)\times\T_A\to A$
is perfect, and $\T_A$ acts faithfully on $S_2(N;A)$. In particular this holds for $A=\Zl$.
\item Let $B$ be Sturm's bound of \Cref{ch:modforms}. Then $T_1,\dots,T_B$ span $\T$ as a $\Z$-module. In general $T_1,\dots,T_d$ do not.
\end{enumerate}
\end{theorem}

\begin{proof}
(i) $\T$ preserves $S$ (\Cref{prop:hecke:hecke-algebra}(iii)), so the values are integers; bilinearity is clear; $\langle ug,t\rangle=a_1(tug)=a_1(utg)=\langle g,ut\rangle$ by commutativity;
and $\langle g,T_n\rangle=a_1(T_ng)=a_n(g)$ by \Cref{thm:hecke:Tn}(iii).

(ii) \emph{Injectivity on the left.} If $\langle g,t\rangle=0$ for all $t$ then $a_n(g)=\langle g,T_n\rangle=0$ for all $n$ and $g=0$. \emph{Injectivity on the right.} If $\langle g,t\rangle=0$ for all $g\in S$,
then for every $g$ and $n$, $a_n(tg)=a_1(T_ntg)=\langle T_ng,t\rangle=0$, so $tg=0$; as $S$ spans $S_2(N)$, $t=0$. \emph{Ranks.} Tensoring with $\Q$, the two injections give
$\dim_\Q S_\Q\le\dim_\Q\T_\Q$ and $\dim_\Q\T_\Q\le\dim_\Q S_\Q$, where $S_\Q=S\otimes\Q$ has dimension $d$; so $\T$ has rank $d$ and both maps become isomorphisms over $\Q$.
\emph{Surjectivity on the left.} Let $\psi\colon\T\to\Z$ be additive. By the isomorphism over $\Q$ there is $g\in S_\Q$, a form with rational coefficients, with $a_1(tg)=\psi(t)$ for all
$t$. Then $a_n(g)=\psi(T_n)\in\Z$ for all $n$, so $g$ has integer coefficients: $g\in S$. Hence $S\to\Hom(\T,\Z)$ is an isomorphism. \emph{Surjectivity on the right.} $\T$ is free (a subgroup of the free group $\Mat_d(\Z)$,
\Cref{prop:hecke:hecke-algebra}(iii)) of rank $d$, so the composite $\T\to\Hom(S,\Z)\to\Hom(\Hom(\T,\Z),\Z)$, the second map
being the dual of the isomorphism just proved, is the canonical isomorphism of a free module of finite rank with its double dual (\Cref{ch:modules}); hence the first
map is an isomorphism.

(iii) The first statements are (ii) and (i) (the actions match because $\langle ug,t\rangle=\langle g,ut\rangle$). A perfect pairing of free $\Z$-modules of finite rank has Gram matrix in
$\GL_d(\Z)$ in any bases, and so does its extension to $A$; hence the pairing $(S\otimes A)\times(\T\otimes A)\to A$ is perfect. If $t\in\T_A$ kills $S_2(N;A)$, then $\langle g,t\rangle=0$ for
all $g$, so $t=0$.

(iv) By (ii) the image of $\T$ in $\Hom(S,\Z)$ is everything, and by (i) the image of $T_n$ is the functional $a_n$. \Cref{ch:modforms} shows (from Sturm's bound) that
the truncation $S\to\Z^B$, $g\mapsto(a_1(g),\dots,a_B(g))$, is injective with torsion-free cokernel, so it is a split injection and its dual $\Z^B\to\Hom(S,\Z)$ is onto: the functionals
$a_1,\dots,a_B$ span $\Hom(S,\Z)$. So $T_1,\dots,T_B$ span $\T$. The examples below show that $T_1,\dots,T_d$ need not.
\end{proof}

\begin{example}[the pairing at small levels]\label{ex:hecke:pairing-table}
For each $N$ below we computed a $\Z$-basis of $S$ by saturating the lattice of $q$-expansions (PARI's \texttt{matrixqz}, with coefficients up to $B\cdot\max(B,41)$), the
matrices of $T_1,\dots,T_B$ in that basis (all integral), the rank of their $\Z$-span (equal to $d$), the determinant of the pairing in $\Z$-bases (always $\pm1$, as
\Cref{thm:hecke:perfect-pairing} says) and the index of $\Z T_1+\cdots+\Z T_d$ in $\T$:
\begin{center}
\begin{tabular}{lcccccccccc}
\toprule
$N$ & 11 & 23 & 33 & 37 & 44 & 77 & 88 & 99 & 121 & 539\\
\midrule
$d=\rank\T$ & 1 & 2 & 3 & 2 & 4 & 7 & 9 & 9 & 6 & 49\\
$B$ & 2 & 4 & 8 & 6 & 12 & 16 & 24 & 24 & 22 & 112\\
$\det\langle g_i,t_j\rangle$ & 1 & 1 & $-1$ & $-1$ & $-1$ & $-1$ & $-1$ & $-1$ & $-1$ & $-1$\\
$[\T:\sum_{n\le d}\Z T_n]$ & 1 & 1 & 1 & 1 & 1 & 5 & 3 & 2 & $\infty$ & $\infty$\\
\bottomrule
\end{tabular}
\end{center}
Here $\infty$ means that $T_1,\dots,T_d$ are linearly dependent. At level $23$, $\T=\Z[T_2]$ with $T_2^2+T_2-1=0$, the ring of integers of $\Q(\sqrt5)$ (\Cref{ch:ant}).
\end{example}

\begin{theorem}[eigenforms are homomorphisms]\label{thm:hecke:eigen-dictionary}
Let $A$ be a commutative ring, and let $\T$ act on $\Hom_\Z(\T,A)$ by $(t\psi)(u)=\psi(tu)$.
\begin{enumerate}[label=(\roman*)]
\item The map $S_2(N;A)\to\Hom_\Z(\T,A)$, $g\mapsto(t\mapsto a_1(tg))$, is an isomorphism of $\T$-modules.
\item Under it the normalised eigenforms with coefficients in $A$ (elements $g$ with $a_1(g)=1$ and $tg\in Ag$ for all $t\in\T$) correspond bijectively to the ring homomorphisms
$\phi\colon\T\to A$, via $\phi\mapsto f_\phi:=\sum_n\phi(T_n)q^n$, and $tf_\phi=\phi(t)f_\phi$.
\item For $A=\C$ (or $\Qbar$): the normalised eigenforms in $S_2(N)$ are the $f_\phi$, $\phi\in\Hom(\T,\Qbar)$; there are finitely many, at most $d$. For $A=\Fbar_\ell$: the mod $\ell$
eigenforms $\sum\phi(T_n)q^n$ correspond to $\Hom(\T,\Fbar_\ell)$, the kernels of these homomorphisms are the maximal ideals of $\T$ of residue characteristic $\ell$, and a maximal
ideal $\mf m$ corresponds to the $[\T/\mf m:\F_\ell]$ homomorphisms $\T\to\T/\mf m\hookrightarrow\Fbar_\ell$, which are Frobenius conjugates. For $A=\mathcal O$, the ring of integers of a number field,
and $f$ a normalised eigenform with coefficients in $\mathcal O$: $\phi_f\colon\T\to\mathcal O$, and the kernel $\mf m_{f,\lambda}$ of its reduction $\T\to\mathcal O/\lambda$ is a maximal ideal of
residue characteristic $\ell$, \emph{the maximal ideal attached to $(f,\lambda)$}.
\end{enumerate}
\end{theorem}

\begin{proof}
(i) is \Cref{thm:hecke:perfect-pairing}(iii): $S\otimes A\simeq\Hom(\T,\Z)\otimes A=\Hom(\T,A)$, $\T$ being free of finite rank; the map is the one displayed. It is $\T$-linear by
\Cref{thm:hecke:perfect-pairing}(i). (ii) Let $g\leftrightarrow\psi$. If $g$ is a normalised eigenform with $tg=\phi(t)g$, then $\psi(t)=a_1(tg)=\phi(t)$, and $\phi$ is a ring homomorphism
(\Cref{prop:hecke:eigenforms}(iii) works over any $A$). Conversely, if $\psi$ is a ring homomorphism, then for $t,u\in\T$, $a_1(u\,tg)=\psi(ut)=\psi(t)\psi(u)=a_1(u\,\psi(t)g)$; by (i), $tg=\psi(t)g$, and
$a_1(g)=\psi(1)=1$. The $q$-expansion is $\sum\psi(T_n)q^n$ by \Cref{thm:hecke:perfect-pairing}(i). (iii) Ring homomorphisms $\T\to\Qbar$ are determined by the images of a $\Z$-basis of
$\T$, which are roots of their characteristic polynomials; there are at most $\dim\T_\Q=d$ of them (\Cref{ch:modules}). For $\Fbar_\ell$: the kernel of $\T\to\Fbar_\ell$ is a prime
ideal with finite quotient, hence maximal, of residue characteristic $\ell$; conversely $\T/\mf m$ is a finite field of characteristic $\ell$, which has $[\T/\mf m:\F_\ell]$ embeddings
into $\Fbar_\ell$, permuted transitively by Frobenius. The last statement is the case $\T\to\mathcal O\to\mathcal O/\lambda$.
\end{proof}

\begin{paperlink}\label{pl:hecke:max-ideal}
In Lemma lem:crit, ``the kernel of $\T\to\mathcal O/\lambda$, a maximal ideal of residue characteristic $\ell$'' is \Cref{thm:hecke:eigen-dictionary}(iii) for $\phi_g$; in Lemma lem:hecke,
``$\mf m$ the maximal ideal of $\T\otimes\Z_\ell$ containing $\ell$ and $\ker\phi_f$'' is $\mf m_{f,\ell}$ (for rational $f$, $\lambda=\ell$), viewed in $\T_\ell$ (\Cref{prop:hecke:local-components}).
\end{paperlink}

\begin{corollary}[consequences of the pairing]\label{cor:hecke:pairing-consequences}
In the notation of \Cref{thm:hecke:perfect-pairing}:
\begin{enumerate}[label=(\roman*)]
\item a form $g\in S_2(N;A)$ is determined by the values $a_1(tg)$, $t\in\T$;
\item for an ideal $\mf a\subseteq\T$, $S[\mf a]=\{g\in S:\mf ag=0\}$ corresponds to $\Hom(\T/\mf a,\Z)$; it is saturated in $S$, of rank $\rank_\Z(\T/\mf a)$;
\item for a ring homomorphism $\phi\colon\T\to\Z$, $S[\ker\phi]=\Z f_\phi$; in particular, for a newform $f$ with rational coefficients and $I=\ker\phi_f$, the forms in $S$ killed by $I$ are
$\Z f$;
\item if $I,J\subseteq\T$ are ideals with $I\cap J=0$, $\T/I$ and $\T/J$ torsion-free and $\T/(I+J)$ finite, then $S[I]\oplus S[J]$ has finite index in $S$, equal to $\#\T/(I+J)$.
\end{enumerate}
\end{corollary}

\begin{proof}
(i) is \Cref{thm:hecke:eigen-dictionary}(i). (ii) $\mf ag=0$ iff $\langle\mf ag,u\rangle=\langle g,\mf au\rangle=0$ for all $u$ (by (i)), iff $\langle g,\cdot\rangle$ vanishes on $\mf a\T=\mf a$. Saturation: if
$mg$ vanishes on $\mf a$ then so does $g$. The rank of $\Hom(\T/\mf a,\Z)$ is that of $\T/\mf a$. (iii) $\T/\ker\phi\simeq\Z$, so $S[\ker\phi]$ has rank one and corresponds to the functionals
$c\phi$, $c\in\Z$; the form attached to $\phi$ is $f_\phi$ (\Cref{thm:hecke:eigen-dictionary}(ii)). (iv) The map $\T\to\T/I\times\T/J$, $t\mapsto(t,t)$, is injective because $I\cap J=0$,
and its cokernel is $\T/(I+J)$ (send $(x,y)$ to $x-y$). Both sides are free (the quotients are torsion-free and finitely generated) and, the cokernel being finite, of the same
rank; so in bases the map is a square integer matrix $M$ with $\abs{\det M}=\#\T/(I+J)$ (Smith normal form, \Cref{ch:modules}). The dual map $\Hom(\T/I,\Z)\oplus\Hom(\T/J,\Z)\to\Hom(\T,\Z)$
has matrix $M^{\mathsf T}$, so its cokernel also has order $\abs{\det M}$. Under $S\simeq\Hom(\T,\Z)$ its image is $S[I]\oplus S[J]$, by (ii).
\end{proof}

The hypotheses of (iv) hold in Lemma lem:crit: $\T/I\simeq\Z$ for $I=\ker\phi_f$, and $J=\Ann_\T(f^\perp)$ is the kernel of the action on $f^\perp$, so $\T/J$ is a subring of $\End(f^\perp)$,
torsion-free; $I\cap J=0$ and the finiteness of $\T/(I+J)$ come from \Cref{thm:hecke:newform-splitting}. Torsion-freeness is what makes $\T/I\times\T/J$ free, so that the
determinant argument applies.

\begin{definition}[congruence ideal and congruence number]\label{def:hecke:congruence-ideal}
Let $\mathcal O$ be $\Z$, $\Z_{(\ell)}$, $\Zl$ or the completion $\mathcal O_\lambda$ of the ring of integers of a number field, and $\phi\colon\T_{\mathcal O}\to\mathcal O$ a surjective ring homomorphism. Its
\emph{(Hecke) congruence ideal} is $\eta(\phi)=\phi\bigl(\Ann_{\T_{\mathcal O}}(\ker\phi)\bigr)\subseteq\mathcal O$ (\Cref{ch:modules}). For a newform $f\in S_2(N)$ with rational coefficients, the
\emph{congruence number} $r_f$ is the largest integer $r$ such that $f\equiv g\pmod r$ for some $g\in S_2(N;\Z)$ orthogonal to $f$ for the Petersson product. For an elliptic curve
$E$ with newform $f_E$, $r_E=r_{f_E}$.
\end{definition}

\begin{proposition}\label{prop:hecke:rf-finite}
For a newform $f\in S_2(N)$ with rational coefficients, $L=f^\perp\cap S_2(N;\Z)$ is saturated of rank $d-1$, $\Z f\oplus L$ has finite index in $S_2(N;\Z)$, and $r_f$ is finite: a
congruence $f\equiv g\pmod r$ with $g\in L$ gives the element $(f-g)/r$ of $S_2(N;\Z)$ whose class in $S_2(N;\Z)/(\Z f\oplus L)$ has order $r$.
\end{proposition}

\begin{proof}
By \Cref{thm:hecke:newform-splitting} below (whose proof does not use this proposition), $f^\perp$ is the kernel of an idempotent $e\in\T\otimes\Q$; hence $f^\perp$ is spanned by forms with
rational coefficients, $f^\perp\cap S_2(N;\Q)$ has dimension $d-1$, and $L$, the intersection of a lattice with a subspace, is saturated of rank $d-1$. Also $f\notin f^\perp$, so
$\Z f\oplus L$ has rank $d$, hence finite index. If $f\equiv g\pmod r$ with $g\in L$, then $h=(f-g)/r\in S_2(N;\Z)$; if $mh\in\Z f\oplus L$ then, projecting to $\Q f$ along $f^\perp$, $m/r\in\Z$; so the
class of $h$ has order exactly $r$, and $r$ is bounded by the index.
\end{proof}

The equality $r_f\Z=\eta(\phi_f)$ is the paper's Lemma lem:crit, proved in \Cref{ch:modgalrep} from \Cref{thm:hecke:perfect-pairing}, \Cref{cor:hecke:pairing-consequences} and
\Cref{thm:hecke:newform-splitting}; its relation with the modular degree (the theorem of Agashe, Ribet and Stein) is \Cref{ch:moddeg}. The paper also uses $\eta(\pi)$ for the
eigencharacter $\pi$ of a depleted form $g^S$, which is not a newform (``$\eta^{\mathrm{Hecke}}(g^S)$'', Step 1 of the proof of Theorem thm:quant); \Cref{prop:hecke:complements} says
when the same arguments apply. Computed values of $r_f$ are in \Cref{ex:hecke:congruence-numbers}.

\begin{paperlink}\label{pl:hecke:crit-pairing}
The proof of Lemma lem:crit uses, in this order: the perfect pairing (\Cref{thm:hecke:perfect-pairing}); ``under it the forms killed by $I$ are $\Z f$'' (\Cref{cor:hecke:pairing-consequences}(iii));
``the forms killed by $J$ are $L=f^\perp\cap S_2(\Gamma_0(N);\Z)$'' (with $J=\Ann_\T(f^\perp)$: $J$ kills $L$ by definition, and a form killed by
$J$ is killed by the idempotent $e\in J\otimes\Q$ of \Cref{thm:hecke:newform-splitting}(ii), so lies in $\ker e=f^\perp$); ``Dualising the inclusion $\T\subset\T/I\times\T/J$
... shows that $\Z f\oplus L$ has index $r'$'' (\Cref{cor:hecke:pairing-consequences}(iv)); and ``Every step holds verbatim over $\Z_\ell$, the pairing being perfect after tensoring with
$\Z_\ell$'' (\Cref{thm:hecke:perfect-pairing}(iii)). The full proof is in \Cref{ch:modgalrep}.
\end{paperlink}

\section{Petersson adjoints}\label{sec:hecke:adjoints}

This section computes the adjoints of the Hecke operators for the Petersson product. Its use in the paper is the sentence of the proof of Lemma lem:crit: ``the
$f$-line is $\T$-stable, and so is $f^\perp$, because $f$ is an eigenvector of the adjoints of the $T_n$ as well''. For $(n,N)=1$ the operator $T_n$ is self-adjoint; $U_p$ is not,
but its adjoint is $w_NU_pw_N$, and a newform is an eigenvector of $w_N$. The cohomological analogue, with the cup product in place of the Petersson product, is the
statement of the paper's Section 4.1 that $w_{M^\Sigma}T_nw_{M^\Sigma}^{-1}$ is the adjoint of $T_n$ (\Cref{ch:moddeg}).

For a subgroup $\Gamma'\subseteq\SL_2(\Z)$ of finite index containing $-1$ and cusp forms $f,g$ of weight $k$ for $\Gamma'$ we write $\langle f,g\rangle_{\Gamma'}=\int_{\Gamma'\backslash\HH}f\bar gy^k\,d\mu$, $d\mu=dx\,dy/y^2$. We
use two facts from \Cref{ch:modforms}: the function $f\bar g\,y^k$ is $\Gamma'$-invariant, and for $\alpha\in\GL_2^+(\Q)$, $\langle f|_k\alpha,g|_k\alpha\rangle_{\alpha^{-1}\Gamma'\alpha}=\langle f,g\rangle_{\Gamma'}$ (change of
variables $z\mapsto\alpha z$, using $\operatorname{Im}(\alpha z)^{k/2}\abs{f(\alpha z)}=(\operatorname{Im}z)^{k/2}\abs{(f|_k\alpha)(z)}$), where $\alpha^{-1}\Gamma'\alpha$ need not lie in $\SL_2(\Z)$ but is commensurable with it.

\begin{proposition}[adjoints of double coset operators]\label{prop:hecke:double-coset-adjoint}
Let $\Gamma=\Gamma_0(N)$, $\alpha\in\GL_2^+(\Q)$ and $\alpha'=(\det\alpha)\alpha^{-1}$. For $f,g\in S_k(N)$,
\[
\langle[\Gamma\alpha\Gamma]_kf,g\rangle=\langle f,[\Gamma\alpha'\Gamma]_kg\rangle.
\]
\end{proposition}

\begin{proof}
Let $\Gamma_\alpha=\Gamma\cap\alpha^{-1}\Gamma\alpha$ and $\Gamma=\bigsqcup_{i=1}^r\Gamma_\alpha\gamma_i$, so that $\Gamma\alpha\Gamma=\bigsqcup\Gamma\alpha\gamma_i$ (\Cref{prop:hecke:double-coset}(i)). The function $f|_k\alpha$ is
invariant under $\alpha^{-1}\Gamma\alpha\supseteq\Gamma_\alpha$, so $F(z)=(f|_k\alpha)(z)\overline{g(z)}y^k$ is $\Gamma_\alpha$-invariant, and since $g|_k\gamma_i=g$, the $i$-th term of $[\Gamma\alpha\Gamma]_kf\cdot\bar g\,y^k$ is
$(\det\alpha)^{k/2-1}F(\gamma_iz)$. If $\mathcal F$ is a fundamental domain for $\Gamma$, then $\bigcup_i\gamma_i\mathcal F$ is one for $\Gamma_\alpha$ (both groups contain $-1$), so
\[
\langle[\Gamma\alpha\Gamma]_kf,g\rangle=(\det\alpha)^{k/2-1}\sum_i\int_{\mathcal F}F(\gamma_iz)\,d\mu=(\det\alpha)^{k/2-1}\langle f|_k\alpha,g\rangle_{\Gamma_\alpha}.
\]
Now $\alpha'\alpha=(\det\alpha)\cdot1$, and positive scalars act trivially in even weight, so $g=(g|_k\alpha')|_k\alpha$. By the change of variables,
$\langle f|_k\alpha,(g|_k\alpha')|_k\alpha\rangle_{\Gamma_\alpha}=\langle f,g|_k\alpha'\rangle_{\alpha\Gamma_\alpha\alpha^{-1}}$, and $\alpha\Gamma_\alpha\alpha^{-1}=\alpha\Gamma\alpha^{-1}\cap\Gamma=\Gamma\cap\alpha'^{-1}\Gamma\alpha'=\Gamma_{\alpha'}$ (conjugation by a scalar is
trivial). Applying the first computation to $\alpha'$ and $g$, $\langle[\Gamma\alpha'\Gamma]_kg,f\rangle=(\det\alpha')^{k/2-1}\langle g|_k\alpha',f\rangle_{\Gamma_{\alpha'}}$. Since $\det\alpha'=\det\alpha$ is real, taking complex
conjugates gives $\langle f,[\Gamma\alpha'\Gamma]_kg\rangle=(\det\alpha)^{k/2-1}\langle f,g|_k\alpha'\rangle_{\Gamma_{\alpha'}}=\langle[\Gamma\alpha\Gamma]_kf,g\rangle$.
\end{proof}

\begin{theorem}[self-adjointness]\label{thm:hecke:self-adjoint}
For $(n,N)=1$ the operator $T_n$ is self-adjoint on $S_k(N)$. Consequently $\T^{(N)}\otimes\C$ is a commutative algebra of self-adjoint operators: the eigenvalues of every $T_n$,
$(n,N)=1$, are real; $S_k(N)$ is the orthogonal direct sum of the simultaneous eigenspaces of $\T^{(N)}$; and the generalised eigenspaces of $\T^{(N)}$ are eigenspaces.
\end{theorem}

\begin{proof}
For a prime $p\nmid N$, $\sm100p'=\sm p001$, and $\Gamma\sm p001\Gamma=\Gamma\sm100p\Gamma$ (\Cref{prop:hecke:coset-decomposition}(ii)); so $T_p^\dagger=T_p$ by
\Cref{prop:hecke:double-coset-adjoint}. For $(n,N)=1$, $T_n$ is a polynomial with integer coefficients in the $T_p$, $p\mid n$ (\Cref{thm:hecke:Tn}), and a real polynomial in
commuting self-adjoint operators is self-adjoint. The rest is the spectral theorem for a finite commuting family of self-adjoint operators on a finite-dimensional Hermitian
space: each is diagonalisable with real eigenvalues in an orthonormal basis, and commuting ones are simultaneously so; eigenspaces for distinct systems are orthogonal
because $\lambda\langle v,w\rangle=\langle tv,w\rangle=\langle v,tw\rangle=\mu\langle v,w\rangle$ with $\lambda\ne\mu$ real.
\end{proof}

\begin{lemma}[the trace from level $Mp$ to level $M$]\label{lem:hecke:trace}
Let $p\nmid M$, $\Gamma_0(M)=\bigsqcup_{i=0}^{p}\Gamma_0(Mp)\gamma_i$, and $\Tr(g)=\sum_ig|_k\gamma_i$ for $g\in S_k(Mp)$. Then $\Tr(g)\in S_k(M)$, and for $h\in S_k(M)$:
(i) $\langle h,g\rangle_{\Gamma_0(Mp)}=\langle h,\Tr g\rangle_{\Gamma_0(M)}$; (ii) $\Tr(h(pz))=p^{1-k}T_p^{(M)}h$; (iii) $\langle h,h\rangle_{\Gamma_0(Mp)}=(p+1)\langle h,h\rangle_{\Gamma_0(M)}$.
\end{lemma}

\begin{proof}
$\Tr(g)$ is $\Gamma_0(M)$-invariant because right multiplication by $\gamma\in\Gamma_0(M)$ permutes the cosets, and cuspidal by the argument of \Cref{prop:hecke:double-coset}. (i) As
in the proof of \Cref{prop:hecke:double-coset-adjoint}, $\bigcup\gamma_i\mathcal F$ is a fundamental domain for $\Gamma_0(Mp)$ when $\mathcal F$ is one for $\Gamma_0(M)$, and $h|_k\gamma_i=h$. (ii) Let
$\alpha=\sm p001$; then $\Gamma_0(M)\cap\alpha^{-1}\Gamma_0(M)\alpha=\Gamma_0(Mp)$ (conjugation by $\alpha$ sends $\sm abcd$ to $\sm{a}{pb}{c/p}d$), so by \Cref{prop:hecke:double-coset}(i) and
\Cref{prop:hecke:coset-decomposition}(ii), $p^{k/2-1}\sum_ih|_k\alpha\gamma_i=[\Gamma_0(M)\alpha\Gamma_0(M)]_kh=T_p^{(M)}h$. Since $h(pz)=p^{-k/2}h|_k\alpha$, $\Tr(h(pz))=p^{-k/2}p^{1-k/2}T_p^{(M)}h$. (iii) is (i)
with $g=h$, since $\Tr h=(p+1)h$.
\end{proof}

\begin{proposition}[the adjoint of $U_p$]\label{prop:hecke:Up-adjoint}
Let $p\mid N$.
\begin{enumerate}[label=(\roman*)]
\item $U_p^\dagger=w_NU_pw_N$ on $S_k(N)$.
\item Let $h$ be a newform of level $M$ and $p\nmid M$. In $S_k(Mp)$, with all products at level $Mp$,
\[
\langle h(pz),h\rangle=\frac{p^{1-k}a_p(h)}{p+1}\langle h,h\rangle,\qquad\langle h(pz),h(pz)\rangle=p^{-k}\langle h,h\rangle.
\]
\item $U_p$ is in general neither self-adjoint nor normal (\Cref{ex:hecke:level-22}).
\end{enumerate}
\end{proposition}

\begin{proof}
(i) With $W=W_N$, a direct computation gives $W\sm100pW^{-1}=\sm p001$, and $W$ normalises $\Gamma$ (\Cref{ch:modforms}). Hence $\Gamma\sm p001\Gamma=W(\Gamma\sm100p\Gamma)W^{-1}=\bigsqcup\Gamma W\alpha_iW^{-1}$, and
$[\Gamma\sm p001\Gamma]_kg=\sum_ig|_kW\alpha_iW^{-1}=w_N^{-1}U_pw_Ng$ (with $w_Ng=g|_kW$). As $w_N^{-1}=w_N$ (\Cref{conv:hecke:standing}(f)), \Cref{prop:hecke:double-coset-adjoint} with $\alpha=\sm100p$,
$\alpha'=\sm p001$ gives $U_p^\dagger=w_NU_pw_N$.

(ii) By \Cref{lem:hecke:trace}, $\langle h(pz),h\rangle_{Mp}=\langle\Tr(h(pz)),h\rangle_M=p^{1-k}a_p(h)\langle h,h\rangle_M$ and $\langle h,h\rangle_{Mp}=(p+1)\langle h,h\rangle_M$ (the first equality is the complex conjugate of \Cref{lem:hecke:trace}(i)). For the second formula, $h(pz)=p^{-k/2}h|_k\alpha$ with $\alpha=\sm p001$, and
$\alpha^{-1}\Gamma_0(M)\alpha$ contains $\Gamma_0(Mp)$ with the same index $p+1$ as $\Gamma_0(Mp)$ in $\Gamma_0(M)$ (the two groups have the same covolume); so
$\langle h|_k\alpha,h|_k\alpha\rangle_{\Gamma_0(Mp)}=(p+1)\langle h|_k\alpha,h|_k\alpha\rangle_{\alpha^{-1}\Gamma_0(M)\alpha}=(p+1)\langle h,h\rangle_{\Gamma_0(M)}=\langle h,h\rangle_{\Gamma_0(Mp)}$, and the factor $p^{-k}$ comes from $(p^{-k/2})^2$.
\end{proof}

\begin{example}[level $22$]\label{ex:hecke:level-22}
Take $h=f_{11}$, $p=2$, $N=22$, and the basis $f_{11}(z),f_{11}(2z)$ of $S_2(22)$. By \Cref{prop:hecke:Up-adjoint}(ii) with $a_2(f_{11})=-2$, the Gram matrix of the Petersson product is
$\langle f_{11},f_{11}\rangle_{22}\cdot G$, $G=\sm1{-1/3}{-1/3}{1/4}$. The matrix of $U_2$ (columns: images of the basis vectors) is $A=\sm{-2}1{-2}0$, since $U_2f_{11}=-2f_{11}-2f_{11}(2z)$ and
$U_2f_{11}(2z)=f_{11}$. With $\langle x,y\rangle=x^{\mathsf T}G\bar y$ in coordinates, the adjoint has matrix $G^{-1}A^{\mathsf T}G=\sm0{-1/2}4{-2}$. The involution $w_{22}$ exchanges the two lines:
$f_{11}|_2W_{22}=-2f_{11}(2z)$ and $f_{11}(2z)|_2W_{22}=-\frac12f_{11}$ (\Cref{ch:modforms}, from $\varepsilon_{11}(f_{11})=-1$), so $w_{22}$ has matrix $\sm0{-1/2}{-2}0$, and indeed
$\sm0{-1/2}{-2}0\sm{-2}1{-2}0\sm0{-1/2}{-2}0=\sm0{-1/2}4{-2}$. Finally $AA^\dagger=\sm4{-1}01\ne\sm10{-4}4=A^\dagger A$: $U_2$ is not normal. (The Gram matrix was also computed with PARI's
\texttt{mfpetersson}, to $40$ digits.)
\end{example}

\begin{corollary}[stable orthogonal complements]\label{cor:hecke:perp-stable}
Let $V\subseteq S_k(N)$ be a subspace stable under $\T$ and under $w_N$. Then $V^\perp$ is stable under $\T$ and $w_N$. In particular, if $f$ is an eigenform of $\T$ and an
eigenvector of $w_N$, then $f^\perp$ is $\T$-stable.
\end{corollary}

\begin{proof}
For $g\in V^\perp$ and $v\in V$: $\langle T_ng,v\rangle=\langle g,T_n^\dagger v\rangle$ with $T_n^\dagger$ equal to $T_n$ ($(n,N)=1$) or to a product of $w_NU_pw_N$ and such $T_m$ (\Cref{thm:hecke:Tn},
\Cref{thm:hecke:self-adjoint}, \Cref{prop:hecke:Up-adjoint}); in both cases $T_n^\dagger V\subseteq V$, so the product vanishes. Also $w_N$ is self-adjoint (\Cref{ch:modforms}).
\end{proof}

\begin{paperlink}\label{pl:hecke:adjoints}
``The $f$-line is $\T$-stable, and so is $f^\perp$, because $f$ is an eigenvector of the adjoints of the $T_n$ as well'' (proof of Lemma lem:crit): by \Cref{prop:hecke:Up-adjoint}(i),
$U_p^\dagger f=w_NU_pw_Nf=\varepsilon_N(f)^2a_p(f)f=a_p(f)f$, using $f|_2W_N=\varepsilon_N(f)f$ (\Cref{prop:hecke:newform-local}); this is \Cref{cor:hecke:perp-stable} with $V=\C f$. A second proof,
which needs no adjoints of $U_p$, is in \Cref{thm:hecke:newform-splitting}: $f^\perp$ is the sum of the eigenspaces of $\T^{(N)}$ other than the line of $f$.
\end{paperlink}

\section{Oldforms and newforms: the theory of Atkin and Lehner}\label{sec:hecke:newforms}

This section explains what ``the newform $f=f_E$'' of the paper's Section 1 is, and why the forms $h(dz)$ of Lemma lem:old span a subspace orthogonal to $f$. It
states the theorems of Atkin and Lehner (newform bases, multiplicity one, the decomposition of $S_k(N)$) and strong multiplicity one, with sketches and references; and it
proves the consequences the paper uses: $f$ is an eigenvector of $w_N$, its coefficients are real, its Galois conjugates are newforms, and the dimension of the new
subspace.

\begin{definition}[oldforms and newforms]\label{def:hecke:old-new}
For $M\mid N$ and $d\mid N/M$, $V_d\colon S_k(M)\to S_k(N)$, $h\mapsto h(dz)=d^{-k/2}h|_k\sm d001$ (\Cref{ch:modforms}). The \emph{old subspace} is
\[
S_k^{\mathrm{old}}(N)=\sum_{p\mid N}\bigl(S_k(N/p)+V_pS_k(N/p)\bigr),
\]
equivalently the span of all $V_dS_k(M)$ with $M\mid N$, $M\ne N$ and $d\mid N/M$. The \emph{new subspace} $S_k^{\mathrm{new}}(N)$ is its orthogonal complement for the Petersson product. A
\emph{newform of level $N$} is a normalised eigenform of $\T_N$ lying in $S_k^{\mathrm{new}}(N)$.
\end{definition}

The two descriptions of the old subspace agree: $S_k(N/p)+V_pS_k(N/p)$ is of the second kind, and conversely, for $M\ne N$, $d\mid N/M$ and $h\in S_k(M)$, either $d=1$ and
$h\in S_k(N/p)$ for any $p\mid N/M$, or $p\mid d$ for some prime and $V_dh=V_p(V_{d/p}h)$ with $V_{d/p}h\in S_k(Md/p)\subseteq S_k(N/p)$.

\begin{proposition}\label{prop:hecke:old-new}
\begin{enumerate}[label=(\roman*)]
\item For $(n,N)=1$, $M\mid N$, $d\mid N/M$ and $h\in S_k(M)$: $T_n^{(N)}(h(dz))=(T_n^{(M)}h)(dz)$.
\item For $h\in S_k(M)$ and $d\mid N/M$: $(h(dz))|_kW_N=d^{-k/2}e^{k/2}(h|_kW_M)(ez)$ with $e=N/(Md)$. So $w_N$ maps $V_dS_k(M)$ onto $V_{N/(Md)}S_k(M)$.
\item $S_k^{\mathrm{old}}(N)$ and $S_k^{\mathrm{new}}(N)$ are stable under all of $\T$ (the $U_p$ included) and under $w_N$.
\end{enumerate}
\end{proposition}

\begin{proof}
(i) $T_n$ is a polynomial in the $T_p$, $p\mid n$, each equal to $\mathcal U_p+p^{k-1}V_p$ at both levels, and these commute with $V_d$ for $(d,p)=1$ (\Cref{lem:hecke:UV}). (ii) A direct
computation gives $\sm d001W_N=d\cdot W_M\sm e001$, and positive scalars act trivially in even weight; so $(h|_k\sm d001)|_kW_N=(h|_kW_M)|_k\sm e001$, and $h|_k\sm c001=c^{k/2}h(cz)$. Since $w_M$
is a bijection of $S_k(M)$, the image of $V_dS_k(M)$ is $V_eS_k(M)$. (iii) The old subspace is stable under $T_n$, $(n,N)=1$, by (i), and under $w_N$ by (ii). For $p\mid N$, $U_p=\mathcal U_p$ on
$q$-expansions (\Cref{prop:hecke:Tp-q-expansion}); for $h\in S_k(M)$ and $d\mid N/M$: if $p\mid d$, $\mathcal U_pV_dh=V_{d/p}h$; if $p\nmid d$, $\mathcal U_pV_dh=V_d\mathcal U_ph$, and $\mathcal U_ph$ is $U_p^{(M)}h$ if $p\mid M$, or
$T_p^{(M)}h-p^{k-1}V_ph$ if $p\nmid M$ (then $p\mid N/M$ and $V_{dp}h$ is again old, as $dp\mid N/M$). So the old subspace is $\T$-stable, and the new subspace, its orthogonal complement, is
$\T$-stable and $w_N$-stable by \Cref{cor:hecke:perp-stable}.
\end{proof}

\begin{theorem}[Atkin, Lehner, Li]\label{thm:hecke:atkin-lehner}
Let $k\ge2$ be even.
\begin{enumerate}[label=(\roman*)]
\item (Newform basis.) $S_k^{\mathrm{new}}(N)$ has a basis of newforms; distinct newforms are orthogonal.
\item (Multiplicity one.) Let $f$ be a newform of level $M$ and $L$ a multiple of $M$. The subspace of $g\in S_k(L)$ with $T_ng=a_n(f)g$ for all $(n,L)=1$ is $\operatorname{span}\{f(dz):d\mid L/M\}$, of
dimension $\sigma_0(L/M)$. In particular two newforms of levels dividing $L$ with the same eigenvalues for all $T_n$, $(n,L)=1$, are equal.
\item (Decomposition.) $S_k(N)=\bigoplus_{M\mid N}\bigoplus_{d\mid N/M}V_dS_k^{\mathrm{new}}(M)$: the forms $h(dz)$, $h$ a newform of level $M\mid N$ and $d\mid N/M$, form a basis of $S_k(N)$,
and the old subspace is spanned by those with $M\ne N$.
\item For a newform $f$ of level $M\mid N$ write $V_f=\operatorname{span}\{f(dz):d\mid N/M\}\subseteq S_k(N)$. Then $S_k(N)=\bigoplus_fV_f$ orthogonally, $V_f$ is the eigenspace of $\T^{(N)}$ for the
system of $f$, and for a newform $f$ of level $N$, $V_f=\C f$ and $f^\perp=\bigoplus_{h\ne f}V_h$.
\end{enumerate}
We do not prove (i) to (iii); see \cite{AtkinLehner70} (for trivial character, the main theorems), \cite{Li75} (any character), \cite[Sections 5.6 to 5.8]{DiamondShurman} and
\cite[Section 4.6]{Miyake}. Part (iv) follows: by (ii) the spaces $V_f$ are the eigenspaces of $\T^{(N)}$ for pairwise distinct systems, and they are orthogonal by
\Cref{thm:hecke:self-adjoint}.
\end{theorem}

\begin{proof}[Sketch of proof]
The route of Diamond and Shurman. The key is the \emph{Main Lemma}: if $g\in S_k(N)$ has $a_n(g)=0$ for every $n$ prime to $N$, then $g\in\sum_{p\mid N}V_pS_k(N/p)$. Granting it, let $g\in
S_k^{\mathrm{new}}(N)$ be an anemic eigenform with $T_ng=\lambda_ng$ for $(n,N)=1$. If $a_1(g)=0$, then $a_n(g)=a_1(T_ng)=0$ for $(n,N)=1$, so $g$ is old by the Main Lemma, hence $g=0$. So every
nonzero anemic eigenform in the new subspace has $a_1\ne0$; two of them with the same system, normalised, have a difference with $a_1=0$, which is $0$; this is multiplicity one
in $S_k^{\mathrm{new}}(N)$. As $\T$ preserves $S_k^{\mathrm{new}}(N)$ and commutes with $\T^{(N)}$, each anemic eigenspace there is a line stable under $\T$, so its normalised generator is
an eigenform of all of $\T$: a newform; the anemic eigenspaces exhaust $S_k^{\mathrm{new}}(N)$ by \Cref{thm:hecke:self-adjoint}, which gives (i). Part (iii) follows by induction on $N$,
applying the same argument to the old subspace, and the statement for other levels in (ii) needs in addition that newforms of different levels have different anemic
systems, which comes out of the same induction. The Main Lemma itself is proved by an induction on the number of prime factors of $N$, using the operators $\mathcal U_p$, $V_p$
and the trace of \Cref{lem:hecke:trace} \cite[Section 5.7]{DiamondShurman}.
\end{proof}

\begin{theorem}[strong multiplicity one]\label{thm:hecke:strong-multiplicity-one}
\begin{enumerate}[label=(\roman*)]
\item If $f$ and $g$ are newforms of weight $k$ and levels $M$, $M'$ with $a_p(f)=a_p(g)$ for all but finitely many primes $p$, then $f=g$ (and $M=M'$).
\item (Weight $2$.) The same holds if $a_p(f)=a_p(g)$ for all $p$ in a set of primes of Dirichlet density one.
\item Let $f$ be a newform of level $M$, $L$ a multiple of $M$, and $h\in S_k(L)$ an eigenvector of $\T^{(L)}$ for a system $\psi$ with $\psi(T_p)\ne a_p(f)$ for infinitely many $p$. Then $h$ is
orthogonal to every $f(dz)$, $d\mid L/M$. If $h'$ is a newform of level dividing $L$ different from $f$, then its system satisfies this hypothesis.
\end{enumerate}
Part (i) is stated with a reference: \cite[Theorem 4.6.19]{Miyake} (the theorem number to be verified), \cite{Li75}; it is a consequence of (ii) of \Cref{thm:hecke:atkin-lehner}
applied at the level $L=\lcm(M,M')$ once one knows that finitely many $T_p$ may be ignored, which is the content of the theorem. Part (ii) is sketched below; (iii) is proved here.
\end{theorem}

\begin{proof}[Sketch of proof of (ii), and proof of (iii)]
(ii) Let $\lambda$ be a prime above some $\ell$ of a number field containing the coefficients of $f$ and $g$. Their $\lambda$-adic Galois representations (\Cref{ch:modgalrep}) are unramified
outside $\ell MM'$ and satisfy $\tr\rho_f(\Frob_p)=a_p(f)$ there. By Chebotarev (\Cref{ch:galois}) the Frobenius elements at a set of primes of density one are dense in the Galois group, so
$\tr\rho_f=\tr\rho_g$ everywhere by continuity; by Brauer-Nesbitt (\Cref{ch:reps}) $\rho_f^{\mathrm{ss}}\simeq\rho_g^{\mathrm{ss}}$, so $a_p(f)=a_p(g)$ for all $p\nmid\ell MM'$, and (i) applies.
(iii) By \Cref{thm:hecke:self-adjoint} $h$ and each $f(dz)$ are eigenvectors of the self-adjoint operators $T_p$, $p\nmid L$, with eigenvalues $\psi(T_p)$ and $a_p(f)$ (\Cref{prop:hecke:old-new}(i));
these differ for some $p\nmid L$, since they differ for infinitely many $p$. Eigenvectors of a self-adjoint operator for distinct eigenvalues are orthogonal. The last sentence is (i).
\end{proof}

\begin{paperlink}\label{pl:hecke:smo}
Strong multiplicity one is used in three places. In the proof of Theorem thm:qual (b) and (c): ``$g=f'$ by strong multiplicity one'', two newforms with the same $a_p$ for $p\nmid qN$
(\Cref{thm:hecke:strong-multiplicity-one}(i)). In the proof of Theorem thm:twist: ``a system of $T_p$-eigenvalues ... differing from that of $g$ at infinitely many $p$'', so that $\iota(h)\in f^\perp$
(\Cref{thm:hecke:strong-multiplicity-one}(iii) and \Cref{prop:hecke:twist-transport}). In Step 1 of the proof of Theorem thm:quant: ``strong multiplicity one identifies the forms with
the $T_p$-eigenvalues of $g$, $p\nmid N$, with the span of the $g(dz)$, $d\mid N/N_S$'', which is \Cref{thm:hecke:atkin-lehner}(ii) with $L=N$, together with (i) to exclude newforms of
other levels. In Lemma lem:old, ``orthogonal to the newform $f$ (whether $M<N$, or $M=N$ and $h\ne f$)'' is \Cref{thm:hecke:atkin-lehner}(iv).
\end{paperlink}

\begin{proposition}[local behaviour of newforms]\label{prop:hecke:newform-local}
Let $f$ be a newform of weight $k$ and level $N$.
\begin{enumerate}[label=(\roman*)]
\item $f|_kW_N=\varepsilon_N(f)f$ with $\varepsilon_N(f)=\pm1$. More generally (Atkin-Lehner) $f$ is an eigenvector of every $w_d$, $d\,\|\,N$, with eigenvalue $\varepsilon_d(f)=\pm1$.
\item $a_p(f)$ is real for every $p\mid N$; hence all $a_n(f)$ are real.
\item If $p^2\mid N$ then $a_p(f)=0$; if $p\,\|\,N$ then $a_p(f)=-p^{k/2-1}\varepsilon_p(f)$, so in weight $2$, $a_p(f)=-\varepsilon_p(f)\in\{\pm1\}$.
\item For $f=f_E$, the newform of an elliptic curve of conductor $N$ (\Cref{ch:moddeg}), $a_p(f)=1$, $-1$, $0$ for split multiplicative, nonsplit multiplicative, additive reduction at
$p$ (\Cref{ch:ellcurves}).
\end{enumerate}
Part (i) for $w_N$ and part (ii) are proved here given \Cref{thm:hecke:atkin-lehner}; (i) for general $d$ and (iii) are stated with a reference \cite{AtkinLehner70}, as cited by the
paper (``a newform on $\Gamma_0(M)$ with $p^2\mid M$ has $a_p=0$ \cite[Theorem 3]{AtkinLehner70}'', proof of Lemma lem:old), and \cite[Sections 5.8, 6.6]{DiamondShurman}.
\end{proposition}

\begin{proof}
(i) For $p\nmid N$, $w_NT_pw_N^{-1}=[\Gamma\sm p001\Gamma]_k=T_p$ by the computation in the proof of \Cref{prop:hecke:Up-adjoint}(i) and \Cref{prop:hecke:coset-decomposition}(ii); so $w_N$ commutes
with $\T^{(N)}$. By \Cref{prop:hecke:old-new}(iii) $w_Nf$ lies in $S_k^{\mathrm{new}}(N)$, in the eigenspace of $\T^{(N)}$ for the system of $f$, which is $\C f$ by \Cref{thm:hecke:atkin-lehner}(ii). So
$w_Nf=\varepsilon f$, and $\varepsilon^2=1$ because $w_N$ is an involution. (ii) $a_p\langle f,f\rangle=\langle U_pf,f\rangle=\langle f,w_NU_pw_Nf\rangle=\langle f,\varepsilon^2a_pf\rangle=\bar a_p\langle f,f\rangle$ by \Cref{prop:hecke:Up-adjoint}(i). The
$a_p$, $p\nmid N$, are real by \Cref{thm:hecke:self-adjoint}, and the recursion and multiplicativity (\Cref{prop:hecke:eigenforms}(ii)) give the rest. (iv) is the comparison of
\Cref{ch:moddeg} (modularity, $a_p(f_E)=a_p(E)$ at all $p$) with the values of $a_p(E)$ at bad primes in \Cref{ch:ellcurves}.
\end{proof}

\begin{proof}[Sketch of proof of (iii)]
For $p\,\|\,N$: the trace of $f$ to level $N/p$ vanishes, since $\langle\Tr f,h\rangle_{N/p}=\langle f,h\rangle_N=0$ for every $h\in S_k(N/p)$ (\Cref{lem:hecke:trace}(i), $f$ being new). A coset computation
like that of \Cref{lem:hecke:trace} expresses this trace as $\Tr f=f+p^{1-k/2}U_p(f|_kW_p)$ (\cite{AtkinLehner70}, \cite{Li75}; we do not reproduce the computation); with
$f|_kW_p=\varepsilon_pf$ this gives $(1+p^{1-k/2}\varepsilon_pa_p)f=0$. Check: $f_{11}$, $p=11$, $k=2$: $\varepsilon_{11}=-1$, $a_{11}=1$. For $p^2\mid N$ we refer to \cite[Theorem 3]{AtkinLehner70}. A conceptual reason, through
\Cref{ch:modgalrep}: the local component of $f$ at $p$ has trivial central character and conductor exponent at least $2$, so it is a ramified principal series, a ramified
twist of the Steinberg representation or supercuspidal, and in each case its local $L$-factor is $1$; as $L_p(f,s)=(1-a_pp^{-s})^{-1}$ (\Cref{prop:hecke:euler-product}), $a_p=0$.
\end{proof}

Numerical checks of (iii) and (iv) (PARI, \texttt{ellap} on Cremona's curves): $a_{11}(\curve{11a1})=1$ with $\varepsilon_{11}=-1$; $a_{37}(\curve{37a1})=-1$ with $\varepsilon_{37}=+1$ (root number
$-\varepsilon_{37}=-1$, rank $1$); $a_{11}=0$ for \curve{121a1}, \curve{121b1}, \curve{121c1}, \curve{121d1}; $a_7=0$ for \curve{539c1}, \curve{539d1}; $a_2=0$ for \curve{44a1}, \curve{88a1}, \curve{20a1};
$a_7(\curve{77c1})=-1$, $a_{11}(\curve{77c1})=1$, $a_5(\curve{20a1})=-1$.

\begin{corollary}[Galois conjugates of newforms]\label{cor:hecke:galois-newforms}
For $g\in S_k(N)$ and $\sigma\in\Aut(\C)$ let $g^\sigma$ be defined as in \Cref{prop:hecke:coefficient-field}(ii) (apply $\sigma$ to the coordinates in a basis of forms with rational
coefficients; in weight $2$, $a_n(g^\sigma)=\sigma(a_n(g))$).
\begin{enumerate}[label=(\roman*)]
\item $S_k^{\mathrm{old}}(N)$ and $S_k^{\mathrm{new}}(N)$ are stable under $g\mapsto g^\sigma$, and both are spanned by forms with rational coefficients.
\item If $f$ is a newform of level $N$, so is $f^\sigma$.
\item If $f$ is a newform, $K_f=\Q(a_n(f):n\ge1)$ is totally real.
\end{enumerate}
\end{corollary}

\begin{proof}
We give the proof in weight $2$, where $S_2(N;\Z)$ is available; the general case is the same with a rational basis of $S_k(N)$ (\Cref{ch:modforms}). (i) Each $S_2(N/p)$ has a basis of
forms with integer coefficients, and $V_p$ preserves integrality; so the old subspace is spanned by forms with rational coefficients and is $\sigma$-stable. By
\Cref{thm:hecke:atkin-lehner}(iii) and (ii), $S_2(N)=S_2^{\mathrm{old}}(N)\oplus S_2^{\mathrm{new}}(N)$, the systems of $\T^{(N)}$ occurring in the two summands are disjoint, and the new subspace is the
sum of the eigenspaces of the systems that do not occur in the old one. The map $g\mapsto g^\sigma$ sends the eigenspace of a system $\psi$ to that of $\sigma\circ\psi$ (as in
\Cref{prop:hecke:coefficient-field}(ii)) and preserves the old subspace; so it permutes the systems occurring in the old subspace, hence also the others, and preserves the new
subspace. A subspace $W$ of $\C^d$ (coordinates in a $\Z$-basis of $S_2(N;\Z)$) stable under every $\sigma$ has a reduced row echelon basis, unique, hence fixed by every $\sigma$, hence
rational. (ii) $f^\sigma$ is a normalised eigenform (\Cref{prop:hecke:coefficient-field}(ii)) in the new subspace. (iii) An embedding $\tau\colon K_f\to\C$ extends to some $\sigma\in\Aut(\C)$, and
$\tau(a_n(f))=a_n(f^\sigma)$ is real by (ii) and \Cref{prop:hecke:newform-local}(ii).
\end{proof}

\begin{corollary}[dimensions]\label{cor:hecke:dimensions}
\[
\dim S_k(N)=\sum_{M\mid N}\sigma_0(N/M)\dim S_k^{\mathrm{new}}(M),\qquad\dim S_k^{\mathrm{new}}(N)=\sum_{M\mid N}\beta(N/M)\dim S_k(M),
\]
where $\beta=\mu*\mu$ is the multiplicative function with $\beta(p)=-2$, $\beta(p^2)=1$ and $\beta(p^r)=0$ for $r\ge3$ ($\mu$ the Moebius function).
\end{corollary}

\begin{proof}
The first formula counts the basis of \Cref{thm:hecke:atkin-lehner}(iii). It says $D=\sigma_0*D^{\mathrm{new}}=\mathbf1*\mathbf1*D^{\mathrm{new}}$ for the arithmetic functions $D(N)=\dim S_k(N)$ and
$D^{\mathrm{new}}(N)=\dim S_k^{\mathrm{new}}(N)$; Moebius inversion twice gives $D^{\mathrm{new}}=\mu*\mu*D$ (\Cref{ch:rings}). $\beta$ is multiplicative, and $\beta(p^r)=\sum_{i+j=r}\mu(p^i)\mu(p^j)$ gives the values.
\end{proof}

\begin{example}[dimensions in weight $2$]\label{ex:hecke:dimensions}
Computed with PARI's \texttt{mfdim} for the full and the new space, and checked against \Cref{cor:hecke:dimensions}:
\begin{center}
\begin{tabular}{lcccccccccccc}
\toprule
$N$ & 11 & 22 & 33 & 44 & 77 & 88 & 99 & 121 & 242 & 539 & 980 & 1089\\
\midrule
$\dim S_2(N)$ & 1 & 2 & 3 & 4 & 7 & 9 & 9 & 6 & 22 & 49 & 145 & 109\\
$\dim S_2^{\mathrm{new}}(N)$ & 1 & 0 & 1 & 1 & 5 & 3 & 4 & 4 & 10 & 34 & 13 & 41\\
\bottomrule
\end{tabular}
\end{center}
For instance $S_2(121)$ is spanned by the four rational newforms of the classes 121a to 121d and the two oldforms $f_{11}(z)$, $f_{11}(11z)$; and $S_2(88)$ by the three newforms of
level $88$, $f_{44}(z)$, $f_{44}(2z)$ and $f_{11}(2^jz)$, $0\le j\le3$ ($9=3+2+4$).
\end{example}

\section{$U_p$ on oldforms: stabilisation, depletion, non-semisimplicity}\label{sec:hecke:Up}

This section is the linear algebra behind the proof of Lemma lem:old and Step 1 of the proof of Theorem thm:quant, and behind the depleted forms $f^\Sigma$ and the levels
$M^\Sigma$ of the paper's Section 4.1. The operator $U_p$ acts on the span of the $h(p^jz)$ as a shift with one modified column; its eigenvalues are the two roots of the
Hecke polynomial, or $a_p(h)$, and $0$; and it is not always semisimple.

\begin{proposition}[$U_p$ on the span of $h(p^jz)$]\label{prop:hecke:Up-oldspan}
Let $h$ be a newform of weight $k$ and level $M$, $p$ a prime, $e\ge0$, and $V_{h,p,e}=\operatorname{span}\{h(p^jz):0\le j\le e\}\subseteq S_k(Mp^e)$, of dimension $e+1$
(\Cref{thm:hecke:atkin-lehner}(iii)). Let $U_p=\mathcal U_p$ be the operator of a level divisible by $p$ and by $Mp^e$ (assume $e\ge1$ or $p\mid M$).
\begin{enumerate}[label=(\roman*)]
\item $U_ph(p^jz)=h(p^{j-1}z)$ for $j\ge1$; $U_ph=a_p(h)h-p^{k-1}h(pz)$ if $p\nmid M$; $U_ph=a_p(h)h$ if $p\mid M$.
\item $h(p^ez)$ is a cyclic vector for $U_p$ on $V_{h,p,e}$, so its characteristic and minimal polynomials coincide; they are
\[
X^{e-1}(X^2-a_p(h)X+p^{k-1})\quad(p\nmid M,\ e\ge1),\qquad X^{e}(X-a_p(h))\quad(p\mid M).
\]
\item Every eigenvalue has a one-dimensional eigenspace, spanned as follows. For $p\nmid M$, writing $X^2-a_pX+p^{k-1}=(X-\alpha_p)(X-\beta_p)$: the \emph{$p$-stabilisations} $h_\alpha=h-\beta_ph(pz)$
(eigenvalue $\alpha_p$) and $h_\beta=h-\alpha_ph(pz)$ (eigenvalue $\beta_p$), and, if $e\ge2$, the \emph{$p$-depleted form} $h^{\{p\}}=h-a_ph(pz)+p^{k-1}h(p^2z)$ (eigenvalue $0$). For $p\,\|\,M$: $h$
(eigenvalue $a_p$, nonzero by \Cref{prop:hecke:newform-local}(iii)) and, if $e\ge1$, $h^{\{p\}}=h-a_ph(pz)$ (eigenvalue $0$). For $p^2\mid M$: $h=h^{\{p\}}$ (eigenvalue $a_p=0$).
\item $U_p$ is semisimple on $V_{h,p,e}$ iff its characteristic polynomial is squarefree: iff $e\le2$ and $\alpha_p\ne\beta_p$ ($p\nmid M$); $e\le1$ ($p\,\|\,M$); $e=0$ ($p^2\mid M$). Otherwise the
eigenvalue $0$ has a single Jordan block, of size $e-1$, $e$ or $e+1$ respectively (or, if $p\nmid M$, $e\le2$ and $\alpha_p=\beta_p$, the eigenvalue $\alpha_p$ has a block of size $2$).
\item Each eigenvector in (iii) is a normalised eigenform of $\T_{Mp^e}$, with $U_p$-eigenvalue as stated and the eigenvalues of $h$ for the other $T_n$.
\end{enumerate}
\end{proposition}

\begin{proof}
(i) $\mathcal U_pV_p=1$ (\Cref{lem:hecke:UV}) gives the first formula; $\mathcal U_ph=T_p^{(M)}h-p^{k-1}V_ph$ for $p\nmid M$ and $\mathcal U_ph=U_p^{(M)}h$ for $p\mid M$ (\Cref{prop:hecke:Tp-q-expansion}).
(ii) By (i), $U_p^jh(p^ez)=h(p^{e-j}z)$ for $0\le j\le e$, so the iterates of $h(p^ez)$ span $V_{h,p,e}$: it is cyclic, and the minimal polynomial has degree $e+1$, hence equals the
characteristic polynomial (\Cref{ch:modules}). In the basis $v_j=h(p^jz)$ the matrix is the shift $v_j\mapsto v_{j-1}$ with first column $(a_p,-p^{k-1},0,\dots)$ or $(a_p,0,\dots)$; expanding
the determinant of $X-U_p$ along the first column gives the polynomials. (iii) For a cyclic operator every eigenvalue has a single Jordan block, so a one-dimensional
eigenspace. The vectors are eigenvectors by (i): $U_p(h-\beta h(pz))=(a_p-\beta)h-p^{k-1}h(pz)=\alpha(h-\beta h(pz))$ as $\alpha+\beta=a_p$ and $\alpha\beta=p^{k-1}$; $U_p(h-a_ph(pz)+p^{k-1}h(p^2z))=a_ph-p^{k-1}h(pz)-a_ph+p^{k-1}h(pz)=0$;
$U_p(h-a_ph(pz))=a_ph-a_ph=0$ when $p\mid M$. (iv) A cyclic operator is semisimple iff its minimal polynomial, here equal to the characteristic polynomial, is squarefree. (v) The
$T_n$ with $(n,Mp)=1$ and the $U_{p'}$ with $p'\mid M$, $p'\ne p$, commute with $V_p$ (\Cref{lem:hecke:UV}, \Cref{prop:hecke:old-new}(i)), so they act on $V_{h,p,e}$ by the eigenvalues of $h$; each
listed vector has $a_1=1$ since $h(p^jz)$, $j\ge1$, has $a_1=0$; and $T_{p^r}=U_p^r$ at a level divisible by $p$.
\end{proof}

\begin{example}[levels $22$, $44$, $88$, $121$]\label{ex:hecke:Up-matrices}
For $h=f_{11}$ ($a_2=-2$, $a_{11}=1$) the matrices of $U_p$ on $V_{h,p,e}$, columns the images of $h(p^jz)$, are: at level $22$, $\sm{-2}1{-2}0$, characteristic polynomial $X^2+2X+2$,
$\alpha,\beta=-1\pm i$; at level $44$, characteristic polynomial $X(X^2+2X+2)$, with $f_{11}^{\{2\}}=f_{11}+2f_{11}(2z)+2f_{11}(4z)$ spanning the kernel; at level $88$,
\[
\begin{pmatrix}-2&1&0&0\\-2&0&1&0\\0&0&0&1\\0&0&0&0\end{pmatrix},\qquad X^2(X^2+2X+2),
\]
with a Jordan block of size $2$ at $0$; at level $121$ ($p=11\,\|\,11$), $\sm1100$, characteristic polynomial $X(X-1)$, eigenvectors $f_{11}$ and $f_{11}-f_{11}(11z)$. On the whole space, PARI's
\texttt{mfheckemat} gives for $U_2$ on $S_2(22)$, $S_2(44)$, $S_2(88)$ the characteristic polynomials $X^2+2X+2$, $X^2(X^2+2X+2)$, $X^7(X^2+2X+2)$ and the minimal polynomials $X^2+2X+2$,
$X(X^2+2X+2)$, $X^2(X^2+2X+2)$.
\end{example}

\begin{paperlink}\label{pl:hecke:lem-old}
The proof of Lemma lem:old asserts: ``The eigenvalues of $U_p$ there are therefore the roots $\alpha_p,\beta_p$ of $X^2-a_p(h)X+p$, together with $0$ if $e_p\ge2$, when $p\nmid M$; and $a_p(h)$,
together with $0$ if $e_p\ge1$, when $p\mid M$''; ``The operator $U_p$ is not semisimple on this span when $0$ is a repeated eigenvalue, but each listed eigenvalue has an eigenvector'';
and lists the stabilisations and depleted forms. All of this is \Cref{prop:hecke:Up-oldspan}(ii) to (iv) with $k=2$. (When $p^2\mid M$ the paper's ``$a_p(h)$, together with $0$'' is the single
eigenvalue $0$, since $a_p(h)=0$.) The eigenvalues $\alpha_p,\beta_p$ are those of the arithmetic Frobenius on $V_\ell$ of an elliptic curve with newform $h$, by the
convention of Appendix B (\Cref{ch:notation}).
\end{paperlink}

\begin{corollary}[eigenvalues of $U_p$ in weight $2$]\label{cor:hecke:Up-eigenvalues}
Let $p\mid N$. Every eigenvalue of $U_p$ on $S_2(N)$ is $0$, or $\pm1$, or a root of $X^2-a_p(h)X+p$ for a newform $h$ of level $M\mid N$ with $p\nmid M$. For every prime $\lambda$ of $\Qbar$
not above $p$, every nonzero eigenvalue of $U_p$ is a $\lambda$-adic unit. Hence a generalised eigenform of $\T$ whose $U_p$-eigenvalue $\mu$ satisfies $\mu\equiv0\pmod\lambda$ has $\mu=0$.
\end{corollary}

\begin{proof}
By \Cref{thm:hecke:atkin-lehner}(iii), $S_2(N)=\bigoplus_hV_h$, and on $V_h=\operatorname{span}\{h(dz)\}$ the operator $U_p$ acts through the factor $V_{h,p,e_p}$, $e_p=v_p(N/M_h)$ (write $d=p^jd'$ and use
$\mathcal U_pV_{d'}=V_{d'}\mathcal U_p$). Its eigenvalues there are listed in \Cref{prop:hecke:Up-oldspan}: $0$, $a_p(h)=\pm1$ for $p\,\|\,M_h$ (\Cref{prop:hecke:newform-local}(iii)), or the roots $\alpha,\beta$. These
roots are algebraic integers with $\alpha\beta=p$, so at $\lambda\nmid p$ both are units; and $\pm1$ are units. A nonzero eigenvalue is therefore $\not\equiv0\pmod\lambda$.
\end{proof}

The characteristic polynomials of $U_p$ on the whole space at the levels of the paper (PARI's \texttt{mfheckemat}):
\begin{itemize}
\item level $88$, $U_2$: $X^7(X^2+2X+2)$; minimal polynomial $X^2(X^2+2X+2)$;
\item level $99$, $U_3$: $X^6(X+1)(X^2+X+3)$; minimal polynomial $X(X+1)(X^2+X+3)$;
\item level $121$, $U_{11}$: $X^5(X-1)$; minimal polynomial $X(X-1)$;
\item level $539$, $U_7$: $X^{42}(X-1)^3(X+1)^2(X^2+2X+7)$; minimal polynomial $X(X-1)(X+1)(X^2+2X+7)$;
\item level $980$, $U_7$: $X^{126}(X-1)^8(X+1)^9(X^2-2X+7)$; minimal polynomial $X(X-1)(X+1)(X^2-2X+7)$;
\item level $1089$, $U_3$: $X^{76}(X-1)^{10}(X+1)^{11}(X^2-2X+3)^2(X^2+X+3)^4$; minimal polynomial\\ $X(X-1)(X+1)(X^2-2X+3)(X^2+X+3)$.
\end{itemize}

\begin{paperlink}\label{pl:hecke:units}
In Step 1 of the proof of Theorem thm:quant: ``The generalised eigenforms of $\T$ whose eigencharacter is congruent to $\pi$ modulo $\mf m$ have $U_q\equiv0$ at $q\in S$'', and then only
systems with $U_q=0$ survive, because ``the others have $U_q$-eigenvalue $a_q(g)=\pm1$ (multiplicative case) or a root $\alpha,\beta$ of $X^2-a_q(g)X+q$ (good case), which are $\ell$-adic units since
$\ell\ne q$''. The passage from ``$\equiv0$'' to ``$=0$'' is \Cref{cor:hecke:Up-eigenvalues}.
\end{paperlink}

\begin{remark}[non-semisimplicity and multiplicity]\label{rem:hecke:non-semisimple}
(i) At level $88$ the minimal polynomial of $U_2$ is $X^2(X^2+2X+2)$: $U_2$ is not semisimple, and $U_2(U_2^2+2U_2+2)$ is a nonzero nilpotent element of $\T_{88}\otimes\Q$. So the full Hecke
algebra need not be reduced, and the step ``$\wp\cap I=0$'' in the proof of the paper's Lemma lem:ddt417 must be proved, not deduced from reducedness (as the revised paper does,
\Cref{ch:congruence}). The Jordan block comes from $V_{f_{11},2,3}$ ($e=3$, \Cref{prop:hecke:Up-oldspan}(iv)); on $V_{f_{44},2,1}$ ($2^2\mid44$) $U_2$ is nilpotent of order $2$ as well. (ii) A
large multiplicity of the eigenvalue $0$ does not by itself make $U_p$ non-semisimple: at level $539$ the eigenvalue $0$ of $U_7$ has multiplicity $42$, but the minimal polynomial
is squarefree, because no single $V_{h,7,e}$ has a repeated root ($539=7^2\cdot11$, so $e\le2$ for $7\nmid M_h$, $e\le1$ for $7\,\|\,M_h$, $e=0$ for $7^2\mid M_h$). The paper's ``the eigenvalue $0$ with
multiplicity one'' (Step 1 of the proof of Theorem thm:quant) is the algebraic multiplicity on the span of the $g(dz)$, not on all of $S_2(\Gamma_0(N))$. (iii) (Coleman-Edixhoven.) For a
newform of weight $2$ and level prime to $p$, $\alpha_p\ne\beta_p$ \cite{ColemanEdixhoven}; so $U_p$ is semisimple on the $p$-old part of level $Mp$. For rational $a_p$ this is
elementary: $\alpha_p=\beta_p$ would force $a_p^2=4p$.
\end{remark}

\begin{definition}[depletion and the level $M^\Sigma$]\label{def:hecke:depletion}
Let $f$ be a newform of weight $k$ and level $M$, and $\Sigma$ a finite set of primes. For a prime $p$ put $\delta_p=\delta_p(f)=2$ if $p\nmid M$, $1$ if $p\,\|\,M$, and $0$ if $p^2\mid M$. The
\emph{$\Sigma$-depleted form} is $f^\Sigma=\sum_{(n,\Sigma)=1}a_n(f)q^n$, and $M^\Sigma=M\prod_{p\in\Sigma}p^{\delta_p}$.
\end{definition}

\begin{proposition}[depleted forms]\label{prop:hecke:depletion}
In the notation of \Cref{def:hecke:depletion}:
\begin{enumerate}[label=(\roman*)]
\item $\delta_p$ is the number of roots of the Hecke polynomial at $p$, that is, the degree in $p^{-s}$ of $L_p(f,s)^{-1}$ (\Cref{prop:hecke:euler-product}); for $f=f_E$ in weight $2$,
$\delta_p=2,1,0$ for good, multiplicative and additive reduction, and $\delta_p=\dim V_\ell E^{I_p}$ for $p\ne\ell$ (\Cref{ch:galrep}), the paper's definition in Section 1.1.
\item $f^\Sigma=\bigl(\prod_{p\in\Sigma}D_p\bigr)f$ with $D_p=1-a_pV_p+p^{k-1}V_p^2$ ($p\nmid M$), $D_p=1-a_pV_p$ ($p\,\|\,M$), $D_p=1$ ($p^2\mid M$).
\item $f^\Sigma\in S_k(M^\Sigma)$, and $M^\Sigma$ is the least level $L$ with $f^\Sigma\in S_k(L)$.
\item $f^\Sigma$ is a normalised eigenform of the full Hecke algebra of level $M^\Sigma$, with $U_pf^\Sigma=0$ for $p\in\Sigma$ and the eigenvalues of $f$ at the other $T_n$, and its generalised
eigenspace in $S_k(M^\Sigma)$ is the line $\C f^\Sigma$.
\item At a level $N$ with $M^\Sigma\mid N$, $f^\Sigma$ is still a normalised eigenform of $\T_N$, but its generalised eigenspace can be larger.
\item $L(f^\Sigma,s)=L(f,s)\prod_{p\in\Sigma}L_p(f,s)^{-1}$.
\end{enumerate}
\end{proposition}

\begin{proof}
(i) is \Cref{prop:hecke:euler-product} and \Cref{prop:hecke:newform-local}(iii),(iv); the Galois description is \Cref{ch:galrep}. (ii) The operators $D_p$ for different $p$ commute
(\Cref{lem:hecke:UV}), and $D_p$ only changes the coefficients at multiples of $p$. The $n$-th coefficient of $D_pf$, for $n=p^rm$ with $p\nmid m$, is $a_m$ times
$a_{p^r}-a_pa_{p^{r-1}}+p^{k-1}a_{p^{r-2}}$ (good case), $a_{p^r}-a_pa_{p^{r-1}}$ (case $p\,\|\,M$) or $a_{p^r}$ (case $p^2\mid M$), by multiplicativity, with $a_{p^{-1}}=a_{p^{-2}}=0$; for $r\ge1$ this vanishes by
\Cref{prop:hecke:eigenforms}(ii) and $a_p=0$ when $p^2\mid M$, and for $r=0$ it is $a_m$. (iii) By (ii), $f^\Sigma$ is a combination of the $f(dz)$ with $d\mid\prod_{p\in\Sigma}p^{\delta_p}$, so it lies in
$S_k(M^\Sigma)$. Its coefficient on $f(d_0z)$, $d_0=\prod_{p\in\Sigma}p^{\delta_p}$, is the product of the leading coefficients of the $D_p$, namely $p^{k-1}$ or $-a_p$, nonzero. If $f^\Sigma\in S_k(L)$,
compare its expansions in the bases of \Cref{thm:hecke:atkin-lehner}(iii) at the levels $L$ and $\lcm(L,M^\Sigma)$: the basis at level $L$ is part of the basis at level $\lcm(L,M^\Sigma)$, and
coordinates are unique, so $f(d_0z)$ occurs at level $L$, that is, $Md_0\mid L$. (iv) The coefficients of $f^\Sigma$ are multiplicative, vanish at $p^r$ ($r\ge1$, $p\in\Sigma$), so that $b_{p^r}=0^r$,
and agree with those of $f$ elsewhere; by \Cref{lem:hecke:eigen-criterion} at the level $M^\Sigma$, $f^\Sigma$ is a normalised eigenform with the stated eigenvalues. A generalised
eigenvector for its eigencharacter is in particular an eigenvector of $\T^{(M^\Sigma)}$ for the system of $f$ (\Cref{thm:hecke:self-adjoint}), so it lies in $V_f=\operatorname{span}\{f(dz):d\mid d_0\}$
(\Cref{thm:hecke:atkin-lehner}(ii)). Writing $d=\prod_pp^{j_p}$, the map $\bigotimes_{p\in\Sigma}V_{f,p,\delta_p}\to V_f$, $\otimes_p h(p^{j_p}z)\mapsto f(dz)$ (with $h=f$ in each factor) is an isomorphism under which $U_p$
acts on the $p$-th factor (because $\mathcal U_pV_{d'}=V_{d'}\mathcal U_p$ for $(d',p)=1$). On $V_{f,p,\delta_p}$ the characteristic polynomial of $U_p$ is $X(X^2-a_pX+p^{k-1})$, $X(X-a_p)$ with $a_p\ne0$, or $X$,
so $0$ is a simple root; the joint generalised eigenspace for the eigenvalues $0$ is the tensor product of these lines, a line. (v) The eigenform property is (iv) with
\Cref{lem:hecke:eigen-criterion} at level $N$. For the second statement see \Cref{ex:hecke:level-88}: $f_{11}^{\{2\}}$ spans its generalised eigenspace at level $44=M^{\{2\}}$, but lies in a
Jordan block of $U_2$ at level $88$. (vi) is \Cref{prop:hecke:euler-product} applied to $f^\Sigma$ and $f$.
\end{proof}

\begin{remark}[why $\delta_p$]\label{rem:hecke:why-delta}
Depletion at $p$ removes the Euler factor at $p$ (\Cref{prop:hecke:depletion}(vi)), and each root of the Hecke polynomial removed costs one factor of $p$ in the level: $D_p$ is a polynomial
of degree $\delta_p$ in $V_p$. In the language of \Cref{prop:hecke:Up-oldspan}, the eigenvalue $0$ of $U_p$ appears in $V_{f,p,e}$ exactly when $e\ge\delta_p$, and with multiplicity one when
$e=\delta_p$. For an additive prime nothing is removed and nothing is added: this is the paper's ``adding an additive prime to $\Sigma$ changes neither the level nor the ideal''
(paragraph after Theorem thm:dfg14).
\end{remark}

\begin{example}[depleted forms]\label{ex:hecke:depleted}
$f_{11}^{\{2\}}=f_{11}+2f_{11}(2z)+2f_{11}(4z)$, of level $44$ ($a_2=-2$; checked to $n=4000$); $f_{11}^{\{11\}}=f_{11}-f_{11}(11z)$, of level $121$ ($a_{11}=1$); for $f$ of class 121a, $\Sigma=\{11\}$:
$f^\Sigma=f$, level $121$ ($\delta_{11}=0$); for $g$ of class 77c and $\Sigma=\{7\}$: $g^\Sigma=g-a_7(g)g(7z)=g+g(7z)$, level $539$; for $g=f_{11}$ and $\Sigma=\{7\}$: $g^\Sigma=g+2g(7z)+7g(49z)$, level $539$
($a_7(f_{11})=-2$).
\end{example}

\begin{paperlink}\label{pl:hecke:section41}
In Section 4.1 of the paper, ``Put $M^\Sigma=M\prod_{p\in\Sigma}p^{\delta_p}$ and let $f^\Sigma=\sum_{(n,\Sigma)=1}a_n(f)q^n$ be the $\Sigma$-depleted eigenform, of level $M^\Sigma$'' is
\Cref{def:hecke:depletion} and \Cref{prop:hecke:depletion}(iii), (iv). The proof of Theorem thm:quant: ``$N_S^\Sigma=N_S\prod_{q\in S}q^{\delta_q}=N$: the $S$-depleted form $g^S$ lives at level
$N$'', and in Step 1 ``on that span the operators $U_q$, $q\in S$, have the eigenvalue $0$ with multiplicity one'', is \Cref{prop:hecke:depletion}(iv) for $g$ at level $N=N_S^S$.
\end{paperlink}

\section{The structure of $\T$, generalised eigenspaces and congruences}\label{sec:hecke:structure}

This section assembles the previous ones into a description of $\T\otimes\Qbar$ and of the generalised eigenspaces of $\T$. It proves the decomposition $\T\otimes\Q=\Q\times\T'$ that
the proof of Lemma lem:crit writes down without proof; it separates the $\T$-stable complement of an eigenform from its Petersson complement, the distinction on which
Step 1 of the proof of Theorem thm:quant depends; it describes the local components $\T_{\mf m}$ of $\T_\ell$ and their effect on forms; and it explains how congruences at
primes become congruences at all $n$.

\begin{theorem}[the structure of $\T$]\label{thm:hecke:structure}
Let $k$ be even and $N\ge1$. For a newform $h$ of level $M_h\mid N$ put $e_p=e_p(h)=v_p(N/M_h)$ and $V_h=\operatorname{span}\{h(dz):d\mid N/M_h\}$.
\begin{enumerate}[label=(\roman*)]
\item $S_k(N)=\bigoplus_hV_h$, the sum over the newforms $h$ of levels dividing $N$; $V_h$ is the eigenspace of $\T^{(N)}$ for the system of $h$, and as a module over $\C[U_p:p\mid N]$,
$V_h\simeq\bigotimes_{p\mid N}V_{h,p,e_p}$, with $U_p$ acting on the $p$-th factor.
\item Let $m_{h,p}$ be the characteristic polynomial of $U_p$ on $V_{h,p,e_p}$ (\Cref{prop:hecke:Up-oldspan}(ii); it is $X-a_p(h)$ when $e_p=0$). Then
\[
\T\otimes\C\simeq\prod_hA_h,\qquad A_h\simeq\bigotimes_{p\mid N}\C[X]/(m_{h,p}(X)),
\]
where $A_h$ is the image of $\T\otimes\C$ in $\End(V_h)$; $\dim A_h=\dim V_h$, $\rank_\Z\T=\dim S_k(N)$, and $S_k(N)$ is free of rank one over $\T\otimes\C$.
\item The generalised eigenspaces of $\T$ in $S_k(N)$ are the joint generalised eigenspaces $V_h(\mu)$, $\mu=(\mu_p)_{p\mid N}$ a choice of a root $\mu_p$ of each $m_{h,p}$; $V_h(\mu)$ has dimension the
product of the multiplicities of the $\mu_p$. $\T\otimes\Q$ is reduced iff every $m_{h,p}$ is squarefree.
\end{enumerate}
\end{theorem}

\begin{proof}
(i) is \Cref{thm:hecke:atkin-lehner}(iii),(iv); the tensor structure is the argument in the proof of \Cref{prop:hecke:depletion}(iv): $f(dz)\leftrightarrow\otimes_ph(p^{v_p(d)}z)$, and $\mathcal U_p$ commutes
with $V_{d'}$ for $(d',p)=1$. (ii) On $V_h$ every $T_n$ with $(n,N)=1$ is a scalar, and every $T_n$ is a product of such an operator and powers of the $U_p$ (\Cref{thm:hecke:Tn}); so $A_h$ is
the algebra generated by the $U_p$ on $\bigotimes_pV_{h,p,e_p}$. Each $U_p$ has a cyclic vector on its factor, so the algebra it generates there is $\C[X]/(m_{h,p})$ (\Cref{ch:modules}); and
the map $\bigotimes_p\C[X]/(m_{h,p})\to\bigotimes_p\End(V_{h,p,e_p})=\End(V_h)$ is injective, a tensor product of injective linear maps over a field. Hence $A_h\simeq\bigotimes_p\C[X]/(m_{h,p})$, of dimension
$\prod_p\deg m_{h,p}=\prod_p(e_p+1)=\dim V_h$. The map $\T\otimes\C\to\prod_hA_h$ is injective because $\T\otimes\C$ acts faithfully on $S_k(N)=\bigoplus V_h$ ($\T$ preserves the lattice $S_k(N;\Z)$ of forms with integer
coefficients by \Cref{thm:hecke:Tn}(iii), which spans $S_k(N)$ in every even weight by the integral basis theorem of \Cref{ch:modforms}; so $\T\subseteq\End_\Z S_k(N;\Z)$, and tensoring this
inclusion with the flat $\Z$-module $\C$ gives $\T\otimes\C\subseteq\End_\C S_k(N)$). It is surjective: the systems of the $V_h$ are pairwise distinct (\Cref{thm:hecke:atkin-lehner}(ii)), so by Lagrange
interpolation in finitely many $T_n$, $(n,N)=1$, which act semisimply (\Cref{thm:hecke:self-adjoint}), there are idempotents $e_h\in\T^{(N)}\otimes\C$ acting as $1$ on $V_h$ and $0$ on the other $V_{h'}$;
then $e_h(\T\otimes\C)$ maps onto $A_h$ and to $0$ in the other factors. The dimension count gives $\rank\T=\sum_h\dim V_h=\dim S_k(N)$ (for $k=2$ this is also
\Cref{thm:hecke:perfect-pairing}). Each $V_h$ is cyclic over $A_h$, generated by $h((N/M_h)z)$ (a tensor product of cyclic vectors), so $S_k(N)=\bigoplus V_h$ is cyclic over $\prod A_h$, hence free of
rank one by the dimension count. (iii) The generalised eigenspaces of the commutative algebra $A_h$ acting on $\bigotimes_pV_{h,p,e_p}$ are the tensor products of the generalised eigenspaces of
the $U_p$; the algebra $\bigotimes_p\C[X]/(m_{h,p})$ is reduced iff each factor is, iff each $m_{h,p}$ is squarefree; and $\T\otimes\Q$ is reduced iff $\T\otimes\C$ is (\Cref{ch:modules}: a finite
algebra over a perfect field is reduced iff it stays reduced after extension of scalars).
\end{proof}

\begin{theorem}[the newform splitting]\label{thm:hecke:newform-splitting}
Let $f\in S_2(N)$ be a newform with rational coefficients, $\phi_f\colon\T\to\Z$ its eigencharacter, and $f^\perp$ its Petersson orthogonal complement.
\begin{enumerate}[label=(\roman*)]
\item $f^\perp$ is the sum of the eigenspaces of $\T^{(N)}$ other than the line $\C f$; it is $\T$-stable and spanned by forms with rational coefficients.
\item There is $u\in\T$ with $uf^\perp=0$ and $\phi_f(u)\ne0$; then $e=u/\phi_f(u)\in\T\otimes\Q$ is an idempotent, acting as $1$ on $f$ and $0$ on $f^\perp$.
\item Let $\T'$ be the image of $\T\otimes\Q$ in $\End(f^\perp)$. Then $t\mapsto(\phi_f(t),t|_{f^\perp})$ is an isomorphism of $\Q$-algebras $\T\otimes\Q\simeq\Q\times\T'$, and $\T'$ acts faithfully on $f^\perp$.
\item With $I=\ker\phi_f$ and $J=\Ann_\T(f^\perp)$: $I=\T\cap(0\times\T')$, $J=\T\cap(\Q\times0)$, $I\cap J=0$, $\phi_f(J)\ne0$, and $\T/(I+J)\simeq\Z/\phi_f(J)$ is finite cyclic.
\end{enumerate}
\end{theorem}

\begin{proof}
(i) By \Cref{thm:hecke:self-adjoint}, $S_2(N)$ is the orthogonal sum of the eigenspaces of $\T^{(N)}$, and the eigenspace of the system of $f$ is $\C f$ (\Cref{thm:hecke:atkin-lehner}(ii) with $L=M=N$).
So $f^\perp$ is the sum of the others. Each eigenspace of $\T^{(N)}$ is $\T$-stable, as $\T$ is commutative. Rationality follows from (ii): $f^\perp=\ker e$ with $e\in\T\otimes\Q$, and $\T\otimes\Q$
preserves $S_2(N;\Q)$.

(ii) Let $\psi_1,\dots,\psi_s$ be the systems of $\T^{(N)}$ occurring in $f^\perp$, all different from the system of $f$ on $\T^{(N)}$. For each $i$ choose a prime $p_i\nmid N$ with $\psi_i(T_{p_i})\ne a_{p_i}(f)$: two
systems that agree on all $T_p$, $p\nmid N$, agree on $\T^{(N)}$ by \Cref{prop:hecke:hecke-algebra}(ii) and \Cref{prop:hecke:eigenforms}(ii). Let $P_i\in\Z[X]$ be a nonzero integer multiple of the minimal
polynomial of the algebraic integer $\psi_i(T_{p_i})$ over $\Q$, and $u=\prod_iP_i(T_{p_i})\in\T$. On the eigenspace of $\psi_i$, $T_{p_i}$ is the scalar $\psi_i(T_{p_i})$, so $P_i(T_{p_i})$ kills it; the factors
commute, so $u$ kills $f^\perp$. And $\phi_f(u)=\prod_iP_i(a_{p_i}(f))\ne0$: $P_i$ is irreducible up to a constant, and a root $a_{p_i}(f)\in\Q$ would force $P_i$ to be linear with root $\psi_i(T_{p_i})=a_{p_i}(f)$.
Then $e=u/\phi_f(u)$ acts as $1$ on $f$ and $0$ on $f^\perp$; as $S_2(N)=\C f\oplus f^\perp$, $e^2=e$ on $S_2(N)$, hence in $\T\otimes\Q$ (faithfulness).

(iii) $\T\otimes\Q=e(\T\otimes\Q)\times(1-e)(\T\otimes\Q)$. For $t\in\T\otimes\Q$, $te-\phi_f(t)e$ kills $f$ and $f^\perp$, so $te=\phi_f(t)e$ and $e(\T\otimes\Q)=\Q e\simeq\Q$ via $\phi_f$. The element $1-e$ acts as $0$ on
$f$ and $1$ on $f^\perp$, so $(1-e)(\T\otimes\Q)$ maps isomorphically onto $\T'$ (injective because an element of it acting by $0$ on $f^\perp$ also kills $f$, hence is $0$). The resulting
isomorphism is $t\mapsto(\phi_f(t),t|_{f^\perp})$.

(iv) $\T$ is torsion-free, so $\T\subseteq\T\otimes\Q=\Q\times\T'$. An element $t$ lies in $I$ iff its first coordinate vanishes, and in $J$ iff its second does; so $I\cap J=0$. As $u\in J$, $\phi_f(J)\ne0$.
The homomorphism $\phi_f$ is onto $\Z$ ($\phi_f(T_1)=1$), with kernel $I$, so $\T/(I+J)\simeq\Z/\phi_f(J)$, finite cyclic.
\end{proof}

\begin{paperlink}\label{pl:hecke:splitting}
The proof of Lemma lem:crit begins: ``Writing $\T\otimes\Q=\Q\times\T'$ with $\T'$ acting faithfully on $f^\perp$ (the $f$-line is $\T$-stable, and so is $f^\perp$, because $f$ is an eigenvector of the
adjoints of the $T_n$ as well), one has $I=\T\cap(0\times\T')$ and $J=\T\cap(\Q\times0)$, so $I\cap J=0$ and $\T/(I+J)\simeq\Z/\phi_f(J)$ is cyclic''. Every clause is \Cref{thm:hecke:newform-splitting}. The
input that the paper leaves implicit is multiplicity one for the newform $f$ (\Cref{thm:hecke:atkin-lehner}(ii)): without it the local factor of $\T\otimes\Q$ through which $\phi_f$ factors could be
larger than $\Q$, as happens for the depleted form $f_{11}^{\{2\}}$ at level $88$ (\Cref{ex:hecke:level-88}).
\end{paperlink}

\begin{proposition}[$\T$-stable versus Petersson complements]\label{prop:hecke:complements}
Let $g\in S_k(N)$ be a normalised eigenform of $\T=\T_N$ with eigencharacter $\pi$ whose generalised eigenspace is the line $\C g$ (a newform of level $N$; or $f^\Sigma$ at level $M^\Sigma$,
\Cref{prop:hecke:depletion}(iv)). Let $g^c$ be the $\T$-stable complement of $\C g$, the sum of the other generalised eigenspaces.
\begin{enumerate}[label=(\roman*)]
\item $S_k(N)=\C g\oplus g^c$, and $\T\otimes\C=\C\times\T''$ with $\T''$ acting faithfully on $g^c$. If $\pi$ takes values in $\Q$ (so $k=2$ and $g$ has rational coefficients), the idempotent of the
first factor lies in $\T\otimes\Q$, $g^c$ is spanned by forms with rational coefficients, $g^c\cap S_2(N;\Z)$ is saturated of rank $d-1$, and the conclusions of
\Cref{thm:hecke:newform-splitting}(iii),(iv) hold with $g^c$ in place of $f^\perp$.
\item If $g$ is a newform of level $N$, $g^c=g^\perp$.
\item If $g$ is not a newform, $g^c\ne g^\perp$ in general: for $g=f_{11}^{\{2\}}$ at level $44$, the stabilisation $h=f_{11}-\beta f_{11}(2z)$ lies in $g^c$ and $\langle g,h\rangle=\frac5{12}\langle f_{11},f_{11}\rangle_{44}\ne0$.
\end{enumerate}
\end{proposition}

\begin{proof}
(i) The decomposition of a finite commutative algebra and its module into generalised eigenspaces (\Cref{ch:modules}) gives an idempotent $e_g\in\T\otimes\C$ projecting onto the
generalised eigenspace $\C g$ along $g^c$, and $e_g(\T\otimes\C)=\C e_g$ because that eigenspace is a line on which $\T$ acts by $\pi$. If $\pi$ is $\Q$-valued, the idempotent of $\T\otimes\Q$
attached to the maximal ideal $\ker(\pi\otimes\Q)$ acts on $S_k(N)\otimes\C$ as the projection onto the sum of the generalised eigenspaces of the conjugates of $\pi$, which is $\C g$ ($\pi$
is its own only conjugate); so $e_g\in\T\otimes\Q$, and $g^c=\ker e_g$ is defined over $\Q$. The rest is as in \Cref{prop:hecke:rf-finite} and \Cref{thm:hecke:newform-splitting}. (ii) is
\Cref{thm:hecke:newform-splitting}(i): both are the sum of the other eigenspaces. (iii) $h$ is an eigenvector of $U_2$ with eigenvalue $\alpha\ne0$, while $g$ has $U_2$-eigenvalue $0$, so $h$
lies in another generalised eigenspace. The Gram matrix of $f_{11}(z)$, $f_{11}(2z)$, $f_{11}(4z)$ at level $44$ is $\langle f_{11},f_{11}\rangle_{44}$ times
\[
\begin{pmatrix}1&-1/3&1/24\\-1/3&1/4&-1/12\\1/24&-1/12&1/16\end{pmatrix}
\]
(PARI's \texttt{mfpetersson}, $40$ digits; the $2\times2$ block is \Cref{prop:hecke:Up-adjoint}(ii)), and with $g=(1,2,2)$ and $h=(1,-\beta,0)$ in these coordinates, $\langle g,h\rangle=\frac5{12}\langle f_{11},f_{11}\rangle_{44}$
for both $\beta=-1\pm i$.
\end{proof}

\begin{paperlink}\label{pl:hecke:complements}
The proof of Theorem thm:twist speaks of ``$g^\perp$, the $\T$-stable complement of the $g$-line (the sum of the generalised eigenspaces of the other systems of eigenvalues)'': at level $N_S$,
where $g$ is a newform, the two notions agree (\Cref{prop:hecke:complements}(ii)). Step 1 of the proof of Theorem thm:quant uses ``the $\T$-stable complement of the $g^S$-line (not the
Petersson complement, which is not $\T$-stable when the $U_p$ are involved)''; \Cref{prop:hecke:complements}(iii) shows that the two differ there, and (i) shows that the argument of
Lemma lem:crit runs with the $\T$-stable complement for any eigencharacter whose generalised eigenspace is a line (the proof of the lemma in \Cref{ch:modgalrep} uses
only the conclusions of \Cref{thm:hecke:newform-splitting}).
\end{paperlink}

\begin{proposition}[local components]\label{prop:hecke:local-components}
Let $\ell$ be a prime, $k=2$, $S_\ell=S_2(N;\Zl)$, and $\T_\ell=\T\otimes\Zl=\prod_{\mf m}\T_{\mf m}$ the decomposition over the maximal ideals $\mf m$ of $\T_\ell$ (those of $\T$ containing $\ell$), with
idempotents $e_{\mf m}$ (\Cref{ch:modules}).
\begin{enumerate}[label=(\roman*)]
\item $S_\ell=\bigoplus_{\mf m}e_{\mf m}S_\ell$, and $e_{\mf m}S_\ell\otimes\Qbar_\ell$ is the sum of the generalised eigenspaces of the $\phi\colon\T\to\Qbar_\ell$ whose reduction has kernel $\mf m$.
\item A ring homomorphism $\phi\colon\T_\ell\to\mathcal O_\lambda$ factors through exactly one $\T_{\mf m}$; the other factors lie in $\ker\phi$, and $\Ann_{\T_\ell}(\ker\phi)=\Ann_{\T_{\mf m}}(\ker\phi\cap\T_{\mf m})$, so
$\eta(\phi)=\eta(\phi|_{\T_{\mf m}})$.
\item $\langle e_{\mf m}g,t\rangle=\langle g,e_{\mf m}t\rangle$; the pairing of \Cref{thm:hecke:perfect-pairing} restricts to perfect pairings $e_{\mf m}S_\ell\times\T_{\mf m}\to\Zl$, and
$e_{\mf m}S_\ell\simeq\Hom_{\Zl}(\T_{\mf m},\Zl)$.
\item $e_{\mf m}$ preserves $S_\ell$ and every $\T$-stable subspace of $S_2(N)\otimes\Ql$. If $g\in e_{\mf m}S_\ell$, $h\in S_\ell\cap W$ with $W$ a $\T$-stable subspace, and $g\equiv h\pmod{\ell^t}$, then
$g\equiv e_{\mf m}h\pmod{\ell^t}$ with $e_{\mf m}h\in e_{\mf m}S_\ell\cap W$.
\end{enumerate}
\end{proposition}

\begin{proof}
(i) The first statement holds because the $e_{\mf m}$ are orthogonal idempotents with sum $1$ acting on $S_\ell$. For the second, $S_\ell\otimes\Qbar_\ell$ is the sum of the generalised
eigenspaces $V(\phi)$, $\phi\colon\T\to\Qbar_\ell$ (\Cref{ch:modules}); such $\phi$ take values in the integral closure of $\Zl$, since $\T$ is finite over $\Z$. On $V(\phi)$ the idempotent $e_{\mf m}$
has the single eigenvalue $\phi(e_{\mf m})\in\{0,1\}$, and an idempotent with one eigenvalue acts by it (its minimal polynomial divides $X^2-X$). Finally $\phi(e_{\mf m})=1$ iff $e_{\mf m}$ is not
in the maximal ideal $\phi^{-1}(\mf m_{\overline\Z_\ell})$ of $\T_\ell$, iff that ideal is $\mf m$. (ii) $\phi(e_{\mf m})\in\{0,1\}$ and $\sum_{\mf m}\phi(e_{\mf m})=1$, so exactly one $\phi(e_{\mf m})$ is $1$, and $e_{\mf m'}\in\ker\phi$
for $\mf m'\ne\mf m$, so $\T_{\mf m'}=e_{\mf m'}\T_\ell\subseteq\ker\phi$. If $a$ kills $\ker\phi$, it kills each $e_{\mf m'}$, so $a=e_{\mf m}a\in\T_{\mf m}$; and an element of $\T_{\mf m}$ kills $\ker\phi$ iff it kills
$\ker\phi\cap\T_{\mf m}$, since it kills the other factors. (iii) $\langle e_{\mf m}g,t\rangle=a_1(te_{\mf m}g)=\langle g,e_{\mf m}t\rangle$. Hence $\langle e_{\mf m}S_\ell,\T_{\mf m'}\rangle=\langle S_\ell,e_{\mf m}e_{\mf m'}\T_\ell\rangle=0$ for $\mf m\ne\mf m'$: the
perfect pairing $S_\ell\times\T_\ell\to\Zl$ (\Cref{thm:hecke:perfect-pairing}(iii)) is block diagonal for the two decompositions, so each block is perfect. (iv) $e_{\mf m}\in\T_\ell$ acts on $S_\ell$, and a
$\T$-stable $\Ql$-subspace is $\T_\ell$-stable. If $g-h=\ell^tx$ with $x\in S_\ell$, then $g-e_{\mf m}h=e_{\mf m}(g-h)=\ell^te_{\mf m}x$.
\end{proof}

\begin{paperlink}\label{pl:hecke:local}
After Lemma lem:hecke the revised paper says: ``the congruence ideal is computed in $\T_{\mf m}$: the other local factors of $\T\otimes\Z_\ell$ lie in the kernel $\wp$, $\phi_f$ factors through
$\T_{\mf m}$, and $\operatorname{Ann}(\wp)$ decomposes along the idempotents'': \Cref{prop:hecke:local-components}(ii) (the other factors lie in the kernel, not in the annihilator). Step 1 of the proof of
Theorem thm:quant: ``projecting $h$ to the $\mf m$-component (an idempotent of $\T_\ell$, so the projection stays in the complement and the congruence survives)'':
\Cref{prop:hecke:local-components}(iv). The perfectness of the cohomological pairing $\hat\delta_!$ on each local component is the same argument with $\hat\delta_!$ in place of $\langle g,t\rangle$
(\Cref{ch:moddeg}).
\end{paperlink}

\begin{proposition}[congruences]\label{prop:hecke:congruences}
Let $\mathcal O$ be the ring of integers of a number field and $\lambda$ a prime of $\mathcal O$.
\begin{enumerate}[label=(\roman*)]
\item For normalised eigenforms $f,g$ of $\T_N$ with coefficients in $\mathcal O$ the following are equivalent: (a) $a_n(f)\equiv a_n(g)\pmod\lambda$ for all $n\ge1$; (b) the same for all primes $n=p$,
including $p\mid N$; (c) $\phi_f(t)\equiv\phi_g(t)\pmod\lambda$ for all $t\in\T$. They imply (d) $\mf m_{f,\lambda}=\mf m_{g,\lambda}$, and are implied by it when
$\T/\mf m_{f,\lambda}=\F_\ell$ (for instance when $f$ has rational coefficients).
\item Let $f$, $g$ be newforms of weight $k$ and levels $M_f$, $M_g$ with coefficients in $\mathcal O$, and $\Sigma$ a finite set of primes containing every prime dividing $M_fM_g$. If
$a_p(f)\equiv a_p(g)\pmod\lambda$ for all $p\notin\Sigma$, then $f^\Sigma\equiv g^\Sigma\pmod\lambda$, both forms lying in $S_k(L)$, $L=\lcm(M_f^\Sigma,M_g^\Sigma)$.
\end{enumerate}
\end{proposition}

\begin{proof}
(i) (a)$\Rightarrow$(b) is trivial. (b)$\Rightarrow$(a): by \Cref{prop:hecke:eigenforms}(ii) the $a_n$ of both forms are given by the same integer polynomials in the $a_p$ (multiplicativity,
$a_{p^r}=a_p^r$ for $p\mid N$, the recursion with the same $p^{k-1}$ for $p\nmid N$). (a)$\Leftrightarrow$(c): $\T$ is the $\Z$-span of the $T_n$ (\Cref{prop:hecke:hecke-algebra}(ii)) and $\phi_f(T_n)=a_n(f)$. (c)$\Rightarrow$(d): the two reductions $\T\to\mathcal O/\lambda$ coincide, so they have the same kernel. If the kernels coincide and $\T/\mf m=\F_\ell$, both reductions are the unique
ring homomorphism $\T\to\T/\mf m=\F_\ell\subseteq\mathcal O/\lambda$. (Without that hypothesis (d) only determines the reduction up to an automorphism of the residue field: Galois-conjugate
forms $f$, $f^\sigma$ with $\sigma\lambda=\lambda$ have the same maximal ideal.) (ii) For $(n,\Sigma)=1$ write $n=\prod p^r$ with $p\notin\Sigma$, so $p\nmid M_fM_g$; both $a_{p^r}(f)$ and $a_{p^r}(g)$ are given by the same recursion in $a_p$,
so they are congruent, and by multiplicativity $a_n(f)\equiv a_n(g)$. For $(n,\Sigma)>1$ both coefficients of the depleted forms vanish. The levels are \Cref{prop:hecke:depletion}(iii).
\end{proof}

\begin{remark}[Sturm versus the recursion]\label{rem:hecke:sturm-vs-recursion}
Two tools turn finitely many congruences into congruences at all $n$, and they should not be confused. The recursion (\Cref{prop:hecke:congruences}(i)) transfers congruences at all
primes, including those dividing the level, to all $n$, between eigenforms of the full $\T_N$. Sturm's bound (\Cref{ch:modforms}) transfers a finite check, $n\le B$, to all $n$, between any two
forms of the same space $S_k(N)$. Neither turns a congruence at the primes $p\notin\Sigma$ alone into one at all $n$: that requires depletion (\Cref{prop:hecke:congruences}(ii)) or a choice of
$U_p$-eigenvalues (the proof of Lemma lem:old, \Cref{ch:modgalrep}). The congruences of the paper's Example of Section 3 hold at all $n$ by Sturm's bound (\Cref{ch:modforms}).
\end{remark}

\section{Twisting}\label{sec:hecke:twists}

This section treats the twist of a modular form by a Dirichlet character, in the generality the paper needs: quadratic characters $\chi_{q^*}$ and their products $\chi_S$. It
proves the elementary properties of the twist operator (level, integrality, intertwining with Hecke operators, twisting twice), states the theorem of Atkin and Li (the
paper's Lemma lem:twist (iii)), and proves the transport of congruences along the twist map $\iota$ used in the proofs of Theorems thm:twist and thm:quant. In this section
$q$ denotes an odd prime; the $q$-expansion variable appears only as $e^{2\pi inz}$.

\begin{proposition}[the twist operator]\label{prop:hecke:twist-operator}
Let $\chi$ be a primitive Dirichlet character of conductor $r$, $\mathfrak g(\bar\chi)=\sum_{u\bmod r}\bar\chi(u)e^{2\pi iu/r}$ the Gauss sum of $\bar\chi$, and $f=\sum a_ne^{2\pi inz}\in S_k(M)$.
\begin{enumerate}[label=(\roman*)]
\item $f\otimes\chi=\sum_n\chi(n)a_ne^{2\pi inz}=\mathfrak g(\bar\chi)^{-1}\sum_{u\bmod r}\bar\chi(u)\,f\bigl(z+\tfrac ur\bigr)$.
\item $f\otimes\chi\in S_k(\Gamma_0(L),\chi^2)$ with $L=\lcm(M,r^2)$; for quadratic $\chi$, $f\otimes\chi\in S_k(L)$.
\item If $\chi$ is quadratic and $f\in S_2(M;\Z)$, then $f\otimes\chi\in S_2(L;\Z)$.
\item (Intertwining.) Let $\chi$ be quadratic. For a prime $p\nmid r$, $T_p^{(L)}(f\otimes\chi)=\chi(p)\,(T_p^{(M)}f)\otimes\chi$, where on each side $T_p$ means $U_p$ if $p$ divides the level; for
$p\mid r$, $U_p^{(L)}(f\otimes\chi)=0$. Consequently $f\mapsto f\otimes\chi$ maps the (generalised) eigenspace of $\T^{(rM)}$ for a system $(\lambda_p)$ into the (generalised) eigenspace for
$(\chi(p)\lambda_p)$, and a normalised eigenform of $\T_M$ to a normalised eigenform of $\T_L$.
\item (Twisting twice.) $(f\otimes\chi)\otimes\bar\chi=\sum_{(n,r)=1}a_ne^{2\pi inz}=f^{\{p\mid r\}}$, the depletion at the primes dividing $r$. In particular $(g\otimes\chi_{q^*})\otimes\chi_{q^*}=g^{\{q\}}$, not $g$.
\item (The paper's $\iota$.) For a finite set $S$ of odd primes, $D=\prod_{q\in S}q$ and $\chi_S=\prod_{q\in S}\chi_{q^*}$, the map $\iota=\iota_S\colon S_2(M)\to S_2(\lcm(M,D^2))$, $h\mapsto h\otimes\chi_S$, is
$\C$-linear, maps $S_2(M;\Z)$ into $S_2(\lcm(M,D^2);\Z)$, kills every $h(dz)$ with $(d,D)>1$, satisfies $\iota(h)=\iota(h^S)$ and $\iota\circ\iota(h)=h^S$, and maps the generalised eigenspace of $\T^{(DM)}$ for
$(\lambda_p)$ into that for $(\chi_S(p)\lambda_p)$.
\end{enumerate}
\end{proposition}

\begin{proof}
(i) For a primitive character, $\sum_{u\bmod r}\bar\chi(u)e^{2\pi inu/r}=\chi(n)\mathfrak g(\bar\chi)$ for every integer $n$, including $(n,r)>1$, where both sides vanish, and $\mathfrak g(\bar\chi)\ne0$
($\abs{\mathfrak g}^2=r$) (\Cref{ch:rings}). Insert $f(z+u/r)=\sum_na_ne^{2\pi inu/r}e^{2\pi inz}$.

(ii) Let $\gamma=\sm abcd\in\Gamma_0(L)$ and $u\in\Z$. Choose $u'\equiv d^2u\pmod r$. Then
\[
\gamma'=\Mtwo1{u/r}01\gamma\Mtwo1{-u'/r}01=\Mtwo{a+cu/r}{\ b+(du-au')/r-cuu'/r^2}{c}{d-cu'/r}
\]
is integral: $r^2\mid c$, and $du-au'\equiv du-ad^2u=du(1-ad)\equiv0\pmod r$, because $ad\equiv1\pmod r$ (as $r\mid c$). It has determinant $1$ and lower left entry $c\equiv0\pmod M$, so $\gamma'\in\Gamma_0(M)$ and
$f|_k\sm1{u/r}01\gamma=f|_k\gamma'\sm1{u'/r}01=f|_k\sm1{u'/r}01$. As $u$ runs modulo $r$ so does $u'$ ($d$ is a unit mod $r$), and $\bar\chi(u)=\bar\chi(a^2u')=\chi(d)^2\bar\chi(u')$ since $u\equiv a^2u'$. Hence
$(f\otimes\chi)|_k\gamma=\chi(d)^2(f\otimes\chi)$. Holomorphy and vanishing at the cusps hold for each $f|_k\sm1{u/r}01$, by the argument of \Cref{prop:hecke:double-coset}(ii). For quadratic $\chi$,
$\chi(d)^2=1$ as $(d,r)=1$.

(iii) The coefficients $\chi(n)a_n$ are integers. (iv) By \Cref{prop:hecke:Tp-q-expansion}: if $p\nmid L$ (so $p\nmid M$), the $n$-th coefficient of $T_p^{(L)}(f\otimes\chi)$ is
$\chi(pn)a_{pn}+p^{k-1}\chi(n/p)a_{n/p}=\chi(p)\chi(n)\bigl(a_{pn}+p^{k-1}a_{n/p}\bigr)$, using $\chi(n/p)=\chi(n)\chi(p)^{-1}=\chi(n)\chi(p)$ for $p\mid n$ ($\chi(p)=\pm1$); this is $\chi(p)a_n((T_p^{(M)}f)\otimes\chi)$. If $p\mid L$ and
$p\nmid r$, then $p\mid M$ and both sides are $\mathcal U_p$: $\chi(pn)a_{pn}=\chi(p)\chi(n)a_{pn}$. If $p\mid r$, $\chi(pn)=0$. For the consequence: $T_p\iota=\chi(p)\iota T_p$ gives
$(T_p-\chi(p)\lambda_p)^m\iota=\chi(p)^m\iota(T_p-\lambda_p)^m$, and the normalised eigenform statement follows from \Cref{lem:hecke:eigen-criterion} (the coefficients $\chi(n)a_n$ are again
multiplicative and satisfy the recursions, with $b_{p^r}=0$ for $p\mid r$).

(v) $\bar\chi(n)\chi(n)a_n$ is $a_n$ for $(n,r)=1$ and $0$ otherwise. (vi) Linearity and integrality are (iii); the coefficients of $h(dz)$ are supported on multiples of $d$, where $\chi_S$ vanishes if
$(d,D)>1$; $\iota(h)=\iota(h^S)$ and $\iota\circ\iota(h)=h^S$ are (v); the last statement is (iv) for $\chi_S$, which is primitive of conductor $D$ (\Cref{ch:rings}).
\end{proof}

\begin{theorem}[Atkin and Li]\label{thm:hecke:atkin-li}
Let $g$ be a newform of weight $2$ on $\Gamma_0(M)$, $q$ an odd prime and $\chi=\chi_{q^*}$.
\begin{enumerate}[label=(\roman*)]
\item If $q\nmid M$, $g\otimes\chi$ is a newform on $\Gamma_0(Mq^2)$.
\item If $q\,\|\,M$, $g\otimes\chi$ is a newform on $\Gamma_0(Mq)$.
\item If $q^2\mid M$, there is a unique newform $g'$ with $a_p(g')=\chi(p)a_p(g)$ for all $p\ne q$; its level divides $M$ and has the same exponent as $M$ at every $p\ne q$, and
$g\otimes\chi=(g')^{\{q\}}$, which equals $g'$ iff $a_q(g')=0$.
\end{enumerate}
In cases (i) and (ii) the level is $\lcm(M,q^2)$: this is the paper's Lemma lem:twist (iii). Parts (i), (ii) are stated with a reference, \cite[Theorem 3.1]{AtkinLi78} as cited by the
paper (the theorem number and the exact form of the hypothesis there to be verified), \cite{Li75} for the case $q\nmid M$, \cite[Section 3.6]{Shimura71}; part (iii) is proved here.
\end{theorem}

\begin{proof}[Sketch of proof of (i), (ii), and proof of (iii)]
(i), (ii): through \Cref{ch:modgalrep}. The level of a newform is the conductor of its $\lambda$-adic Galois representation $\rho_g$ (Carayol \cite{Carayol86}), and the naive twist
$g\otimes\chi$, an eigenform by \Cref{prop:hecke:twist-operator}(iv), has the system of eigenvalues of $\rho_g\otimes\chi$ at $p\ne q$. If $\rho_g$ is unramified at $q$ (case (i)) or its restriction
to inertia at $q$ is unipotent and nontrivial (case (ii), Steinberg), then $\rho_g\otimes\chi$ has no inertia invariants at $q$ and is tamely ramified there, so its conductor exponent at $q$ is
$2$; at $p\ne q$ the conductor is unchanged. So the newform $g'$ attached to $\rho_g\otimes\chi$ has level $\lcm(M,q^2)$, hence $a_q(g')=0$ (\Cref{prop:hecke:newform-local}(iii)); at the primes $p\mid M$,
$p\ne q$, $a_p(g')=\chi(p)a_p(g)$ by the argument given for (iii) below, run at the level $\lcm(M,q^2)$ in place of $M$; so $g'$ has the coefficients $\chi(n)a_n(g)$, and $g'=g\otimes\chi$.
(iii) By \Cref{prop:hecke:twist-operator}(ii),(iv), $g\otimes\chi\in S_2(\lcm(M,q^2))=S_2(M)$ is a normalised eigenform of $\T_M$, with $T_n$-eigenvalue $\chi(n)a_n(g)$ for $(n,M)=1$, $U_p$-eigenvalue
$\chi(p)a_p(g)$ for $p\mid M$, $p\ne q$, and $U_q$-eigenvalue $0$. By \Cref{thm:hecke:structure}(i) it lies in $V_{g'}$ for a newform $g'$ of level $M'\mid M$ with $a_p(g')=\chi(p)a_p(g)$ for $p\nmid M$,
unique by \Cref{thm:hecke:strong-multiplicity-one}(i). Symmetrically, $g'\otimes\chi$ lies in $S_2(\lcm(M',q^2))\subseteq S_2(M)$ and is an eigenvector of $\T^{(M)}$ with the system of $g$, so it is a multiple of $g$ by
\Cref{thm:hecke:atkin-lehner}(ii), in fact $g$ (compare $a_1$). Thus $g\in S_2(\lcm(M',q^2))$, and this forces $M\mid\lcm(M',q^2)$: a newform of level $M$ lies in $S_2(L)$ only if $M\mid L$, by the
uniqueness of coordinates in the bases of \Cref{thm:hecke:atkin-lehner}(iii) (the basis at level $L$ is part of the basis at level $\lcm(L,M)$). Hence $v_p(M')=v_p(M)$ for every
$p\ne q$. For $p\mid M$, $p\ne q$, the factor $V_{g',p,0}$ of $V_{g'}\simeq\bigotimes_pV_{g',p,e_p}$ is the line of $g'$, on which $U_p$ acts by $a_p(g')$; so $a_p(g')=\chi(p)a_p(g)$ there as well. At $q$, the
$U_q$-eigenvalue of $g\otimes\chi$ is $0$, and by \Cref{prop:hecke:Up-oldspan}(iii) the $0$-eigenvector of $U_q$ on $V_{g',q,e_q}$ is the depleted form. A normalised eigenform of $\T_M$ in $V_{g'}$ is
determined by its $U_p$-eigenvalues (each joint eigenspace is a tensor product of lines), so $g\otimes\chi=(g')^{\{q\}}$. Finally $(g')^{\{q\}}=g'$ iff $a_{q^r}(g')=0$ for $r\ge1$, iff $a_q(g')=0$.
\end{proof}

\begin{example}[levels of twists]\label{ex:hecke:twist-table}
Conductors of twisted curves, computed with PARI (\texttt{elltwist}, \texttt{ellglobalred}, minimal model matched to Cremona's tables); the newform of $E^{(q)}$ is the newform twist
of $f_E$, equal to the naive twist in cases (i) and (ii):
\begin{center}
\begin{tabular}{llll}
\toprule
case & curve, character & twist & level\\
\midrule
$q\nmid M$ & \curve{11a1}, $\chi_{-7}$ & \curve{539d2} & $11\cdot7^2$\\
 & \curve{20a1}, $\chi_{-7}$ & \curve{980g2} & $20\cdot7^2$\\
 & \curve{37a1}, $\chi_5$ & \curve{925e1} & $37\cdot5^2$\\
$q\,\|\,M$ & \curve{11a1}, $\chi_{-11}$ & \curve{121d2} & $11^2$\\
 & \curve{14a1}, $\chi_{-7}$ & \curve{98a3} & $2\cdot7^2$\\
 & \curve{77c1}, $\chi_{-7}$ & \curve{539c1} & $11\cdot7^2$\\
 & \curve{15a1}, $\chi_{-3}$ & \curve{45a4} & $5\cdot3^2$\\
$q^2\mid M$ & \curve{539c1}, $\chi_{-7}$ & \curve{77c1} & level drops\\
 & \curve{539d1}, $\chi_{-7}$ & \curve{11a3} & level drops\\
 & \curve{980i1}, $\chi_{-7}$ & \curve{140b1} & level drops\\
 & \curve{121a1}, $\chi_{-11}$ & \curve{121c2} & level kept\\
 & \curve{99a1}, $\chi_{-3}$ & \curve{99c1} & level kept\\
 & \curve{49a1}, $\chi_{-7}$ & \curve{49a3} & level kept\\
\bottomrule
\end{tabular}
\end{center}
For all $n\le3000$ we also checked the coefficient identities
\[
a_n(\curve{121d1})=\chi_{-11}(n)a_n(\curve{11a1}),\qquad a_n(\curve{539d1})=\chi_{-7}(n)a_n(\curve{11a1}),
\]
\[
a_n(\curve{539c1})=\chi_{-7}(n)a_n(\curve{77c1}),\qquad a_n(\curve{121c1})=\chi_{-11}(n)a_n(\curve{121a1}),
\]
\[
a_n(\curve{980g1})=\chi_{-7}(n)a_n(\curve{20a1});
\]
the last fails with \curve{980i1} in place of \curve{980g1}. So $f_{\text{20a}}\otimes\chi_{-7}$ is the form of
980g, and the paper's congruence ``$f_{\text{980i}}\equiv f_{\text{20a}}\otimes\chi_{-7}=f_{\text{980g}}\pmod5$'' (Example of Section 3) is a congruence between two different newforms of level $980$, as the paper
writes it. (As $5\mid980$, the paper uses this example only as an illustration.) In case $q^2\mid M$ the naive twist is a depletion: $a_7(\curve{77c1})=-1$, while $f_{\text{539c}}\otimes\chi_{-7}$ has
$a_7=0$.
\end{example}

\begin{paperlink}\label{pl:hecke:atkin-li}
Lemma lem:twist (iii) is \Cref{thm:hecke:atkin-li}(i),(ii); the paper cites \cite[Theorem 3.1]{AtkinLi78} and adds, in the proof of Lemma lem:twist, that ``in the case $q\,\|\,M$, the local
component of $g$ at $q$ is an unramified twist of the Steinberg representation, whose twist by a ramified quadratic character has conductor exponent $2$'', which is the sketch
above in local language. The hypothesis $q^2\nmid M$ cannot be dropped (\Cref{thm:hecke:atkin-li}(iii) and the last rows of \Cref{ex:hecke:twist-table}). The proof of Theorem thm:qual (b) uses
(ii) (``$h=g\otimes\chi$ is a newform on $\Gamma_0(Mq)$'') and (c) uses (i) (``a newform on $\Gamma_0(Mq^2)$''); its parenthetical ``Note that $h\otimes\chi$ is the $q$-depletion of $g$, not $g$'' is
\Cref{prop:hecke:twist-operator}(v).
\end{paperlink}

\begin{proposition}[transport along the twist]\label{prop:hecke:twist-transport}
Let $S$, $D$, $\chi_S$ and $\iota$ be as in \Cref{prop:hecke:twist-operator}(vi). Let $g$ be a newform of weight $2$ and level $M$, $f$ a newform of level $N_f$ with $a_p(f)=\chi_S(p)a_p(g)$ for almost all $p$,
and $L$ a common multiple of $\lcm(M,D^2)$ and $N_f$.
\begin{enumerate}[label=(\roman*)]
\item If $h\in g^\perp$ (at level $M$), then $\iota(h)\in S_2(L)$ is orthogonal to every $f(dz)$, $d\mid L/N_f$. If $L=N_f$, then $\iota(h)\in f^\perp$.
\item If moreover $\iota(g)=f$ as $q$-series and $N_f=L$, a congruence $g\equiv h\pmod r$ with $h\in g^\perp\cap S_2(M;\Z)$ gives $f\equiv\iota(h)\pmod r$ with $\iota(h)\in f^\perp\cap S_2(N_f;\Z)$. Hence $r_g\mid r_f$
when both have rational coefficients.
\item Conversely, suppose $D^2\mid N_f$ and $\iota(g)=f$. Then $\iota(f)=g^S$. If $g^S\in S_2(N_f)$ and its generalised eigenspace there is a line (\Cref{prop:hecke:depletion}(iv)), then $\iota$
maps $f^\perp$ into the $\T$-stable complement of $\C g^S$ in $S_2(N_f)$, and a congruence $f\equiv h\pmod r$ with $h\in f^\perp\cap S_2(N_f;\Z)$ gives $g^S\equiv\iota(h)\pmod r$ with $\iota(h)$ integral in that
complement.
\end{enumerate}
\end{proposition}

\begin{proof}
(i) By \Cref{thm:hecke:atkin-lehner}(iv) at level $M$, $g^\perp=\bigoplus_{h'\ne g}V_{h'}$, the sum over newforms $h'\ne g$ of levels dividing $M$; it suffices to treat $h\in V_{h'}$, an eigenvector of
$\T^{(M)}$ for the system of $h'$. By \Cref{prop:hecke:twist-operator}(vi), $\iota(h)$ is an eigenvector of the $T_p$, $p\nmid L$, with eigenvalues $\chi_S(p)a_p(h')$. These differ from
$a_p(f)=\chi_S(p)a_p(g)$ at infinitely many $p$, because $a_p(h')\ne a_p(g)$ for infinitely many $p$ (\Cref{thm:hecke:strong-multiplicity-one}(i)). So \Cref{thm:hecke:strong-multiplicity-one}(iii) applies. If
$L=N_f$, the only $f(dz)$ is $f$. (ii) Apply $\iota$, which is linear and preserves integrality and congruences coefficientwise. (iii) $\iota(f)=\iota\circ\iota(g)=g^S$ by
\Cref{prop:hecke:twist-operator}(v). At level $N_f$, $\lcm(N_f,D^2)=N_f$, so $\iota$ maps $S_2(N_f)$ to itself. By \Cref{thm:hecke:atkin-lehner}(iv), $f^\perp=\bigoplus_{h'\ne f}V_{h'}$ over newforms of
levels dividing $N_f$, and $\iota(V_{h'})$ lies in the eigenspace of $\T^{(N_f)}$ for the system $\chi_S(p)a_p(h')$. If that were the system of $g$, then $a_p(h')=\chi_S(p)a_p(g)=a_p(f)$ for almost all $p$,
and $h'=f$ by \Cref{thm:hecke:strong-multiplicity-one}(i). So $\iota(f^\perp)$ lies in the sum of the eigenspaces of $\T^{(N_f)}$ for systems other than that of $g$, which is contained in the sum of
the generalised eigenspaces of $\T$ other than that of $g^S$ (every generalised eigenspace of $\T$ lies in one eigenspace of $\T^{(N_f)}$), that is, in the $\T$-stable complement of $\C g^S$.
\end{proof}

\begin{remark}[the engine of Theorems thm:twist and thm:quant]\label{rem:hecke:transport-engine}
\Cref{prop:hecke:twist-transport}(ii) is the proof of Theorem thm:twist; (iii) with \Cref{prop:hecke:complements} gives the reverse of Step 1 of the proof of Theorem thm:quant. There
$g=f_{E_S}$ has level $N_S$, $f=f_E$ has level $N=N_S^S$ with $q^2\mid N$ for $q\in S$, $\iota(g)=f$ (the paper: ``both are multiplicative, $a_p(f)=\chi_S(p)a_p(g)$ for $p\nmid D$ since the twist is
unramified outside $S$, and $a_{q^k}(f)=0$ for $q\in S$ since $q^2\mid N$''), and $g^S$ has a line as generalised eigenspace at level $N$ (\Cref{prop:hecke:depletion}(iv)). By (iii), a
congruence of modulus $r_E$ for $f$ gives one for $g^S$ in its $\T$-stable complement, so by the argument of Lemma lem:crit for the eigencharacter $\pi$ of $g^S$
(\Cref{prop:hecke:complements}(i), \Cref{ch:modgalrep}) $r_E$ divides the order of $\T/(\ker\pi+\Ann_\T\ker\pi)$, the generator of $\eta^{\mathrm{Hecke}}(g^S)$. Together with the paper's Step 1, which
gives $v_\ell(r_E)\ge v_\ell(\eta^{\mathrm{Hecke}}(g^S))$ for $\ell\notin S$, this makes Step 1 an equality for every $\ell\notin S$, in particular for every $\ell\nmid N$, with no hypothesis on $S_f$. The
revised paper's Remark rem:afterquant (2) still calls this question open (``Whether Step 1 is an equality without that hypothesis is open''). At level $539$ the Hecke congruence modules
of $(\text{77c})^{\{7\}}$ and $(\text{11a})^{\{7\}}$ have orders $288=r(\curve{539c1})$ and $504=r(\curve{539d1})$ (\Cref{ex:hecke:twist-pairs}). The theorems themselves are proved in
\Cref{ch:congruence}.
\end{remark}

\section{The $L$-function of a newform; Rankin-Selberg}\label{sec:hecke:lfunctions}

This section attaches to a normalised eigenform its Dirichlet series and Euler product, states the functional equation, and states the Rankin-Selberg formula that relates
the Petersson norm $\langle f,f\rangle$ to the symmetric square $L$-function. In the paper these appear in Section 1 (``the Petersson norm'', ``$\langle f,f\rangle$ is the special value
$L(\Sym f,2)$ up to the local factors at the bad primes'') and in Section 1.1 ($\delta_p$ as the degree of a local factor). The exact formula with all bad factors is
\Cref{ch:lfunctions}'s.

\begin{proposition}[Euler products]\label{prop:hecke:euler-product}
Let $f$ be a normalised eigenform of $\T_N$ of weight $k$. The series $L(f,s)=\sum_{n\ge1}a_n(f)n^{-s}$ converges absolutely for $\operatorname{Re}s>k/2+1$ (\Cref{ch:modforms}), and there
\[
L(f,s)=\prod_{p\nmid N}\bigl(1-a_pp^{-s}+p^{k-1-2s}\bigr)^{-1}\prod_{p\mid N}\bigl(1-a_pp^{-s}\bigr)^{-1}.
\]
Write $L_p(f,s)$ for the factor at $p$. For a newform of weight $2$: $L_p(f,s)^{-1}=(1-\alpha_pp^{-s})(1-\beta_pp^{-s})$ for $p\nmid N$, $(1\mp p^{-s})$ for $p\,\|\,N$ with $a_p=\pm1$, and $1$ for $p^2\mid N$. Moreover
$L(f^\Sigma,s)=L(f,s)\prod_{p\in\Sigma}L_p(f,s)^{-1}$ for a newform $f$, and for quadratic $\chi$ of conductor $r$ the factor of $L(f\otimes\chi,s)$ at $p\nmid rN$ is
$(1-\chi(p)a_pp^{-s}+p^{k-1-2s})^{-1}$.
\end{proposition}

\begin{proof}
By multiplicativity (\Cref{prop:hecke:eigenforms}(ii)) and absolute convergence,
\[
L(f,s)=\prod_p\sum_{r\ge0}a_{p^r}p^{-rs}.
\]
The series $\sum_ra_{p^r}X^r$ is the image under $\phi_f$ of $\sum_rT_{p^r}X^r$, which is $(1-a_pX+p^{k-1}X^2)^{-1}$ or $(1-a_pX)^{-1}$ by \Cref{thm:hecke:Tn}(v); put $X=p^{-s}$. The weight-$2$ factors use \Cref{prop:hecke:newform-local}(iii). For $f^\Sigma$, whose
coefficients at $p\in\Sigma$ vanish, the factor at $p\in\Sigma$ is $1$, and the others are those of $f$. For $f\otimes\chi$, the coefficients $\chi(n)a_n$ are multiplicative, and
\[
\sum_r\chi(p)^ra_{p^r}X^r=(1-a_p\chi(p)X+p^{k-1}\chi(p)^2X^2)^{-1},\qquad\chi(p)^2=1.
\]
\end{proof}

\begin{theorem}[functional equation]\label{thm:hecke:functional-equation}
Let $f\in S_k(N)$ with $f|_kW_N=\varepsilon_Nf$, $\varepsilon_N=\pm1$ (for instance a newform, \Cref{prop:hecke:newform-local}(i)). Then
\[
\Lambda(f,s)=N^{s/2}(2\pi)^{-s}\Gamma(s)L(f,s)=N^{s/2}\int_0^\infty f(iy)y^{s-1}\,dy
\]
extends to an entire function, bounded in vertical strips, with $\Lambda(f,s)=i^k\varepsilon_N\Lambda(f,k-s)$. In weight $2$ the sign is $-\varepsilon_N$.
\end{theorem}

\begin{proof}[Sketch of proof]
For $\operatorname{Re}s$ large, integrate $f(iy)=\sum a_ne^{-2\pi ny}$ termwise against $y^{s-1}$, using $\int_0^\infty e^{-2\pi ny}y^{s-1}dy=(2\pi n)^{-s}\Gamma(s)$. The relation $f|_kW_N=\varepsilon_Nf$ at $z=iy$ reads
\[
f\Bigl(\frac{i}{Ny}\Bigr)=\varepsilon_Ni^kN^{k/2}y^kf(iy).
\] Split the integral at $y=1/\sqrt N$ and substitute $y\mapsto1/(Ny)$ on $(0,1/\sqrt N)$:
\[
\Lambda(f,s)=N^{s/2}\int_{1/\sqrt N}^\infty f(iy)y^{s-1}dy+i^k\varepsilon_NN^{(k-s)/2}\int_{1/\sqrt N}^\infty f(iy)y^{k-s-1}dy.
\]
Both integrals converge for all $s$, because $f(iy)$ decays exponentially as $y\to\infty$; this gives the continuation, the bound in vertical strips, and the symmetry $s\leftrightarrow k-s$ with
the factor $i^k\varepsilon_N$ (apply the formula to $k-s$ and use $(i^k\varepsilon_N)^2=1$ for even $k$). This is the Mellin transform of RHB's integral transforms; see
\cite[Section 5.10]{DiamondShurman}, \cite[Section 4.3]{Miyake}.
\end{proof}

Checks: \curve{11a1} has $\varepsilon_{11}=-1$ (\Cref{ch:modforms}), sign $+1$ and rank $0$; \curve{37a1} has $\varepsilon_{37}=+1$, sign $-1$ and rank $1$ (PARI's \texttt{ellrootno}).

\begin{theorem}[Rankin-Selberg]\label{thm:hecke:rankin-selberg}
Let $f\in S_k(N)$ and $D(f,s)=\sum_{n\ge1}\abs{a_n(f)}^2n^{-s}$.
\begin{enumerate}[label=(\roman*)]
\item $D(f,s)$ converges for $\operatorname{Re}s>k$, continues meromorphically to $\C$, is holomorphic on $\operatorname{Re}s\ge k$ except for a simple pole at $s=k$, and
\[
\operatorname{Res}_{s=k}D(f,s)=\frac{3\,(4\pi)^k}{\pi\,\Gamma(k)\,\mu(N)}\,\langle f,f\rangle,\qquad\text{in weight }2:\quad\operatorname{Res}_{s=2}D(f,s)=\frac{48\pi}{\mu(N)}\langle f,f\rangle,
\]
with the unnormalised Petersson product.
\item At a good prime $p$ of a newform $f$ of weight $2$ with real coefficients,
\[
\sum_ra_{p^r}^2X^r=\frac{1+pX}{(1-\alpha_p^2X)(1-pX)(1-\beta_p^2X)};
\]
at $p\,\|\,N$ the series is $(1-X)^{-1}$, and at $p^2\mid N$ it is $1$.
\item Hence, for a newform of weight $2$ with rational coefficients,
\[
D(f,s)=\frac{\zeta^{(N)}(s-1)\,L^{(N)}(\Sym f,s)}{\zeta^{(N)}(2s-2)}\prod_{p\mid N}\sum_ra_{p^r}^2p^{-rs},
\]
where $L(\Sym f,s)$ is
normalised as in the paper and in PARI's \texttt{lfunsympow} (Appendix B, \Cref{ch:notation}) and the superscript $(N)$ removes the factors at $p\mid N$. If $N$ is squarefree,
$D(f,s)=\zeta^{(N)}(s-1)L(\Sym f,s)/\zeta^{(N)}(2s-2)$ and
\[
\langle f,f\rangle=\frac{N\,L(\Sym f,2)}{8\pi^3}.
\]
\item In general $D(f,s)$ is \emph{not} the symmetric square $L$-function: it has the extra factor $1+p^{1-s}$ at every good prime, and at the bad primes the relation with $L(\Sym f,s)$
involves correction factors (\Cref{ch:lfunctions}; \cite[Lemma 2.12]{DFG04}, \cite{Shimura76}).
\end{enumerate}
\end{theorem}

\begin{proof}[Sketch of proof of (i); proof of (ii) and (iii)]
(i) Let $E(z,s)=\sum_{\gamma\in\Gamma_\infty\backslash\Gamma_0(N)}\operatorname{Im}(\gamma z)^s$ be the non-holomorphic Eisenstein series of $\Gamma_0(N)$ at $\infty$ ($\Gamma_\infty$ the stabiliser of $\infty$, with $\pm1$). Unfolding
the integral of $\abs f^2y^kE(z,s)$ over $\Gamma_0(N)\backslash\HH$ to the strip $\Gamma_\infty\backslash\HH=\{0\le x\le1\}$ and integrating in $x$ (Parseval) gives
\[
\int_{\Gamma_0(N)\backslash\HH}\abs{f}^2y^kE(z,s)\,d\mu=\int_0^\infty\sum_n\abs{a_n}^2e^{-4\pi ny}y^{s+k-2}dy=\frac{\Gamma(s+k-1)}{(4\pi)^{s+k-1}}D(f,s+k-1).
\]
$E(z,s)$ continues meromorphically with a simple pole at $s=1$ of residue $1/\vol(\Gamma_0(N)\backslash\HH)=3/(\pi\mu(N))$, constant in $z$. Taking residues at $s=1$ gives (i)
\cite{Rankin39}, \cite[Section 1]{Zagier85}, \cite{Shimura76}. (ii) Write $\alpha=\alpha_p$, $\beta=\beta_p$, $\alpha\beta=p$; then $a_{p^r}=(\alpha^{r+1}-\beta^{r+1})/(\alpha-\beta)$ (\Cref{rem:hecke:chebyshev}), and summing four
geometric series,
\[
\sum_ra_{p^r}^2X^r=\frac1{(\alpha-\beta)^2}\Bigl(\frac{\alpha^2}{1-\alpha^2X}-\frac{2p}{1-pX}+\frac{\beta^2}{1-\beta^2X}\Bigr)=\frac{1+pX}{(1-\alpha^2X)(1-pX)(1-\beta^2X)},
\]
an identity of rational functions checked by clearing denominators (it also holds when $\alpha=\beta$, by continuity). At $p\,\|\,N$, $a_{p^r}^2=1$; at $p^2\mid N$, $a_{p^r}=0$ for $r\ge1$. (iii) At a good $p$,
$L_p(\Sym f,s)^{-1}=(1-\alpha^2p^{-s})(1-p^{1-s})(1-\beta^2p^{-s})$ (Appendix B), and $1+p^{1-s}=(1-p^{2-2s})/(1-p^{1-s})$ is the factor of $\zeta(s-1)/\zeta(2s-2)$; multiply over $p\nmid N$. For squarefree $N$ the
local factors at $p\mid N$ are $(1-p^{-s})^{-1}$, which is $L_p(\Sym f,s)$ at a multiplicative prime. The residue of $\zeta^{(N)}(s-1)$ at $s=2$ is $\prod_{p\mid N}(1-p^{-1})$ and
$\zeta^{(N)}(2)=\frac{\pi^2}6\prod_{p\mid N}(1-p^{-2})$, so $\operatorname{Res}_{s=2}D=\frac6{\pi^2}\prod_{p\mid N}(1+p^{-1})^{-1}L(\Sym f,2)=\frac{6N}{\pi^2\mu(N)}L(\Sym f,2)$; compare with (i).
\end{proof}

\begin{example}[numerical checks]\label{ex:hecke:rankin-numerics}
With PARI, to $30$ to $40$ digits ($\langle f,f\rangle$ is $\mu(N)$ times the value of \texttt{mfpetersson}; $L(\Sym f,s)$ comes from \texttt{lfunsympow} and its bad factors from
\texttt{lfuneuler}): for \curve{11a1},
$\langle f,f\rangle=0.0469001478734951878\ldots$ and $L(\Sym f,2)=1.0575992445909578\ldots$, and $11\,L(\Sym f,2)/(8\pi^3)$ agrees with $\langle f,f\rangle$ to all digits; for \curve{37a1},
$\langle f,f\rangle=0.3717541475106960\ldots=37L(\Sym f,2)/(8\pi^3)$. The residue formula (i), with $D(f,s)$ written as in (iii) with the local factors at $p\mid N$ computed from (ii), holds to $30$
digits for \curve{11a1}, \curve{37a1}, \curve{20a1}, \curve{121a1} and \curve{99a1}; for the last three, whose levels are not squarefree, $N\,L(\Sym f,2)/(8\pi^3)$ differs from $\langle f,f\rangle$
($0.0542\ldots$ against $0.0813\ldots$ for \curve{20a1}).
\end{example}

\begin{warning}[the Rankin series is not the symmetric square]\label{warn:hecke:rankin}
The paper's introduction (original version) called the naive adjoint $L$-function ``the Euler product of $\sum a_n(f)^2n^{-s}$''. By \Cref{thm:hecke:rankin-selberg}(ii),(iii) this is not right: at
every good prime the Rankin series carries the extra factor $1+p^{1-s}$, and globally $\sum a_n^2n^{-s}=\zeta^{(N)}(s-1)L^{\nv}(\Sym f,s)/\zeta^{(N)}(2s-2)$, where $L^{\nv}$ is the naive symmetric square of
Appendix B (\Cref{ch:notation}). The squarefree formula $\langle f,f\rangle=NL(\Sym f,2)/(8\pi^3)$ uses the unnormalised Petersson product of \Cref{conv:hecke:standing}(d) and PARI's
normalisation of $L(\Sym f,s)$; the general formula with its bad factors is \Cref{ch:lfunctions}'s.
\end{warning}

\section{Worked examples from the paper's world}\label{sec:hecke:examples}

This section puts the theory to work on the levels that the paper's examples live at: $88$ (where the full Hecke algebra is not reduced), the small levels where congruence
numbers and modular degrees can be compared, the twist pairs of Theorem thm:twist, and level $121$. All numbers were computed with PARI/GP 2.17 through \texttt{cypari2}, against
Cremona's tables for conductors below $10{,}000$ \cite{Cremona97}.

\begin{example}[level $88$]\label{ex:hecke:level-88}
$S_2(88)$ has dimension $9$: three newforms of level $88$, the newform $f_{44}$ of level $44$ (old twice: $f_{44}(z)$, $f_{44}(2z)$), and $f_{11}$ (old four times: $f_{11}(2^jz)$, $0\le j\le3$). The operator $U_2$
has characteristic polynomial $X^7(X^2+2X+2)$ and minimal polynomial $X^2(X^2+2X+2)$. On $V_{f_{11},2,3}$ its matrix is the $4\times4$ matrix of \Cref{ex:hecke:Up-matrices}: a Jordan block of size
$2$ at $0$, with kernel spanned by $f_{11}^{\{2\}}=f_{11}+2f_{11}(2z)+2f_{11}(4z)$ and generalised kernel spanned by it and $f_{11}^{\{2\}}(2z)$, and the two stabilisations $f_{11}-\beta f_{11}(2z)$,
$\alpha,\beta=-1\pm i$. On $V_{f_{44},2,1}$, $U_2$ is nilpotent ($a_2(f_{44})=0$, \Cref{prop:hecke:Up-oldspan}). So $\T_{88}\otimes\Q$ is not reduced (\Cref{thm:hecke:structure}(iii)); the step $\wp\cap I=0$ of
Lemma lem:ddt417 must be proved (as the revised paper does) rather than read off from reducedness; and the paper's ``the eigenvalue $0$ with multiplicity one'' (Step 1 of the proof of
Theorem thm:quant) is a statement about the minimal level $M^\Sigma$, which fails one level higher: $f_{11}^{\{2\}}$ spans its generalised eigenspace at level $44$, not at level $88$
(\Cref{prop:hecke:depletion}(v)). Also $\rank\T_{88}=9$, the pairing has determinant $-1$, $[\T_{88}:\Z T_1+\cdots+\Z T_9]=3$, and the congruence number of the newform of 88a is $16$, while
$\deg\varphi_{\curve{88a1}}=8$ ($2^3\mid88$).
\end{example}

\begin{example}[congruence numbers]\label{ex:hecke:congruence-numbers}
For each rational newform $f$ below we computed $r_f$ as the index of $\Z f\oplus(f^c\cap S_2(N;\Z))$ in $S_2(N;\Z)$ (\Cref{prop:hecke:rf-finite}, \Cref{prop:hecke:complements}(ii)), with $f^c$ the image of
$t-\phi_f(t)$ for a random $t\in\T$ for which $\phi_f(t)$ is a simple root of the characteristic polynomial of $t$; and compared with the modular degree of the optimal curve from
Cremona's \texttt{degphi} file:
\begin{center}
\begin{tabular}{lccccccc}
\toprule
curve & 11a1 & 33a1 & 37a1 & 37b1 & 44a1 & 77a1 & 77b1\\
\midrule
$r_f$ & 1 & 3 & 2 & 2 & 2 & 4 & 20\\
$\degphi$ & 1 & 3 & 2 & 2 & 2 & 4 & 20\\
\bottomrule
\end{tabular}

\medskip
\begin{tabular}{lcccccc}
\toprule
curve & 77c1 & 88a1 & 99a1 & 99b1 & 99c1 & 99d1\\
\midrule
$r_f$ & 6 & 16 & 12 & 12 & 12 & 6\\
$\degphi$ & 6 & 8 & 4 & 12 & 12 & 6\\
\bottomrule
\end{tabular}

\medskip
\begin{tabular}{lcccccccc}
\toprule
curve & 121a1 & 121b1 & 121c1 & 121d1 & 539a1 & 539b1 & 539c1 & 539d1\\
\midrule
$r_f$ & 6 & 4 & 6 & 24 & 320 & 192 & 288 & 504\\
$\degphi$ & 6 & 4 & 6 & 24 & 320 & 192 & 288 & 72\\
\bottomrule
\end{tabular}
\end{center}
The three differences (88a1 at $2$, 99a1 at $3$, 539d1 at $7$) are at primes $p$ with $p^2\mid N$, as the theorem of Agashe, Ribet and Stein allows (\Cref{ch:moddeg}). The paper's remark after
the Example of Section 3 quotes ``the curve 99A1, of conductor $9\cdot11$, with modular degree $4$ and congruence number $12$'' from \cite[Table 1]{ARS12}, and adds ``539d1 has $\degphi=72$ and
$r_E=504=7\cdot72$, with $7^2\mid539$''; both agree with the table.
\end{example}

\begin{example}[the twist pairs of Theorem thm:twist]\label{ex:hecke:twist-pairs}
(a) Case (c) ($N_S=N$): \curve{121a1} has type $\kod{II}$ at $11$, $f_{\text{121a}}\otimes\chi_{-11}=f_{\text{121c}}$ (\Cref{ex:hecke:twist-table}; both have $a_{11}=0$, so the naive twist is the newform), and
$r(\text{121a})=r(\text{121c})=6$, as Theorem thm:twist predicts.
(b) Case $\kod I_n^*$: \curve{539c1} has type $\kod I_3^*$ at $7$, and $f_{\text{539c}}=f_{\text{77c}}\otimes\chi_{-7}$ (Atkin-Li, $7\,\|\,77$); $r(\text{77c})=6$ divides $r(\text{539c})=288$. At $\ell=3$: $v_3(288)=2=v_3(6)+1$,
where $1=v_3(7^2-1)$ is the Euler term $e_3(E,7)$ of Theorem thm:quant (the paper's Definition def:terms; \Cref{ch:lfunctions}).
(c) Case $\kod I_0^*$: \curve{539d1} has type $\kod I_0^*$ at $7$, $f_{\text{539d}}=f_{\text{11a}}\otimes\chi_{-7}$ ($7\nmid11$), $r(\text{11a})=1$ divides $504$, and $v_3(504)=2=0+e_3(E,7)$.
(d) Step 1 of the proof of Theorem thm:quant at level $539$, $S=\{7\}$: the depleted forms $g^S$ are $(\text{77c})^{\{7\}}=g+g(7z)$ and $(\text{11a})^{\{7\}}=f_{11}+2f_{11}(7z)+7f_{11}(49z)$ (\Cref{ex:hecke:depleted}),
each spanning its generalised eigenspace in $S_2(539)$; their Hecke congruence modules, computed as the index of $\Z g^S\oplus((g^S)^c\cap S_2(539;\Z))$ with the $\T$-stable complement
(\Cref{prop:hecke:complements}), have orders $288$ and $504$, equal to $r(\text{539c})$ and $r(\text{539d})$, as \Cref{rem:hecke:transport-engine} predicts; and $\iota((\text{77c})^{\{7\}})=f_{\text{539c}}$,
$\iota((\text{11a})^{\{7\}})=f_{\text{539d}}$ (checked to $n=3000$; $\iota$ kills multiples of $7$, so this is \Cref{ex:hecke:twist-table}). At level $44$ the module of $f_{11}^{\{2\}}$ has order $2$, and
at level $121$ that of $f_{11}^{\{11\}}=f_{11}-f_{11}(11z)$ has order $24=r(\text{121d})$.
\end{example}

\begin{example}[level $121$]\label{ex:hecke:level-121}
$S_2(121)$ has dimension $6$: the newforms of 121a, 121b, 121c, 121d (all rational, with $a_{11}=0$ as $11^2\mid121$), and $V_{f_{11},11,1}=\operatorname{span}\{f_{11}(z),f_{11}(11z)\}$. On the latter $U_{11}$ has
matrix $\sm1100$ ($a_{11}(f_{11})=1$, $11\,\|\,11$), eigenvectors $f_{11}$ (eigenvalue $1$) and $f_{11}^{\{11\}}=f_{11}-f_{11}(11z)$ (eigenvalue $0$, the depleted form, of level $M^{\{11\}}=121$). On $S_2(121)$ the
characteristic and minimal polynomials of $U_{11}$ are $X^5(X-1)$ and $X(X-1)$: $U_{11}$ is semisimple. The paper's Example of Section 3, case (b): ``$f_{\text{11a}}\otimes\chi_{-11}$ is the form of 121d,
and indeed $f_{\text{121a}}\equiv f_{\text{121d}}\pmod 3$ for all $n$ (Sturm's bound $22$)''. The first statement is \Cref{thm:hecke:atkin-li}(ii) ($11\,\|\,11$) and \Cref{ex:hecke:twist-table}; the
congruence is proved in \Cref{ch:modforms} by Sturm's bound. Finally $T_1,\dots,T_6$ are linearly dependent although $\rank\T_{121}=6$, while $T_1,\dots,T_{22}$ span $\T_{121}$
(\Cref{ex:hecke:pairing-table}).
\end{example}

\begin{lstlisting}[language=GP]
\\ The computations of this chapter, in outline (PARI/GP 2.17).
mf = mfinit([88,2],1);  A = mfheckemat(mf,2);      \\ U_2 on S_2(88)
factor(charpoly(A))                                  \\ x^7 (x^2+2x+2)
factor(minpoly(A))                                   \\ x^2 (x^2+2x+2)
\\ integral basis: saturate the q-expansion lattice up to B*max(B,41)
B = mfsturm([88,2]); C = mfcoefs(mf, B*41)[2..B*41+1,]; Z = matrixqz(C,-2);
\\ Petersson Gram matrix at level 22 (mfpetersson divides by mu(22) = 36)
F = mfbasis(mfinit([11,2],1))[1]; mf22 = mfinit([22,2],1);
S = [mfsymbol(mf22,F), mfsymbol(mf22,mfbd(F,2))];
G = matrix(2,2,i,j,mfpetersson(S[i],S[j])); G/G[1,1]  \\ [1,-1/3;-1/3,1/4]
\end{lstlisting}

\section{Exercises}\label{sec:hecke:exercises}

\begin{exercise}\label{ex:hecke:ex-cosets}
Let $p\mid N$ and $U_p^\star=[\Gamma_0(N)\sm p001\Gamma_0(N)]_2$. Show that $U_p^\star=w_NU_pw_N=U_p^\dagger$. Compute $U_2^\star$ on $f_{11}(z)$, $f_{11}(2z)$ at level $22$, and show that $U_2^\star$ does not commute with
$U_2$, so that $\T_{22}$ and $U_2^\dagger$ generate a noncommutative algebra. (Answer: $U_2^\star f_{11}=4f_{11}(2z)$ and $U_2^\star f_{11}(2z)=-\frac12f_{11}-2f_{11}(2z)$; $U_2U_2^\star$ and $U_2^\star U_2$ have the
matrices $\sm4{-1}01$ and $\sm10{-4}4$.)
\end{exercise}

\begin{exercise}\label{ex:hecke:ex-tau}
Using only \Cref{thm:hecke:Tn} and $\dim S_{12}(1)=1$, prove $\tau(mn)=\tau(m)\tau(n)$ for coprime $m,n$, and $\tau(p)^2-\tau(p^2)=p^{11}$. Check with PARI that $\tau(2)=-24$, $\tau(4)=-1472$.
(Answer: $576+1472=2048=2^{11}$.)
\end{exercise}

\begin{exercise}\label{ex:hecke:ex-a11}
For \curve{11a1} check with PARI (\texttt{ellan}) that $a_{mn}=a_ma_n$ for coprime $m,n\le200$, that $a_{p^{r+1}}=a_pa_{p^r}-pa_{p^{r-1}}$ for $p\ne11$, and that $a_{11^r}=1$. Explain why $a_{11^r}=a_{11}^r$ is a
statement about the operator $U_{11}$ of level $11$, and why $a_4=a_2^2$ fails ($a_4(f_{11})=2$, $a_2(f_{11})=-2$) although $f_{11}$ lies in $S_2(22)$: at level $22$ the form $f_{11}$ is not an eigenvector
of $U_2$, while its stabilisations are, with $a_{2^r}=\alpha^r$.
\end{exercise}

\begin{exercise}\label{ex:hecke:ex-pairing33}
At level $33$ find a $\Z$-basis of $S_2(33;\Z)$ from $q$-expansions up to the Sturm bound $8$, write the matrices of $T_1,\dots,T_8$, show that they span a lattice of rank $3$, and check that the
pairing matrix $(a_1(t_jg_i))$ has determinant $\pm1$. At level $77$ show that $T_1,\dots,T_7$ span a sublattice of index $5$ in $\T_{77}$. (Answer: see \Cref{ex:hecke:pairing-table}.)
\end{exercise}

\begin{exercise}\label{ex:hecke:ex-local}
(Point of Lemma lem:hecke.) $S_2(37)$ is spanned by the newforms of 37a and 37b, both rational, with $r=2$ (\Cref{ex:hecke:congruence-numbers}). Show that $t\mapsto(\phi_{\text{37a}}(t),\phi_{\text{37b}}(t))$
identifies $\T_{37}$ with $\{(a,b)\in\Z^2:a\equiv b\pmod2\}$. Decompose $\T_{37}\otimes\Z_\ell$ into local factors for $\ell=2$ and $\ell=3$, compute $\Ann(\ker\pi)$ and $\eta(\pi)$ for $\pi=\phi_{\text{37a}}$, and decide in
each case whether the factor through which $\pi$ does not factor lies in $\ker\pi$ or in $\Ann(\ker\pi)$. (Answer: $\T_2$ is local and $\eta(\pi)=2\Z_2$; $\T_3=\Z_3\times\Z_3$, the second factor lies in
$\ker\pi$, and $\eta(\pi)=\Z_3$. Use \Cref{thm:hecke:newform-splitting}(iv) for the index.)
\end{exercise}

\begin{exercise}\label{ex:hecke:ex-level22}
At level $22$, verify the matrices of \Cref{ex:hecke:level-22} by hand from \Cref{prop:hecke:Up-adjoint}(ii) and \Cref{ch:modforms}'s formula for $w_{22}$, and show that the eigenvectors
$f_{11}-\beta f_{11}(2z)$ of $U_2$ are not orthogonal to each other. Conclude that no basis of $S_2(22)$ consists of orthogonal eigenvectors of $U_2$.
\end{exercise}

\begin{exercise}\label{ex:hecke:ex-jordan}
Let $N=176=2^4\cdot11$, so that $\dim S_2(176)=19$ and $\dim S_2^{\mathrm{new}}(176)=5$. Using \Cref{thm:hecke:structure} and \Cref{prop:hecke:Up-oldspan}, list the spaces $V_h$ and the Jordan
blocks of $U_2$ on each, predict the characteristic and minimal polynomials of $U_2$ on $S_2(176)$, and check them with PARI's \texttt{mfheckemat}. (Answer: $X^{17}(X^2+2X+2)$ and $X^3(X^2+2X+2)$;
blocks of size $3$ at $0$ on $V_{f_{11},2,4}$ and $V_{f_{44},2,2}$, of size $2$ on the three $V_{h,2,1}$ with $h$ new of level $88$, of size $1$ on the five newforms of level $176$.)
\end{exercise}

\begin{exercise}\label{ex:hecke:ex-depletion}
For $f=f_{11}$ compute $f^\Sigma$ and $M^\Sigma$ for $\Sigma=\{2\},\{11\},\{2,11\},\{3,5\}$ (answers: $44$, $121$, $484$, $2475$), and check with PARI that $f^\Sigma$ lies in $S_2(M^\Sigma)$ but not in
$S_2(M^\Sigma/p)$ for any $p\in\Sigma$. Show $L(f^\Sigma,s)=L(f,s)\prod_{p\in\Sigma}L_p(f,s)^{-1}$.
\end{exercise}

\begin{exercise}\label{ex:hecke:ex-539}
(Step 1 of Theorem thm:quant.) For $g$ of class 77c and $g=f_{11}$, with $q=7$ and $N=539$, write the span of the $g(dz)$, $d\mid539/M_g$, the matrix of $U_7$ on it, and show that the eigenvalue $0$
has algebraic multiplicity one. Show that all other eigenvalues there are prime to $3$. (Answer: $X(X+1)$ and $X(X^2+2X+7)$.)
\end{exercise}

\begin{exercise}\label{ex:hecke:ex-dims}
Compute $\dim S_2^{\mathrm{new}}(N)$ for $N=88$, $121$, $539$ from the dimensions of $S_2(M)$, $M\mid N$, with \Cref{cor:hecke:dimensions}, and compare with \Cref{ex:hecke:dimensions}. Prove that $\dim S_2^{\mathrm{new}}(p^3)=\dim
S_2(p^3)-2\dim S_2(p^2)+\dim S_2(p)$.
\end{exercise}

\begin{exercise}\label{ex:hecke:ex-twist-hne}
(Theorem thm:qual (b).) Let $g$ be a newform on $\Gamma_0(M)$ with $q\,\|\,M$, $\chi=\chi_{q^*}$, and $h=g\otimes\chi$, a newform on $\Gamma_0(Mq)$ (\Cref{thm:hecke:atkin-li}). Show that $h\otimes\chi=g^{\{q\}}\ne g$. Let
$f$ and $f'$ be newforms with $a_p(f')=\chi(p)a_p(f)$ for almost all $p$, and suppose that the level of $f'$ is not $M$. Prove that $h\ne f$. (Answer: if $h=f$, then $a_p(g)=\chi(p)a_p(h)=a_p(f')$ for
almost all $p$, so $g=f'$ by \Cref{thm:hecke:strong-multiplicity-one}(i), and the levels would agree.)
\end{exercise}

\begin{exercise}\label{ex:hecke:ex-elltwist}
(Lemma lem:twist (iii).) With PARI's \texttt{elltwist} and \texttt{ellglobalred}, find the conductors of the twists of \curve{11a1} by $-11$ and by $-7$, of \curve{14a1} and \curve{77c1} by $-7$, and of
\curve{539c1}, \curve{539d1}, \curve{121a1} by $-7$, $-7$, $-11$. Which satisfy $q^2\nmid M$, and where does the level drop? Show that the twist of \curve{20a1} by $-7$ lies in the class 980g, not 980i.
\end{exercise}

\begin{exercise}\label{ex:hecke:ex-99}
(The remark after the Example of Section 3.) Compute $r_f$ for the newform of 99a by the lattice method of \Cref{ex:hecke:congruence-numbers}, compare with \texttt{ellmoddegree}, and explain why
the discrepancy can only be at $3$. Find the newforms of level $99$ congruent to it modulo $3$ at the primes $p\ne3,11$. (Answer: those of 99b and 99d; not those of 11a, 33a, 99c, which differ at $p=2$.)
\end{exercise}

\begin{exercise}\label{ex:hecke:ex-rankin}
Prove the identity $\sum_ra_{p^r}^2X^r=(1+pX)/((1-\alpha^2X)(1-pX)(1-\beta^2X))$ by clearing denominators, and deduce that $\sum a_n^2n^{-s}$ has the factor $1+p^{1-s}$ beyond $L_p(\Sym f,s)$. Compute
$\langle f,f\rangle$ for \curve{11a1} with \texttt{mfpetersson} and check $\langle f,f\rangle=11L(\Sym f,2)/(8\pi^3)$.
\end{exercise}

\begin{exercise}\label{ex:hecke:ex-fe}
Compute the sign of the functional equation of $L(f,s)$ for the newforms of \curve{11a1}, \curve{37a1}, \curve{37b1}, \curve{77c1} from \Cref{thm:hecke:functional-equation} and the Atkin-Lehner eigenvalues
$\varepsilon_N$ (PARI's \texttt{mfatkineigenvalues}), and compare with \texttt{ellrootno} and the ranks. (Answer: root numbers $+1$, $-1$, $+1$, $+1$.)
\end{exercise}

\begin{exercise}\label{ex:hecke:ex-galois23}
At level $23$ the Hecke algebra is $\T_{23}=\Z[T_2]$ with $T_2^2+T_2-1=0$ (\Cref{ch:ant}). Show that the two normalised eigenforms are Galois conjugate newforms with $K_f=\Q(\sqrt5)$, and list the
maximal ideals of $\T_{23}$ above $5$ and above $11$ together with the corresponding mod $\ell$ eigenforms $\sum\phi(T_n)q^n$ (\Cref{thm:hecke:eigen-dictionary}). (Hint: $5$ ramifies in $\Q(\sqrt5)$, and
$11=\frac{7+\sqrt5}2\cdot\frac{7-\sqrt5}2$ splits; above $5$ there is one maximal ideal, with residue field $\F_5$, and the two eigenforms become congruent modulo it.)
\end{exercise}

\section{Notes and references}\label{sec:hecke:notes}

Hecke operators on modular forms of level one go back to Hecke and, for the multiplicativity of $\tau$, to Mordell; the double coset formalism is Shimura's \cite[Chapter 3]{Shimura71}.
We followed Diamond and Shurman \cite[Chapter 5]{DiamondShurman} for the double cosets, the $q$-expansion formulas, the Petersson adjoints and the theory of newforms, with
the normalisation of \Cref{conv:hecke:double-coset}; Serre \cite[Chapter VII]{SerreCourse} and Koblitz \cite{Koblitz} treat level one, and Miyake \cite[Chapter 4]{Miyake} the general theory
with characters. The section numbers of Diamond-Shurman (5.1 to 5.10, 6.3), Shimura (3.1 to 3.6, 7.2), Miyake (4.3, 4.6), Agashe-Ribet-Stein (Section 3), Darmon-Diamond-Taylor (Section 1)
and Zagier (Section 1) cited in the text are pointers we have not verified against the printed editions (for Diamond-Flach-Guo we use Section 1.7.3, the section the paper cites), and the theorem numbers \cite[Theorem 4.6.19]{Miyake}, \cite[Theorem 3.1]{AtkinLi78} and \cite[Theorem 3]{AtkinLehner70} are quoted as the paper (for the last
two) and the literature cite them, not checked.

The theory of newforms is due to Atkin and Lehner \cite{AtkinLehner70} for trivial character and to Li \cite{Li75} in general; strong multiplicity one in the form of
\Cref{thm:hecke:strong-multiplicity-one}(i) is in Miyake's book and in Li's paper; the density-one form uses the Galois representations of \Cref{ch:modgalrep}. The theorem on twists of
newforms is Atkin and Li \cite{AtkinLi78}. The semisimplicity of $U_p$ on $p$-old forms of weight $2$ is Coleman and Edixhoven \cite{ColemanEdixhoven}.

The perfect pairing $S_2(N;\Z)\times\T\to\Z$ is cited by the paper from Ribet \cite[Theorem 2.2]{Ribet83} and Abbes-Ullmo \cite[Lemma 2.1]{AU96} (see also Agashe-Ribet-Stein
\cite[Section 3]{ARS12} and \cite[Section 1]{DDT}); the proof given here, from the integral basis theorem by a rank argument, is standard. The dictionary between eigenforms and
homomorphisms (\Cref{thm:hecke:eigen-dictionary}) is the starting point of Mazur's and Ribet's use of $\T$; see \cite{DDT}. The notation $f^\Sigma$, $M^\Sigma$ follows Diamond, Flach and Guo
\cite[Section 1.7.3]{DFG04}. The Rankin-Selberg method is Rankin's \cite{Rankin39}; the relation with the symmetric square and the bad factors is Shimura's \cite{Shimura76} and
\cite[Lemma 2.12]{DFG04}; Zagier \cite{Zagier85} gives the computation for elliptic curves. Stein's book \cite{Stein} explains the algorithms, and PARI/GP \cite{PARI} implements them
(\texttt{mfinit}, \texttt{mfheckemat}, \texttt{mfpetersson}, \texttt{mftwist}, \texttt{mfcoefs}); the curve data are Cremona's \cite{Cremona97}.

Two results here go beyond what the paper states. \Cref{thm:hecke:newform-splitting} proves the decomposition $\T\otimes\Q=\Q\times\T'$ that the proof of Lemma lem:crit writes down, and shows
that the $\T$-stability of $f^\perp$ needs no adjoints of the $U_p$. \Cref{prop:hecke:twist-transport}(iii) gives the converse transport along $\iota$, which makes Step 1 of the proof of
Theorem thm:quant an equality for every $\ell\nmid N$ (\Cref{rem:hecke:transport-engine}); the revised paper leaves this open in its Remark rem:afterquant (2).
